% concatenated from zuiddam_2025.tex
%% daj-template.tex v0.33     23 Sep 2016   Alex Russell/Laszlo Babai
%%
%% AUTHOR: Fill in fields (or see warnings) below marked with "AUTHOR"
%% ** Add as few macro / package definitions as possible
%% ** Compile with "pdflatex"; make sure that
%%           daj.cls and tocbase.cls are in the same directory.
%%
%% EDITOR: Fill in fields below marked with "EDITOR"
%%    and check that authors proprely filled in field marked with "AUTHOR"

\documentclass{daj}

\usepackage[utf8]{inputenc}
\usepackage{mathtools,amsmath}
\usepackage{amssymb}
\usepackage{microtype}
\usepackage{enumerate}
\usepackage{accents}
%\usepackage{graphicx}
%\usepackage[table, usenames,dvipsnames]{xcolor}
%\usepackage{tikz}
%\usetikzlibrary{calc}
%\usepackage{tabularx}
%\usepackage{colortbl}
%\usepackage{authblk}
\usepackage{arydshln}
%\usepackage{algorithm}
%\usepackage{algpseudocode}

%\usepackage[colorlinks, pagebackref=true]{hyperref}
% \hypersetup{
%   linkcolor=[rgb]{0.3,0.3,0.6},
%   citecolor=[rgb]{0.2, 0.6, 0.2},
%   urlcolor=[rgb]{0.6, 0.2, 0.2}
% }
% \hypersetup{breaklinks=true}
\def\chapterautorefname{Chapter}
\def\sectionautorefname{Section}
\def\subsectionautorefname{Section}
\def\subsubsectionautorefname{Section}
%\def\algorithmautorefname{Algorithm}

%\allowdisplaybreaks[1]

\usepackage{amsthm}
\usepackage{thmtools, thm-restate} 
\declaretheorem[name=Theorem, parent=section]{theorem}
\declaretheorem[name=Corollary, sibling=theorem]{corollary}
\declaretheorem[name=Proposition, sibling=theorem]{proposition}
\declaretheorem[name=Lemma, sibling=theorem]{lemma}
\declaretheorem[name=Conjecture, sibling=theorem]{conjecture}
\declaretheorem[name=Observation, sibling=theorem]{observation}
\declaretheorem[name=Fact, sibling=theorem]{fact}
\declaretheorem[name=Definition, sibling=theorem, style=definition]{definition}
\declaretheorem[name=Remark, sibling=theorem, style=definition]{remark}
\declaretheorem[name=Problem, sibling=theorem, style=definition]{problem}
\declaretheorem[name=Example, sibling=theorem, style=definition]{example}
\declaretheorem[name=Claim, sibling=theorem, style=remark]{claim}

\DeclareMathOperator{\Diag}{Diag}
\DeclareMathOperator{\supp}{supp}
\DeclareMathOperator{\subrank}{Q}
\DeclareMathOperator{\slicerank}{SR}
\DeclareMathAccent{\wtilde}{\mathord}{largesymbols}{"65}
\DeclareMathOperator{\asympsubrank}{\underaccent{\wtilde}{Q}}
\DeclareMathOperator{\asympslicerank}{\underaccent{\wtilde}{SR}}
\DeclareMathOperator{\asympsupslicerank}{{SR}\!\wtilde{\phantom{X}\!\!}}
\DeclareMathOperator{\asymprank}{\underaccent{\wtilde}{R}}
\DeclareMathOperator{\bordersubrank}{\underline{Q}}
\DeclareMathOperator{\tensorrank}{R}
\newcommand{\FF}{\mathbb{F}}
\newcommand{\KK}{\mathbb{K}}
\newcommand{\RR}{\mathbb{R}}
\newcommand{\CC}{\mathbb{C}}
\newcommand{\NN}{\mathbb{N}}
\newcommand{\ZZ}{\mathbb{Z}}
\newcommand{\eps}{\varepsilon}
\DeclareMathOperator{\GL}{GL}

\DeclareMathOperator{\Mat}{Mat}

\DeclareMathOperator{\crk}{crk}
\DeclareMathOperator{\ncrk}{ncrk}
\DeclareMathOperator{\linspan}{span}
\DeclareMathOperator{\rank}{rank} \DeclareMathOperator{\col}{colspan} 

%\DeclareMathOperator{\Diag}{Diag}
%

\DeclareMathOperator{\maxrank}{max-rank}
\DeclareMathOperator{\minrank}{min-rank}
\DeclareMathOperator{\minrnk}{min-rank}
\DeclareMathOperator{\mincov}{min-cov}
\DeclareMathOperator{\minsupp}{min-supp}
\DeclareMathOperator{\maxsupp}{max-supp}


\newcommand{\id}{\mathrm{Id}}
\newcommand{\overCl}[1]{{{#1}^{\overline{\FF}}}}
\newcommand{\overK}[1]{{#1^{\mathbb{K}}}}

% \makeatletter
% \newcommand{\longdash}[1][2em]{%
%   \makebox[#1]{$\m@th\smash-\mkern-7mu\cleaders\hbox{$\mkern-2mu\smash-\mkern-2mu$}\hfill\mkern-7mu\smash-$}}
% \makeatother

%
  \newcommand{\beq}{\begin{equation}}
  \newcommand{\eeq}{\end{equation}}
  \newcommand{\beqn}{\begin{equation*}}
  \newcommand{\eeqn}{\end{equation*}}
  \newcommand{\beqr}{\begin{eqnarray}}
  \newcommand{\eeqr}{\end{eqnarray}}
  \newcommand{\beqrn}{\begin{eqnarray*}}
  \newcommand{\eeqrn}{\end{eqnarray*}}


% \definecolor{diagonalBlockColour}{rgb}{0.48,0.58,0.87}
% \definecolor{rowDiagBlockColour}{rgb}{0.24,0.49,0.03}

\setcounter{tocdepth}{2}


%%%%%%%%%%%%%%%%%%%%%%%%%%%%%%%%%%%%%%%%%%%%%%%%
%% AUTHOR: Fill in meta-data below:
\dajAUTHORdetails{%
  title = {Discreteness of Asymptotic Tensor Ranks}, %% please capitalize all significant words
  author = {Jop Briët, Matthias Christandl, Itai Leigh, Amir Shpilka, and Jeroen Zuiddam},
    %% Please use the format for commas as follows:
    %% "A", or "A and B", or "A, B, and C", or "A, B, C, and D", etc.
  plaintextauthor = {Jop Briët, Matthias Christandl, Itai Leigh, Amir Shpilka, Jeroen Zuiddam},
    %% An author list in plain text: Use the format
    %% "A", or "A, B", or "A, B, C", etc.
    %% NOTE: No LaTeX code in author names.
    %% NOTE: No "and" at the end--simply comma separated,
    % 
 %% The remaing lines in this section are optional:
    %
    %% IF YOUR TITLE CONTAINS MATH OR LATEX such as accented characters: 
    %% Add a "plain text title";  otherwise comment out the next line:
  %plaintexttitle = {Discreteness of Asymptotic Tensor Ranks}, %%  title without math or LaTeX
    %
    %% ONLY IF YOUR TITLE IS TOO LONG to fit in the page headers, please 
    %% add an abbreviated version of the title, otherwise comment it out:
  %runningtitle = {R\"odl's $n^{\log\log n}$ Bound}, 
    %
    %% ONLY IF YOUR AUTHOR LIST IS TOO LONG to fit in the page headers, 
    %% add an abbreviated version, otherwise comment it out:
  %runningauthor = {Paul Erd\H{o}s, Johan H{\aa}stad, L\'aszl\'o Lov\'asz, and Andrew C-C. Yao},
    %% you can replace first names and/or middle names with initials.
    %
    %% ONLY IF YOUR AUTHOR LIST IS TOO LONG to fit the copyright entry
    %% on the bottom of the front page,
    %% add an abbreviated version, otherwise comment it out:
  %copyrightauthor = {P. Erd\H{o}s, J. H{\aa}stad, L. Lov\'asz, and A. C-C. Yao},
    %% Note that the copyrightauthor  field will seldom be necessary;
    %% for instance, in this example with four authors, it would be 
    %% all right to comment it out and have all authors' full names 
    %% appear on the Copyright line
   %
   %% Include keywords of your choice: comma separated, lower case;
   %% comment out the "keywords" line if you don't wish to provide them
  keywords = {tensors, asymptotic rank, subrank, slice rank, discreteness},
}   %%% END \dajAUTHORdetails

%%%%%%%%%%%%%%%%%%%%%%%%%%%%%%%%%%%%%%%%%%%%%%%%
%%% EDITOR: please fill in the following data:
\dajEDITORdetails{%
   year={2025},
   %volume={XX},
   number={15},
   received={27 September 2023},   % received date: example: 7 January 2017
   revised={17 February 2025},    % Optional revised date (you may comment it out)
   published={10 September 2025},  % published date
   doi={10.19086/da.143834},       % XXX = number of paper, e.g. da006 for paper#6
%                              % or  da0006 (length of string arbitrary)
}   %%% END \dajEDITORdetails

\begin{document}

\begin{frontmatter}[classification=text]
%% EDITOR: this will force the keywords to appear right after the Abstract.
%%   If the abstract is too long and would force the keywords off the
%%   front page, please comment out % [classification=text] above
%%   This way the keywords will be floated on the bottom of the first page
%%   even though the Abstract spills over to the next page.

%%% AUTHOR: Title goes here.  This line is optional.  You must use it
%%   if title has footnote attached or requires nontrivial typesetting,
%%   e.g., inclusion of linebreaks to force nice layout.
\title{Discreteness of Asymptotic Tensor Ranks} %% please capitalize all significant words

%%% AUTHOR:
%%% List all authors. If you wish, place grant acknowledgements in \thanks.
%%% In brackets include a short tag for each author.
\author[briet]{Jop Briët\thanks{Supported by the Dutch Research Council (NWO) as part of the NETWORKS programme (Grant No.~024.002.003).}}
\author[christandl]{Matthias Christandl\thanks{Supported by the European Research Council (ERC Grant Agreement No.~81876), VILLUM FONDEN via the QMATH Centre of Excellence (Grant No.~10059) and the Novo Nordisk Foundation (Grant NNF20OC0059939 ``Quantum for Life'').}}
\author[leigh]{Itai Leigh\thanks{Supported by Erasmus+ International Credit Mobility grant, the Deutsch foundation, Israel Science Foundation (grant number 514/20) and the Len Blavatnik and the Blavatnik Family foundation.}}
\author[shpilka]{Amir Shpilka\thanks{Supported by Israel Science Foundation (grant number 514/20) and the Len Blavatnik and the Blavatnik Family foundation.}}
\author[zuiddam]{Jeroen Zuiddam\thanks{Supported by NWO Veni grant VI.Veni.212.284.}}

%%% AUTHOR: Abstract goes here
\begin{abstract}
    Tensor parameters that are amortized or regularized over large tensor powers, often called ``asymptotic'' tensor parameters, play a central role in several areas including algebraic complexity theory (constructing fast matrix multiplication algorithms), quantum information (entanglement cost and distillable entanglement), and additive combinatorics (bounds on cap sets, sunflower-free sets, etc.). Examples are the asymptotic tensor rank, asymptotic slice rank and asymptotic subrank.
    Recent works (Costa--Dalai, Blatter--Draisma--Rupniewski, Christandl--Gesmundo--Zuiddam) have investigated notions of discreteness (no accumulation points) or ``gaps'' in the values of such tensor parameters. 
    
    We prove a general discreteness theorem for asymptotic tensor parameters of order-three tensors and use this to prove that (1) over any finite field (and in fact any finite set of coefficients in any field), the asymptotic subrank and the asymptotic slice rank have no accumulation points, and (2) over the complex numbers, the asymptotic slice rank has no accumulation points.

    Central to our approach are two new general lower bounds on the asymptotic subrank of tensors, which measures how much a tensor can be diagonalized. The first lower bound says that the asymptotic subrank of any concise three-tensor is at least the cube root of the smallest dimension. The second lower bound says that any concise three-tensor that is ``narrow enough'' (has one dimension much smaller than the other two) has maximal asymptotic subrank.
%
    %Our proofs rely on new lower bounds on the maximum rank in matrix subspaces that are obtained by slicing a three-tensor in the three different directions. We prove that for any concise tensor, the product of any two such maximum ranks must be large, and as a consequence there are always two distinct directions with large max-rank.
\end{abstract}
\end{frontmatter}

%%% AUTHOR: body of paper starts here
\tableofcontents

\newpage
\section{Introduction}

%\ILnote{Comment 6. Added citations below -- there are two that I'm missing and one I'm not sure about}\Jnote{Added some. Only thing I'm wondering is if there is more for entanglement than Watrous. But maybe it is fine.}
Tensor parameters that are amortized or regularized over large tensor powers, often called ``asymptotic'' tensor parameters, play a central role in several areas of theoretical computer science including 
algebraic complexity theory (constructing fast matrix multiplication algorithms \cite{burgisser1997algebraic, blaser2013fast}, and barriers for such constructions \cite{alman2018limits,v017a002}), 
quantum information (entanglement cost and distillable entanglement~\cite{watrous_2018,brandao2016mathematics,vrana2015transformation}), 
and additive combinatorics (bounds on cap sets \cite{tao}, sunflower-free sets \cite{naslund_sawin_2017}, etc.). %
%
Examples are the asymptotic tensor rank (famous for its connection to the matrix multiplication exponent), the asymptotic subrank, and the asymptotic slice rank. %
These asymptotic tensor parameters are of the form $\underaccent{\wtilde}{F}(T) = \lim_{n\to \infty} F(T^{\boxtimes n})^{1/n}$ for some integer-valued function~$F$ (e.g.~tensor rank, subrank or slice rank), where $\boxtimes$ denotes the Kronecker product of tensors. 
The computation of these parameters, which in some applications is the main goal and in others is done to bound other parameters of interest, has turned out to be very challenging, and many questions about them are open despite much interest.


The fundamental question whether the matrix-multiplication exponent~$\omega$ equals~2 is closely related to the question whether asymptotic tensor rank is integral-valued. Contrary to matrix rank, some asymptotic tensor parameters may indeed take non-integral values.
For instance, the asymptotic slice rank of the $W$-tensor (which appears in the study of sunflower-free sets) equals $2^{h(1/3)} \approx 1.889$ \cite{strassen1991degeneration}, %\cite{vrana2015transformation} \ILnote{perhaps \cite{MR1056627} as well, as the paper cites it for the lower bound in the case $k=3$}, 
where $h$ is the binary entropy function, and the asymptotic slice rank of the cap set tensor (which appears in the study of arithmetic progression-free sets or cap sets) is known to be the non-integral value $\approx 2.755$ over the finite field $\FF_3$ \cite{MR3583358, tao, kleinberg2018growth, strassen1991degeneration}. %\ILnote{\cite{tao} shows this as an upper bound. Where is a lower bound showed?}.

This raises the fundamental question, for a given function $F$,
%These observations lead to the following questions which play a central role in this paper:
%\begin{itemize}%
    %\item 
    what values $\underaccent{\wtilde}{F}(T)$ can take when varying $T$ over all tensors of order three with arbitrary dimensions (over any fixed field). %In other words, what is the image of 
    More generally, what is the structure (geometric, algebraic, topological, etc) of the set of values 
    \[
    \{\underaccent{\wtilde}{F}(T) : T \in \FF^{n_1} \otimes \FF^{n_2} \otimes \FF^{n_3}, n_1, n_2, n_3 \in \NN\}?
    \]
    %$\underaccent{\wtilde}{F}$ on $\FF^{n_1} \otimes \FF^{n_2} \otimes \FF^{n_3}$ for all dimensions $n_1, n_2, n_3 \in \NN$? 
    %\item  
    % \ILnote{I think the following paragraph starts with no clear connection or motivation, so I suggest adding to its beginning  something like "Recently, very basic properties of these structures started to be discovered. Blatter, Draisma and Rupniewski showed in \cite{https://doi.org/10.48550/arxiv.2212.12219} 
    % these sets are countable over the complex numbers and in \cite{bdr} that they are well-ordered over any finite field. We focus on a question generalising both these results:"} 
    % \JBnote{I agree that for the CS community, some more motivation from the CS perspective is nice early on. These results by Blatter et al.\ are pretty mathematical though. Might matrix multiplication be an alternative as I suggested above?}
    Does $\underaccent{\wtilde}{F}(T)$ have accumulation points, that is, are there non-trivial sequences %\JBnote{``non-constant'' is a bit confusing to me. I guess we mean that the values $\underaccent{\wtilde}{F}(T_i)$ are distinct. Could we say non-trivial sequences?} 
    of tensors $T_1, T_2, \ldots$ such that~$\underaccent{\wtilde}{F}(T_i)$ converges?
    %to a value that is not in the image of $\underaccent{\wtilde}{F}$? 
    Or is it discrete? 
    %\item 
    What ``gaps'' are there between the possible values? Even when $\FF$ is a finite field, the answers to these questions are a priori not clear at all.
%\end{itemize}

In this paper we prove \emph{discreteness} of asymptotic tensor rank, asymptotic subrank and asymptotic slice rank in several regimes. This means that the values of each of these parameters have no accumulation points.
In fact, the proof of discreteness of asymptotic tensor rank (over any finite field or finite coefficient set in any field) follows from a simple argument. Using that same simple argument, together with several new results about tensors, we obtain discreteness for the other parameters. 
In particular, as a core ingredient, we prove a new result about diagonalizability of tensors. This comes in the form of a lower bound on the asymptotic subrank that relies only on the dimensions of the tensor (as opposed to the well-known laser method for fast matrix multiplication, which relies on much more information about the tensor). As another core ingredient we prove new results about matrix subspaces and their max-rank and min-rank. In particular, we prove that the max-ranks of the matrix spaces obtained by slicing a tensor in the three different directions are strongly related, in such a way that at least two of them must be large.

Our discreteness results show that there is a surprising rigidity in the asymptotic behaviour of tensors.
The discreteness of the above parameters gives rise to the phenomenon of ``rounding'' or ``boosting'' (upper or lower) bounds on them to the next possible value (although making this effective requires more knowledge of the possible values than just knowing discreteness). %\ILnote{Comment 1. I added in the following the restriction to finite subset of coefficients, and changed the theorem to include this as well.}
%\Jnote{I think here it's fine to leave out the detail, to keep the statement easier to read.}
%The discreteness of asymptotic tensor rank over finite sets of coefficients implies, for instance, that if we fix such a finite set that includes $\{0,1\}$, the asymptotic rank of the matrix multiplication tensor is bounded away from (or equal to) the asymptotic tensor rank of any other tensor that has coefficients in this finite set. 
The discreteness of asymptotic tensor rank\footnote{Here we work in the regime where we fix the field to be any finite field, or we use any finite set of coefficients in any field.} implies, for instance, that the asymptotic rank of the matrix multiplication tensor is bounded away from (or equal to) the asymptotic tensor rank of any other tensor. In particular, the matrix multiplication exponent $\omega$ is an isolated point among the exponents of all bilinear maps.
%In particular, the matrix multiplication exponent $\omega$ is an isolated point among the exponents of all bilinear maps that use a finite set of coefficients. 
If such a tensor is ``close enough'' to the matrix multiplication tensor, its exponent must ``snap'' to $\omega$. Similar statements hold for the other asymptotic parameters. In the context of combinatorial applications, this may moreover lead to limitations for the asymptotic slice rank to improve on existing results.

Before stating our results in detail, we will first discuss the various asymptotic ranks, their applications and context.

% a general discreteness theorem for asymptotic tensor parameters of order-three tensors and use this to prove that (1) over any finite field, the asymptotic subrank and the asymptotic slice rank have no accumulation points, and (2) over the complex numbers, the asymptotic slice rank has no accumulation points.

\subsubsection*{Matrix multiplication and asymptotic tensor rank}
Determining the computational complexity of matrix multiplication is a fundamental problem in algebraic complexity theory.  This complexity is controlled by the famous matrix multi\-pli\-cat\-ion exponent $\omega$, which is defined as the infimum over nonnegative real numbers~$\tau$ such that any two $n\times n$ matrices can be multiplied using $\mathcal{O}(n^{\tau})$ arithmetic operations \cite{burgisser1997algebraic,blaser2013fast}. The naive matrix multiplication algorithm gives the upper bound $\omega \leq 3$. In 1969, Strassen proved that $\omega \leq 2.81$  \cite{strassen1969gaussian}. Since then, using many different techniques, the best upper bound has been brought down to $\omega \leq 2.371552$ \cite{le2014powers, DBLP:conf/soda/AlmanW21, duan2023faster, williams2023new}. There is a tantalizing possibility (and many have conjectured) that $\omega = 2$, and routes have been proposed that aim to prove this \cite{cohn2003group, cohn2005group,cohn2013fast,MR3631613,blasiak2022matrix}. It is just as intriguing to consider the possibility that~$\omega > 2$ and $\omega$ giving rise to a new fundamental constant, and there has been much work on this lower bound direction as well \cite{DBLP:conf/stoc/BurgisserI11, MR3081636}.

%\subsection*{Asymptotic tensor rank}
Not only can we currently not determine the value of $\omega$, or decide whether $\omega = 2$ or~$\omega > 2$, there is a much more relaxed problem that we cannot solve. Indeed $\omega$ is naturally described in terms of tensors as the logarithm of the asymptotic tensor rank of the matrix multiplication tensor $\langle 2,2,2\rangle \in \FF^4 \otimes \FF^4 \otimes \FF^4$, that is $\omega = \log \asymprank(\langle 2,2,2\rangle)$, and thus $\omega > 2$ is equivalent to $\asymprank(\langle 2,2,2\rangle) > 4$. The following much more relaxed problem is open:

\begin{problem}[{\cite[Open Problem 15.5]{burgisser1997algebraic}}]\label{prob1}
    Prove that there is a tensor $T \in \FF^n \otimes \FF^n \otimes \FF^n$ with $\asymprank(T) > n$.
\end{problem}




It is possible that for every  %\refnote{3}{concise first use out of: \phantomsection\label{ref:3-first-concise-use},\ref{ref:3-first-concise-use},\ref{ref:3-first-concise-def},\ref{ref:3-concise-use7},\ref{ref:3-concise-def8},\ref{ref:3-concise-use10},\ref{ref:3-concise-def13}} 
tensor $T \in \FF^n \otimes \FF^n \otimes \FF^n$ we have $\asymprank(T) \leq n$ (which would in particular imply $\omega=2$) and that the image of $\asymprank$ over all tensors is simply $\NN$. This naturally leads to the (very general) question:
    What is the structure (geometric, algebraic, topological, etc) of the set
    \[
    S = \{\asymprank(T) : T \in \FF^{n_1} \otimes \FF^{n_2} \otimes \FF^{n_3}, n_1, n_2, n_3 \in \NN\}?
    \]
Is there anything we can prove about $S$ without resolving \autoref{prob1} or determining~$\omega$? Not much is known. 

One known structural result is that $S$ is closed under applying any univariate polynomial with non-negative integer coefficients \cite[Theorem 4.8]{wigderson2022asymptotic}. This statement applies in fact more generally, and in particular also to asymptotic subrank and asymptotic slice rank. This thus says that $S$ has ``many'' elements. 

Our discreteness result says that $S$ does not have ``too many'' elements, and for asymptotic tensor rank over a finite field\footnote{Or in fact over any finite set of coefficients coming from an arbitrary field, see \autoref{th:asympran-disc}.} the proof is surprisingly simple, 
and uses some basic definitions that will play an important role throughout. Let~$\FF$ be any field. Let $T \in \FF^{n_1} \otimes \FF^{n_2} \otimes \FF^{n_3}$ be an order-three tensor over $\FF$ with dimensions $(n_1, n_2, n_3)$. 
The \emph{flattenings} of $T$ are the matrices in $(\FF^{n_1} \otimes \FF^{n_2}) \otimes \FF^{n_3}$, $\FF^{n_1} \otimes (\FF^{n_2} \otimes \FF^{n_3})$ and $\FF^{n_2} \otimes (\FF^{n_1} \otimes \FF^{n_3})$, obtained by naturally grouping the tensor legs of~$T$. 
%
We say two tensors $S \in \FF^{n_1} \otimes \FF^{n_2} \otimes \FF^{n_3}$ and $T \in \FF^{m_1} \otimes \FF^{m_2} \otimes \FF^{m_3}$ are \emph{equivalent} if there are linear maps $A_\ell : \FF^{n_\ell} \to \FF^{m_\ell}$ such that $(A_1 \otimes A_2 \otimes A_3) S = T$ and there are linear maps $B_\ell : \FF^{m_\ell} \to \FF^{n_\ell}$ such that $(B_1 \otimes B_2 \otimes B_3) S = T$.
A tensor $T \in \FF^{n_1} \otimes \FF^{n_2} \otimes \FF^{n_3}$ is called \emph{concise} 
%\refnote{3}{concise definition out of: \phantomsection\label{ref:3-concise-def8},\ref{ref:3-first-concise-use},\ref{ref:3-first-concise-def},\ref{ref:3-concise-use7},\ref{ref:3-concise-def8},\ref{ref:3-concise-use10},\ref{ref:3-concise-def13}} 
if the three flattenings obtained by grouping two of the three tensor legs together, in the three possible ways, each have maximal rank. This means essentially that $T$ cannot be embedded into a smaller tensor space. Every tensor is equivalent to a concise tensor \cite[Section 14.3]{burgisser1997algebraic} (see also \autoref{lem:equiv-concise}).

Here is the proof that the asymptotic rank over finite fields is discrete:
Let $T_1, T_2, \ldots$ be any sequence of tensors such that $T_i \in \FF^{a_i} \otimes \FF^{b_i} \otimes \FF^{c_i}$ and such that $\{\asymprank(T_i) : i \in \NN\}$ is infinite. 
%takes 
%\refnote{1}{sugested to change to: "$\asymprank(T_i)$'s take"} 
%\Jshort{I changed it in a different way.}
We may assume that every~$T_i$ is concise 
%\refnote{3}{first rough definition of concise, out of several:\phantomsection\label{ref:3-first-concise-def},\ref{ref:3-first-concise-use},\ref{ref:3-first-concise-def},\ref{ref:3-concise-use7},\ref{ref:3-concise-def8},\ref{ref:3-concise-use10},\ref{ref:3-concise-def13}}
. Then $\asymprank(T_i) \geq \max \{a_i, b_i, c_i\}$. Since by assumption $\FF$ is finite, there are only finitely many tensors per format $a_i \times b_i \times c_i$, so the set of triples $\{(a_i, b_i, c_i) : i \in \NN\}$ is infinite. In particular, $\max \{a_i, b_i, c_i\}$ is unbounded, and so $\asymprank(T_i)$ is unbounded and cannot converge, which proves the claim. 

While this argument is very simple, a much more subtle argument and new technical results will be needed to deal with the other tensor parameters that we consider.

\subsubsection*{Asymptotic subrank and asymptotic slice rank}

Besides tensor rank, there are many other notions of rank of a tensor that play a role in applications, for instance the subrank, slice rank,  analytic rank~\cite{MR2773103,MR3964143,10.1215/00127094-2022-0086}, geometric rank~\cite{KopMosZui:GeomRankSubrankMaMu}, and G-stable rank~\cite{MR4471036}. We will focus here on the asymptotic subrank and asymptotic slice rank. The subrank was introduced by Strassen \cite{MR882307} in the study of matrix multiplication algorithms. The subrank $\subrank(T)$ is the size of the largest diagonal tensor that can be obtained from $T$ by taking linear combinations of the slices in the three different directions (i.e.~``Gaussian elimination'' for tensors). The slice rank was introduced by Tao \cite{tao} to give a tensor proof of the cap set problem after the first proof of Gijswijt and Ellenberg~\cite{MR3583358}. The slice rank $\slicerank(T)$ is the smallest number of tensors with flattening rank one whose sum is $T$. Tao proved that $\subrank(T) \leq \slicerank(T)$.
Recent works have shown that analytic rank, geometric rank \cite{DBLP:conf/stoc/CohenM21,10.1215/00127094-2022-0086}, and G-stable rank~\cite{MR4471036} are all equal to slice rank, up to a multiplicative constant. These results imply that the asymptotic versions of these parameters are all equal to the asymptotic slice rank, warranting our focus on it.

\emph{Slice rank method in combinatorics.} The proof of the longstanding cap set problem \cite{tao, sawin} (and other related results \cite{naslund_sawin_2017}) can be thought of as upper bounding the independence number of powers of a hypergraph, by constructing a tensor that ``fits'' on the hypergraph and then computing the slice rank of the powers of the tensor, that is, the asymptotic slice rank $\asympslicerank(T)$. Knowing (the structure of) the set
\[
S = \{\asympslicerank(T) : T \in \FF^{n_1} \otimes \FF^{n_2} \otimes \FF^{n_3}, n_1, n_2, n_3 \in \NN\}
\]
thus gives information on what bounds one can prove on the size of combinatorial objects using the slice rank method.
One of our main results is that asymptotic slice rank is discrete, not only over every finite field, but even over the complex numbers. The latter crucially requires a result from \cite{MR4495838} that characterizes asymptotic slice rank in terms of representation-theoretic objects called moment polytopes. In fact it is known by now that the four smallest values in $S$ are $0, 1, 1.889..., 2, 2.686...$ (see next section) and our result says that also the larger values are discrete.

\emph{Matrix multiplication barriers.} Besides the aforementioned combinatorial problems, the asymptotic subrank and asymptotic slice rank appear in several ``barrier results'' for matrix multiplication algorithms \cite{alman2018limits,v017a002,MR3631613}. Matrix multiplication algorithms are usually constructed by a reduction of matrix multiplication to another bilinear map, and these barrier results say what properties that intermediate bilinear map must have to obtain certain upper bounds on $\omega$, or to reach $\omega = 2$. These properties can be phrased in terms of asymptotic subrank or asymptotic slice rank. In particular, the barrier of \cite{v017a002} states that to reach $\omega = 2$, an intermediate tensor $T$ must have $\asymprank(T) = \asympsubrank(T)$, which has led to further research to find tensors with large asymptotic subrank \cite{DBLP:conf/mfcs/BlaserL20}. Our discreteness result says that asymptotic subrank is discrete over every finite field. We do not get this result over the complex numbers because the analogous representation-theoretic ingredient from above is missing here. Intriguingly, it is possible that 
\[
S = \{\asympsubrank(T) : T \in \FF^{n_1} \otimes \FF^{n_2} \otimes \FF^{n_3}, n_1, n_2, n_3 \in \NN\}
\]
equals the analogous set for asymptotic slice rank, that is, the following is open:
\begin{problem}
    Prove that asymptotic slice rank equals asymptotic subrank.
\end{problem}
We do as a side-result prove a new relation between asymptotic subrank and asymptotic slice rank (which we will discuss in more detail in the next section).

% \subsubsection*{Another story line idea}
% \begin{itemize}
%     \item It has been widely conjectured that $\omega = 2$.
%     \item Even more strongly, it is possible that for every tensor the asymptotic rank equals the largest dimension of the tensor (which would imply the above when applied to the matrix multiplication tensor). In particular, then the asymptotic rank would take only values in $\NN$ and would thus be discrete
%     \item Can we prove discreteness of asymptotic tensor rank without proving $\omega = 2$?
%     \item It turns out that we can for any finite set of coefficients (in particular over finite fields) with a relatively simple proof.
%     \item With much more elaborate techniques we extend this result to other asymptotic tensor parameters, in particular asymptotic subrank and asymptotic slice rank. Interestingly, we know that these parameters take values outside of $\NN$, yet we can prove the set of values is discrete!
%     \item Asymptotic subrank and asymptotic slice rank have played important roles in the study of matrix multiplication too and in the study of problems in combinatorics. 
%     \item The discreteness of the above parameters gives rise to the possibility of ``rounding'' or ``boosting'' (upper or lower) bounds on them to the next possible value. Moreover in the context of upper bounding the hypergraph independence number, this may lead to limitations for the asymptotic slice rank to improve on existing bounds (as was the motivation for Costa--Dalai).
% \end{itemize}

% \subsubsection*{The origins and significance of Asymptotic Subrank and Slicerank}
% Strassen, Laser method, Tao, Cap-set, SLOCC and asymptotic transformations, distillable entanglement.

% \subsubsection*{Questions about the asymptotic parameters and some answers}
% Lower bounds -- first lower bounds purely from local dimensions
% Relationship between the parameters -- first "converse" bound
% Structural: how does the image look like -- discreteness in some regimes


\subsubsection*{Previous work on discreteness of asymptotic ranks}
Several works, among which some very recent ones, have investigated notions of discreteness in the values of tensor parameters.

Strassen \cite[Lemma~3.7]{strassen1988asymptotic} proved that for any $k$-tensor $T$ over any field, the asymptotic subrank (and, as a consequence of his method, also the asymptotic slice rank) of~$T$ is equal to $0$, equal to $1$, or at least $2^{2/k}$. This result established the first ``gaps'' in asymptotic tensor parameters.
%
Costa and Dalai \cite{DBLP:journals/jcta/CostaD21} proved that, for any $k$-tensor $T$ over any field, the asymptotic slice rank of $T$ is equal to $0$, equal to $1$ or at least $2^{h(1/k)}$ where~$h$ is the binary entropy function\footnote{The binary entropy function is defined for $p \in (0,1)$ by $h(p) = - p\log_2 p - (1-p) \log_2 (1-p)$ and $h(0) = h(1) = 0$.}.
%
Christandl, Gesmundo and Zuiddam \cite{cgz} extended the result of Costa and Dalai by proving that, for any $k$-tensor $T$ over any field, the asymptotic subrank and asymptotic partition rank of $T$ are equal to~$0$, equal to~$1$ or at least~$2^{h(1/k)}$ (which is a tight bound). Additionally, they prove that for any $3$-tensor $T$ over any field, the asymptotic subrank and asymptotic slice rank of~$T$ are equal to~$0$, equal to $1$, equal to~$2^{h(1/3)}\approx 1.889$ or at least $2$. Gesmundo and Zuiddam~\cite{gesmundo2023gap} extended this result by proving that the next possible value after 2 is $\approx 2.686$.

Blatter, Draisma and Rupniewski \cite{bdr} proved that for any function on $k$-tensors, over any finite field, that is normalized and monotone the set of values that this function takes is well-ordered. The asymptotic (sub)rank and asymptotic slice rank\footnote{Over finite fields, since we do not know whether the limit $\lim_{n\to\infty}\slicerank(T^{\boxtimes n})^{1/n}$ exists in general, when we say asymptotic slice rank we will mean $\liminf_{n\to\infty} \slicerank(T^{\boxtimes n})^{1/n}$.} %
are examples of such functions. This means that the values of any such function do not have accumulation points ``from above'', but leaves open the possibility that there are accumulation points ``from below''.

Christandl, Vrana and Zuiddam \cite{MR4495838} proved that the asymptotic slice rank over the complex numbers takes only finitely many values on tensors of any fixed format, %
and thus only countably infinite many values in general. This is done by characterizing the asymptotic slice rank as an optimization over the moment polytope of the tensor and using the result that there are only finitely many such polytopes per tensor format.

Blatter, Draisma and Rupniewski \cite{https://doi.org/10.48550/arxiv.2212.12219} proved that for any ``algebraic'' tensor invariant over the complex numbers the related asymptotic parameter takes only countably many values. This implies in particular that the asymptotic (sub)rank, asymptotic slice rank, asymptotic geometric rank, and asymptotic partition rank (all over the complex numbers) take only countably many values. 
%\newpage

\subsubsection*{New results}

In this paper:
\begin{itemize}
    \item We prove two general lower bounds on the asymptotic subrank of concise 
    %\refnote{3}{concise use out of: \phantomsection\label{ref:3-concise-use7},\ref{ref:3-first-concise-use},\ref{ref:3-first-concise-def},\ref{ref:3-concise-use7},\ref{ref:3-concise-def8},\ref{ref:3-concise-use10},\ref{ref:3-concise-def13}} 
    tensors that depend only on the dimensions of the tensor. The first says that the asymptotic subrank of any concise tensor is at least the cube root of its smallest dimension. The second says that the asymptotic subrank of any ``narrow enough'' tensor (meaning that one dimension is much smaller than the others) is maximal.
    \item We use those lower bounds to prove that over any finite set of coefficients in any field the asymptotic subrank has no accumulation points (i.e.,~is discrete). We moreover prove that over any finite set of coefficients and over the complex numbers the asymptotic slice rank has no accumulation points. A much simpler argument gives that the asymptotic rank is discrete over any finite set of coefficients. 
    \item As a core ingredient for the above, we prove optimal relations among the maximal rank of any matrix in the span of the slices of a tensor, when considering slicings in the three different directions.
    \item With similar techniques, we prove an upper bound on the difference between asymptotic subrank and asymptotic slice rank.
\end{itemize}

%

%

%




%

%

%

\subsection{Discreteness of asymptotic tensor parameters}

We will now discuss our results in more detail. We begin with some basic definitions (which we expand on in \autoref{subsec:prelim}).
% \refnote{2}{fractal-structure}
% \ILnote{maybe instead: "on which we expand"?}
% \Jshort{Yes}
Let $\FF$ be any field. Let $T \in \FF^{n_1} \otimes \FF^{n_2} \otimes \FF^{n_3}$ be an order-three tensor over $\FF$ with dimensions $(n_1, n_2, n_3)$. The subrank of $T$, denoted by $\subrank(T)$, is the largest number~$r$ such that there are linear maps $L_\ell : \FF^{n_\ell} \to \FF^r$ such that $(L_1 \otimes L_2\otimes L_3)T = \sum_{i=1}^r e_i \otimes e_i \otimes e_i$. In other words, the subrank measures how much a tensor can be diagonalized. 
The flattenings of $T$ are the elements in $(\FF^{n_1} \otimes \FF^{n_2}) \otimes \FF^{n_3}$, $\FF^{n_1} \otimes (\FF^{n_2} \otimes \FF^{n_3})$ and $\FF^{n_2} \otimes (\FF^{n_1} \otimes \FF^{n_3})$, obtained by naturally grouping the tensor legs of~$T$. 
We say $T$ has slice rank one if at least one of the flattenings has rank one. The slice rank of $T$, denoted by $\slicerank(T)$ is the smallest number $r$ such that~$T$ is a sum of $r$ tensors with slice rank one.
The asymptotic subrank of $T$ is defined as $\asympsubrank(T) = \lim_{n \to \infty} \subrank(T^{\boxtimes n})^{1/n}$ where $\boxtimes$ is the Kronecker product on tensors. The limit exists and equals the supremum by Fekete's lemma, since $\subrank$ is super-multiplicative. The asymptotic slice rank of $T$ we define as $\asympslicerank(T) = \liminf_{n\to\infty} \slicerank(T^{\boxtimes n})^{1/n}$. (Over the complex numbers, it is known that the liminf can be replaced by a limit. Over other fields, however this is generally not known.) 

% Subrank and asymptotic subrank were introduced by Strassen \cite{MR882307} to study matrix multiplication algorithms, for which he introduced the asymptotic spectrum of tensors \cite{strassen1988asymptotic,wigderson2022asymptotic}. Slice rank (and its asymptotic analysis) was introduced by Tao~\cite{tao} to give a tensor-proof of the cap set conjecture. Among other applications, these parameters played a central role in recent study of matrix multiplication algorithms~\cite{blaser2013fast} and barriers~\cite{MR3631613,v017a002,alman2018limits}.

%

%



\subsubsection{Discreteness of 
asymptotic slice rank and asymptotic subrank}
%
We prove discreteness (no accumulation points) for the asymptotic slice rank and asymptotic subrank, in several regimes, as follows.

%\Jnote{I changed finite field to finite set of coefficients in the theorem below.}

\begin{theorem}\label{th:discr-finite}
    Over any finite set of coefficients in any field, the asymptotic subrank and the asymptotic slice rank each have no accumulation points.
\end{theorem}

%

\autoref{th:discr-finite} improves the result of Blatter, Draisma and Rupniewski~\cite{bdr} that the asymptotic subrank and asymptotic slice rank over any finite field have no accumulation points ``from above'', that is, have well-ordered sets of values. Indeed our result rules out all accumulation points (so also those ``from below''). We obtain the statement of \autoref{th:discr-finite}  also for asymptotic rank, but with a much shorter proof (\autoref{th:asympran-disc}).

\begin{theorem}\label{th:discr-complex}
    Over the complex numbers, the asymptotic slice rank has no accumulation points.
\end{theorem}

\autoref{th:discr-complex} improves the result of Blatter, Draisma and Rupniewski~\cite{https://doi.org/10.48550/arxiv.2212.12219} and Christandl, Vrana and Zuiddam \cite{MR4495838} that the asymptotic slice rank over the complex numbers takes only countably many values. Our result is indeed strictly stronger, as countable sets may a priori have accumulation points. %

Our results shed light on the recent results of Costa and Dalai \cite{DBLP:journals/jcta/CostaD21} and Christandl, Gesmundo and Zuiddam \cite{cgz} that found gaps between the smallest values of the asymptotic slice rank and asymptotic subrank, and answers positively (in some regimes) the question stated in \cite{cgz} asking whether the values will always be ``gapped''.

Besides the above, we prove discreteness theorems over arbitrary fields for two sub-classes of all tensors. Namely we consider the class of oblique tensors, which are the tensors whose support in some basis is an antichain, and the tight tensors, whose support in some basis can be characterized by algebraic equation (details in \autoref{sec:discreteness}). 
The set of tight tensors is a strict subset of the set of oblique tensors, which is a strict subset of all tensors. 
Both classes originate in the work of Strassen \cite{MR882307,strassen1988asymptotic,strassen1991degeneration}. Examples of tight tensors include the well-known matrix multiplication tensors. Tight tensors also play a central role in the laser method of Strassen to construct fast matrix multiplication algorithms. Our discreteness theorems in these regimes are as follows.

\begin{theorem}
    Over any field, on tight tensors,  the asymptotic subrank and asymptotic slice rank (which are equal) have no accumulation points.
\end{theorem}

\begin{theorem}
    Over any field, on oblique tensors, the asymptotic slice rank has no accumulation points.
\end{theorem}

%
%

\subsubsection{General discreteness theorem}

We prove the above discreteness theorems as an application of a general discreteness theorem that we discuss now. This general theorem gives discreteness for real-valued tensor parameters that satisfy several conditions. 
% To describe these conditions we need the notion of equivalence of tensors. We say two tensors $S \in \FF^{n_1} \otimes \FF^{n_2} \otimes \FF^{n_3}$ and $T \in \FF^{m_1} \otimes \FF^{m_2} \otimes \FF^{m_3}$ are \emph{equivalent} if there are linear maps $A_i : \FF^{n_i} \to \FF^{m_i}$ such that $(A_1 \otimes A_2 \otimes A_3) S = T$ and there are linear maps $B_i : \FF^{m_i} \to \FF^{n_i}$ such that $(B_1 \otimes B_2 \otimes B_3) S = T$.
% A tensor $T \in \FF^{n_1} \otimes \FF^{n_2} \otimes \FF^{n_3}$ is called concise \refnote{3}{concise definition out of: \phantomsection\label{ref:3-concise-def8},\ref{ref:3-first-concise-use},\ref{ref:3-first-concise-def},\ref{ref:3-concise-use7},\ref{ref:3-concise-def8},\ref{ref:3-concise-use10},\ref{ref:3-concise-def13}} if the three flattenings obtained by grouping two of the three tensor legs together, in the three possible ways, each have maximal rank. This means essentially that $T$ cannot be embedded into a smaller tensor space. 

\begin{theorem}\label{gen-no-acc-simple}
    Let $\FF$ be any field. Let $\mathcal{C}$ be any subset of tensors over~$\FF$ such that for every $T \in \mathcal{C}$ there is an $S \in \mathcal{C}$ that is concise and equivalent to $T$.
    Let~$f : \mathcal{C} \to \RR_{\geq0}$ be any function such that
    \begin{enumerate}[\upshape (i)]
    \item $f$ does not change under equivalence of tensors
    \item\label{lb-cond-intro} $f(T) \geq \asympsubrank(T)$ for every $T \in \mathcal{C}$
    %
    %
    %
    \item\label{fin-cond-intro} For every $n_1, n_2, n_3\in \NN$,  $f$ takes finitely many values on $\mathcal{C} \cap (\FF^{n_1} \otimes \FF^{n_2} \otimes \FF^{n_3})$ 
    \item For every $n_1, n_2, n_3\in \NN$, $f(T) \leq \min_\ell n_\ell$ for every $T \in \mathcal{C} \cap (\FF^{n_1} \otimes \FF^{n_2} \otimes \FF^{n_3})$.
    \end{enumerate}
    Then the set of values that $f$ takes on $\mathcal{C}$ has no accumulation points.
\end{theorem}


%

We make some remarks on condition~(\ref{fin-cond-intro}) of \autoref{gen-no-acc-simple}.
When applying \autoref{gen-no-acc-simple} to a given real-valued function $f$ on tensors, it will depend very much on the regime we are working in whether condition~(\ref{fin-cond-intro}) is trivial or not. In particular, over finite fields, condition~(\ref{fin-cond-intro}) is trivial, because then $\FF^{n_1} \otimes \FF^{n_2} \otimes \FF^{n_3}$ contains only finitely many elements. 
Another trivial situation is when~$f$ takes only integral values, as there are only finitely many integers between 0 and $\min_\ell n_\ell$ (in this case the conclusion of the theorem is also trivial). 
However, when $f$ is not integral (say $f$ is the asymptotic slice rank or asymptotic subrank) and the field is not finite (say $\FF$ is the complex numbers) condition~(\ref{fin-cond-intro}) can be very non-trivial. 
For instance, our application of \autoref{gen-no-acc-simple} to asymptotic slice rank over the complex numbers (leading to \autoref{th:discr-complex}) relies on the representation-theoretic characterization of this parameter in terms of moment polytopes \cite{MR4495838} and a result from invariant theory that there are only finitely many such polytopes per choice of~$(n_1, n_2, n_3)$.

Our proof of \autoref{gen-no-acc-simple} relies mainly on new lower bounds on the asymptotic subrank~$\asympsubrank(T)$ of concise tensors $T$, which we will discuss next. Intuitively, these lower bounds will ensure (using condition~(\ref{lb-cond-intro})) that for any infinite sequence of tensors $T_i$ the value of~$f(T_i)$ gets ``pushed up'' so much that it either cannot converge, or eventually becomes constant (when $\min_\ell n_\ell$ is bounded).

%
%
%
%
%
%
%
%
%
%
%
%

\subsection{Lower bounds on asymptotic subrank}

Having discussed our discreteness theorems, we will now discuss two results that are central in the proof of the discreteness theorems, and of independent interest. These results are about lower bounds on the asymptotic subrank. 

The general goal of these results, in the context of the proof of the general discreteness theorem, is to establish that if we have a sequence of tensors $T_i \in \FF^{a_i} \otimes \FF^{b_i} \otimes \FF^{c_i}$ for $i \in \NN$ such that the set of triples $\{(a_i, b_i, c_i) : i \in \NN\}$ is infinite, then the asymptotic subrank of these tensors must be either unbounded, or eventually constant and integral. We will discuss this in detail in the proof of the general discreteness theorem.

% \subsubsection{Previous work}

% Strassen \cite{strassen1991degeneration}, building on the work of Coppersmith and Winograd \cite{MR1056627}, introduced a method to prove (optimal) asymptotic subrank lower bounds for a subclass of structured tensors called ``tight'' tensors. This method formed an integral part of the laser method for constructing fast matrix multiplication algorithms, and similar ideas have also been applied in the context of additive combinatorics \cite{kleinberg2018growth}.
% \JBnote{This currently looks like a somewhat isolated comment. Can we embed it somewhere, for instance to contrast with our results or methods?}

\subsubsection{Concise tensors have large asymptotic subrank}

For concise tensors~$T$ we prove a cube-root lower bound on the asymptotic subrank $\asympsubrank(T)$ in terms of the smallest dimension of the tensor.

\begin{theorem}\label{th:balanced}
Let $T \in \FF^{n_1} \otimes \FF^{n_2} \otimes \FF^{n_3}$ be concise. Then %
$\asympsubrank(T) \geq (\min_\ell n_\ell)^{1/3}$.
\end{theorem}

Strassen \cite{strassen1991degeneration}, building on the work of Coppersmith and Winograd \cite{MR1056627}, introduced a method to prove (optimal) asymptotic subrank lower bounds for a subclass of structured tensors called ``tight'' tensors. This method formed an integral part of the laser method for constructing fast matrix multiplication algorithms, and similar ideas have also been applied in the context of additive combinatorics \cite{kleinberg2018growth}.
We emphasize that \autoref{th:balanced} does not rely on any special structure of the tensor~$T$, unlike previous methods like the laser method that rely on $T$ being tight.

%
We do not know whether the lower bound $(\min_\ell n_\ell)^{1/3}$ in \autoref{th:balanced} is optimal. 
The best upper bound we know is from an example of a concise tensor $T \in \FF^n \otimes \FF^n \otimes \FF^n$ such that $\asympsubrank(T)=2\sqrt{n-1}$ (\autoref{ex:strassenNullAlgebra}). %

Alternatively, \autoref{th:balanced} can be phrased without the conciseness condition if we replace the dimension $n_\ell$ by the flattening rank $\tensorrank(T_\ell)$, as follows: Let $T$ be any tensor. Then $\asympsubrank(T) \geq (\min_\ell \tensorrank_\ell(T))^{1/3}$.

%
%
%

%

For symmetric tensors we can prove the following stronger bound. (In fact, we can prove this stronger bound for a larger class of tensors which we call ``pivot--matched'', as we will explain in \autoref{sub:sqrt-assub-equal-pivots}). %
We recall that a tensor $T=\sum_{i,j,k}T_{i,j,k}\, e_i\otimes e_j \otimes e_k \in \FF^{n} \otimes \FF^{n} \otimes \FF^{n}$ is called symmetric if permuting the three tensor legs does not change the tensor, that is, for every $(i_1,i_2,i_3)\in [n]^3$ and  every permutation  $\sigma\in S_3$ it holds that $T_{i_1,i_2,i_3}=T_{i_{\sigma(1)},i_{\sigma(2)},i_{\sigma(3)}}$. 
%\ILnote{Comment 10. Fixed the above definition of symmetric tensors.}

%\refnote{concise usages from here up to page 12 included, out of: \phantomsection\label{ref:3-concise-use10}, \ref{ref:3-first-concise-use},\ref{ref:3-first-concise-def},\ref{ref:3-concise-use7},\ref{ref:3-concise-def8},\ref{ref:3-concise-use10},\ref{ref:3-concise-def13}}

\begin{theorem}\label{th:balanced-sym}
Let $T \in \FF^{n} \otimes \FF^{n} \otimes \FF^{n}$ be concise and symmetric. Then $\asympsubrank(T) \geq n^{1/2}$.
\end{theorem}

We do not know whether the lower bound $n^{1/2}$ in \autoref{th:balanced-sym} is optimal.

Again, alternatively, \autoref{th:balanced-sym} can be phrased without conciseness if the dimension~$n$ is replaced by the flattening rank $\tensorrank_1(T)$ (for symmetric tensors the three flattening ranks are equal): Let $T \in \FF^n \otimes \FF^n \otimes \FF^n$ be symmetric. Then $\asympsubrank(T)\geq \tensorrank_1(T)^{1/2}$.


%
%
%

\subsubsection{Narrow tensors have maximal asymptotic subrank}

\autoref{th:balanced} implies that if $n_1$, $n_2$ and $n_3$ all grow, and $T \in \FF^{n_1} \otimes \FF^{n_2} \otimes \FF^{n_3}$ is any concise tensor, then $\asympsubrank(T)$ must also grow. This leaves open what happens in the regime where one of the $n_i$ is constant. We will consider the ``narrow'' regime where~$n_3=c$ is constant, and one of the dimensions $n_1, n_2$ is large enough. Here we prove that the asymptotic subrank is maximal, that is, matches the upper bound $c$.

\begin{theorem}\label{th:narrow-intro}
For every integer $c \geq 2$ there is an $N(c) \in \NN$ such that for every~$n_1, n_2$ with $\max\{n_1, n_2\} > N(c)$ and every concise tensor $T \in \FF^{n_1} \otimes \FF^{n_2} \otimes \FF^c$ we have $\asympsubrank(T) = c$.
\end{theorem}

Moreover, for the case $c = 2$ we prove with a direct construction that $N(2) = 2$ and that asymptotic subrank can be replaced by subrank.

\begin{theorem}
Let $n_1,n_2 > 2$. Let $T \in \FF^{n_1} \otimes \FF^{n_2} \otimes \FF^2$ be concise. Then $\subrank(T) = 2$.
\end{theorem}

\subsubsection{Lower bound on asymptotic subrank in terms of slice rank}

Besides the aforementioned bounds on the asymptotic subrank, we use some of the same methods to prove a lower bound on the asymptotic subrank in terms of the asymptotic slice rank.

Slice rank was introduced by Tao \cite{tao}. He proved that for every tensor $T$ we have $\slicerank(T) \geq \subrank(T)$. The gap between $\slicerank(T)$ and $\subrank(T)$ can be large, namely for generic tensors $T \in \FF^n \otimes \FF^n \otimes \FF^n$ (over algebraically closed fields $\FF$) it is known that $\slicerank(T) = n$ while $\subrank(T) = \Theta(\sqrt{n})$ \cite{derksen_et_al:LIPIcs.CCC.2022.9}. 
%
It is, however, not known whether there can be a large gap between $\slicerank(T^{\boxtimes m})$ and $\subrank(T^{\boxtimes m})$ when $m$ is large. In particular, it is possible that $\lim_{n\to\infty} \subrank(T^{\boxtimes n})^{1/n} = \lim_{n\to\infty} \slicerank(T^{\boxtimes n})^{1/n}$. Strassen's work implies this equality for the subset of tight tensors \cite{strassen1991degeneration}. Over the complex numbers, the limit $\lim_{n\to\infty} \slicerank(T^{\boxtimes n})^{1/n}$ is also known to exist and has a characterization in terms of moment polytopes \cite{MR4495838}.

We prove the following:

%\ILnote{Comment 8. I changed the theorem below and its instance at the end of the paper. But we might want to keep the statement here with the asymptotic slice rank as it's more elegant, and leave only the one at the end of the paper as the general one.}
%\Jnote{Yes maybe keep the statement with asymptotic slice rank here, and put the general one in the body of the paper.}
\begin{theorem}
\label{th:q-sr-bound-intro}
Suppose the limit $\asympslicerank(T) = \lim_{n\to\infty} \slicerank(T^{\boxtimes m})^{1/m}$ exists. 
%Denote $\asympsupslicerank(T):=\limsup_{m\to\infty}\slicerank(T^{\boxtimes m})^{1/m}$. 
Then 
\[
%
\asympsubrank(T) \geq \asympslicerank(T)^{2/3}.
\]
\end{theorem}

Our proof of \autoref{th:q-sr-bound-intro} consists of proving that the border subrank of the third power of $T$ is bounded from below in terms of the slice rank of $T$, and applying a field-agnostic Flanders-type lower bound on the max-rank of Haramaty and Shpilka~\cite{haramaty2010structure}.

%

%
%
%
%
%
%

\subsection{Max-rank and min-rank of slice spans}

Our lower bounds on asymptotic subrank discussed in the previous subsection rely on results we prove about the ranks of elements in the span of the slices of a tensor. These may be of independent interest and we discuss some of them here.

\subsubsection{Max-ranks of slice spans}

%
%

%
%

Our proof of \autoref{th:balanced} relies (among other ingredients) on the notion of the max-rank of matrix subspaces, and the relation between max-ranks of matrix subspaces obtained by slicing a tensor in the three possible directions.

For any matrix subspace $\mathcal{A} \subseteq \FF^{n_1} \otimes \FF^{n_2}$, we let $\textnormal{max-rank}(\mathcal{A})$ denote the largest matrix rank of any element of $\mathcal{A}$. 
%
To any tensor $T \in \FF^{n_1} \otimes \FF^{n_2} \otimes \FF^{n_3}$ we may associate three matrix subspaces 
%\refnote{4}{typo in superscripts}\ILnote{fixed}
$\mathcal{A}_1 \subseteq \FF^{n_2} \otimes \FF^{n_3}$, $\mathcal{A}_2 \subseteq \FF^{n_1} \otimes \FF^{n_3}$ and $\mathcal{A}_3 \subseteq \FF^{n_1} \otimes \FF^{n_2}$, obtained by taking the span of the slices of $T$ in each of the three possible directions. We denote the max-ranks of these spaces by $\subrank_\ell(T) = \textnormal{max-rank}(\mathcal{A}_\ell)$, for $\ell \in [3]$. We prove that for a concise tensor $T$ the max-ranks $\subrank_\ell(T)$ cannot all be small, in the following sense.

%
%

%

% \refnote{5}{use different letters in Thm below, as $i,j,k$ are mostly used for slices indices throughout the paper}\ILnote{Maybe we should just use $1,2,3$ and add "and similarly for any permutation of $[3]$" or even without this sentence, as it is clear.}
% \Jshort{Maybe it would be good to make this change, say replace these by $a,b,c$. But then we have to do it everywhere, for instance also in \autoref{def:pivots-tensors}.}

\begin{theorem}\label{th:uncertainty-principle-n-intro}
    Let $T \in \FF^{n_1} \otimes \FF^{n_2} \otimes \FF^{n_3}$ be concise. Let $\ell_1, \ell_2, \ell_3 \in [3]$ be distinct. Suppose that $|\FF|>n_{\ell_3}$. Then
    $\subrank_{\ell_1}(T)\subrank_{\ell_2}(T) \geq n_{\ell_3}.$
\end{theorem}

We have explicit examples of families of tensors (provided later) that show how \autoref{th:uncertainty-principle-n-intro} is essentially optimal:
\begin{itemize}
\item For every $n$ that is a square, 
there is a concise tensor $T \in \FF^{n} \otimes \FF^n \otimes \FF^n$ such that for all $\ell$ we have
$\sqrt{n} \leq \subrank_\ell(T)\leq 2\sqrt{n}$ (\autoref{ex:JopsEx}). In particular, for all $\ell_1 \neq \ell_2$ we have $n\leq\subrank_{\ell_1}(T)\subrank_{\ell_2}(T)\leq4n$. 

\item For every $n$, there is a concise tensor $T \in \FF^n \otimes \FF^n \otimes \FF^n$ such that $\subrank_1(T) = \subrank_3(T) = n$ and $\subrank_2(T) = 2$, so that $\subrank_2(T) \subrank_3(T) = 2n$ (\autoref{ex:strassenNullAlgebra}).

\item For every $c$ and for every $n$  that is a multiple of $c$,
%
there is a concise tensor $T \in \FF^{n} \otimes \FF^n \otimes \FF^n$ such that
%
$\subrank_1(T)=n$, $\subrank_2(T)\leq c+1$ and $\subrank_3(T)\leq n/c+1$ (\autoref{ex:generalisedNullAlg}).
In particular,
$\subrank_2(T)\subrank_3(T)\leq  \tfrac{(c+1)}{c} n+c+1$.
\end{itemize}

It follows from \autoref{th:uncertainty-principle-n-intro} (by a straightforward argument) that $\subrank_\ell(T)$ must be large for at least two directions $\ell$, in the following sense:
% \Jnote{DA Ref 1 comment 1: Should Corollary 1.14 include a condition on the cardinality of F?}
% \Jnote{Added ``Suppose that $|\FF|>\max_i n_i$.''}
\begin{corollary}
    \label{th:rank-two-direc-intro}
    Let $T \in \FF^{n_1} \otimes \FF^{n_2} \otimes \FF^{n_3}$ be concise.
    Suppose that $|\FF|>\max_\ell n_\ell$.
    There are distinct $\ell_1, \ell_2 \in [3]$ such that $\subrank_{\ell_1}(T) \geq (\max_\ell n_\ell)^{1/2}$ and $\subrank_{\ell_2}(T) \geq (\min_\ell n_\ell)^{1/2}$.
\end{corollary}

From \autoref{th:rank-two-direc-intro}  we can prove a preliminary asymptotic subrank lower bound $\asympsubrank(T) \geq (\min_\ell n_\ell)^{1/4}$ using the (easy to prove) fact that $\asympsubrank(T)^2 \geq \subrank_{\ell_1}(T) \subrank_{\ell_2}(T)$ for any distinct ${\ell_1},{\ell_2} \in [3]$. Proving our stronger cube-root lower bound $\asympsubrank(T) \geq (\min_\ell n_\ell)^{1/3}$ of \autoref{th:balanced} requires slightly more work.

Work on max-rank (and its relation to dimension) goes back to Dieudonn\'{e} \cite{MR29360}, Meshulam \cite{meshulam1985maximal} and \cite{flanders1962spaces}. Another relevant line of work here is on commutative and non-commutative rank, which has established strong connections between max-rank and non-commutative rank \cite{cohn1995skewfields,MR2123060,derksen2018ncrk}.

\subsubsection{Min-ranks of slice spans}

Our proof of \autoref{th:narrow-intro} relies on a careful analysis of the min-rank of matrix subspaces, the relation to subrank and the behaviour of min-rank under powering.

For any matrix subspace $\mathcal{A} \subseteq \FF^{n_1} \otimes \FF^{n_2}$, we let $\minrank(\mathcal{A})$ denote the smallest matrix rank of any nonzero element of $\mathcal{A}$. We prove several properties of the min-rank, of which we give a rough outline here (the precise description we defer to later):
\begin{itemize}
    \item If the slices of a tensor have large min-rank, then the tensor has large subrank (\autoref{cla-high-rank-high-subrank}). %
    \item Any concise tensor has a slice of large rank (\autoref{cla:high-rk}).
    \item If a matrix subspace has large max-rank, then we can transform the subspace in a natural fashion such that it has large min-rank and all elements are diagonal (\autoref{claim:matrices-into-diagonal-with-large-rank}).
    \item Min-rank is super-multiplicative under tensor product, as long as at least one of the matrix subspaces is diagonal (\autoref{lem:minrank-diag-supermultiplicative}).
\end{itemize}
A careful combination of the above ingredients leads to a proof of \autoref{th:narrow-intro}.

The min-rank has been investigated before in several different contexts. Amitsur~\cite{amitsur1965generalized} used min-rank to characterize properties of rings of operators. %
Meshulam and Semrl \cite{MR2073900} used the min-rank to study properties of operator spaces. %
%
%
%
In quantum information the rank is a measure of entanglement. Spaces of bipartite states which are all entangled are spaces of matrices over the complex numbers with min-rank strictly greater than~$1$. They were investigated in \cite{wallach2002unentangled} and \cite{parthasarathey2004maxdim}. In \cite{hayden2006generic} a slightly different angle was taken, analysing random subspaces. It was shown that most random subspaces have almost maximal min-rank, which was used for superdense coding in \cite{anura2006superdense}, and was summarised in \cite{hayden2004random}.
 Generalising both of these lines of work,
in \cite{cubitt2008subspacesbounded} the question of dimension versus min-rank was addressed as number of qubits versus guaranteed entanglement in a subspace of states, and their construction is used in a follow-up paper \cite{cubitt2008counterexamples} to show that 
%\refnote{6}{missing "that" here} \ILnote{I don't think "that" is needed here}\Jshort{We can give this one to the reviewer, as either is fine.} 
there are counterexamples to the additivity of $p$-R{\'e}nyi entropies, for all $p\leq p_0$ for some small constant $p_0<1$, utilising the fact that the $0$-R{\'e}nyi entropy is the min-rank. 


%

%
%

%
%
%
%
%
%
%
%

%

%
%
%
%
%

%
%
%
%

%
%
%

%
%
%
%

%

%
%
%
%
%

%

%


%
%

%
%
%
%

%
%


%
%
%

%

%
%

%
%
%

%
%
%
%

%
%
%

%

%




%

%

%
%
%

%

%
%
%

%




%
%
%
%
%
%
%
%
%
%
%
%
%
%
%
%
%
%
%
%
%

%

%
%
%
%
%
%
%
%
%
%
%
%
%
%
%
%
%

%
%
%
%
%
%
%
%
%
%
%
%
%
%
%
%
%
%

%



%

%

%

%
%
%
%
%

%

\subsection{Tensor preliminaries}\label{subsec:prelim}

We summarize in this section some standard tensor preliminaries and terminology that we will use throughout the paper, following the papers of Strassen \cite{MR882307,strassen1988asymptotic,strassen1991degeneration} and the book of Bürgisser, Clausen and Shokrollahi~\cite{burgisser1997algebraic}.


\subsubsection*{Tensors}
%
%
Let $\FF$ be any field. For any $n \in \NN$ we denote by $e_1, \ldots, e_n$ the standard basis elements of the vector space $\FF^n$. For any $n_1, n_2, n_3 \in \NN$, we denote by $\FF^{n_1} \otimes \FF^{n_2} \otimes \FF^{n_3}$ the vector space of order-three tensors $T = \sum_{i \in [n_1]} \sum_{j \in [n_2]} \sum_{k \in [n_3]} T_{i,j,k}\, e_i \otimes e_j \otimes e_k$ with coefficients $T_{i,j,k} \in \FF$. In most of this paper we will be dealing only with order-three tensors, so instead of order-three tensor we will just say tensor. 

%
\subsubsection*{Restriction and equivalence} 
We will be using a natural preordering on tensors called restriction, which is defined as follows. For any two tensors $T \in \FF^{n_1} \otimes \FF^{n_2} \otimes \FF^{n_3}$ and $S \in \FF^{m_1} \otimes \FF^{m_2} \otimes \FF^{m_3}$, we say $T$ \emph{restricts to} $S$ and write $T \geq S$ if there are linear maps $L_\ell : \FF^{n_\ell} \to \FF^{m_\ell}$ such that $(L_1 \otimes L_2 \otimes L_3)T = S$. We say $S$ and $T$ are \emph{equivalent} if $T \geq S$ and $S \geq T$. Most tensor properties we will study are invariant under equivalence and monotone under restriction (see \autoref{lem:monotone-restrict}). Finally, we say $S, T \in \FF^{n_1} \otimes \FF^{n_2} \otimes \FF^{n_3}$ are \emph{isomorphic} if $(L_1 \otimes L_2 \otimes L_3)T = S$ for invertible linear maps $L_\ell : \FF^{n_\ell} \to \FF^{n_\ell}$.
%

\subsubsection*{Flattenings and flattening ranks}
Let $T \in V_1 \otimes V_2 \otimes V_3$. Grouping together $V_2$ and $V_3$ gives an element  $T_{1} \in V_1 \otimes (V_2 \otimes V_3)$, and we similarly obtain $T_2$ and $T_3$. %(for any choice of distinct ${\ell_1},{\ell_2},{\ell_3} \in [3]$). 
%We may think of $T_{\ell_1}$ as an $n_{\ell_1} \times (n_{\ell_2}\cdot n_{\ell_3})$ matrix. We call the $T_{\ell_1}$ 
We may think of the $T_\ell$ as matrices and we call them 
the \emph{flattenings} of $T$.
Let $\tensorrank_{\ell}(T)$ ve the matrix rank of $T_{\ell}$. We call $\tensorrank_1(T)$, $\tensorrank_2(T)$ and $\tensorrank_3(T)$ 
%\refnote{7}{missing plural " 's"}\ILnote{added it}\Jshort{I don't like the plural s so much after symbols so I solved it differently} 
the \emph{flattening ranks} of $T$.
%Equivalently, we may slice $T$ into matrices $A_i = (T_{i,j,k})_{j,k} \in \FF^{n_2} \otimes \FF^{n_3}$, $B_j = (T_{i,j,k})_{i,k} \in \FF^{n_1} \otimes \FF^{n_3}$ and $C_k = (T_{i,j,k})_{i,j} \in \FF^{n_1} \otimes \FF^{n_2}$
The triple $(\tensorrank_1(T), \tensorrank_2(T), \tensorrank_3(T))$ is often called the \emph{multilinear rank} of $T$. 


\subsubsection*{Conciseness} 
%\refnote{3}{concise actual definition, out of: \phantomsection\label{ref:3-concise-def13}, \ref{ref:3-first-concise-use},\ref{ref:3-first-concise-def},\ref{ref:3-concise-use7},\ref{ref:3-concise-def8},\ref{ref:3-concise-use10},\ref{ref:3-concise-def13}}
Let $T \in \FF^{n_1} \otimes \FF^{n_2} \otimes \FF^{n_3}$. From the definition of flattening rank it follows that $\tensorrank_\ell(T) \leq n_\ell$ for every $\ell \in [3]$. We call $T$ \emph{concise} if $\tensorrank_\ell(T) = n_\ell$  for every~$\ell \in [3]$. Conciseness essentially says that the tensor cannot fit in a smaller tensor space. The property of being concise is very mild in the following sense. We call a tensor $S$ a \emph{subtensor} of $T = \sum_{i,j,k} T_{i,j,k}\, e_i \otimes e_j \otimes e_k$ if $S$ is of the form $S = \sum_{i\in I,j\in J,k\in K} T_{i,j,k}\, e_i \otimes e_j \otimes e_k$ for some subsets $I,J,K$.

\begin{lemma}\label{lem:equiv-concise}
    Every tensor is equivalent to a concise tensor, which is unique up to equivalence. Concretely, every tensor is equivalent to a subtensor that is concise.
\end{lemma}
\begin{proof}
    The first statement is a standard fact \cite[Section 14.3]{burgisser1997algebraic}. For the second statement, let $T = \sum_{i,j,k} T_{i,j,k}\, e_i \otimes e_j \otimes e_k \in \FF^{n_1} \otimes \FF^{n_2} \otimes \FF^{n_3}$. Let $A_i = (T_{i,j,k})_{j,k} \in \FF^{n_2} \otimes \FF^{n_3}$ for $i \in [n_1]$ be the 1-slices of~$T$. Then one verifies that the dimension of the span of the~$A_i$ equals the flattening rank $r_1=\tensorrank_1(T)$. Suppose $A_1, \ldots, A_r$ are linearly independent, and let $S$ be the tensor with these matrices as 1-slices. Then $S$ is a subtensor of $T$ so clearly $T \geq S$. Also, since the matrices $A_{r+1}, \ldots, A_{n_1}$ are in the span of the matrices $A_1, \ldots, A_r$ we have that $S \geq T$. Since $S$ and $T$ are thus equivalent, their flattening ranks are equal. Repeating this construction for the other two directions gives the claimed subtensor.
    %
    %Zij T een n1 x n2 x n3 tensor. Zij A_1, ..., A_{n1} de 1-slices van T, dus T = \sum_{i=1}^{n1} e_i \otimes A_i. De flattening rank R1(T) is de dimensie van de span van de A_i. Stel A_1 is in de span van A_2, ..., A_{n1}. Laat S = \sum_{i=2}^{n1} e_i \otimes A_i. Dit is een subtensor van T. Er geldt T \geq S en S \geq T. De eerste restrictie T \geq S is duidelijk, want S is een subtensor van T. De tweede restrictie is makkelijk in te zien de slices van S zijn A_2, ..., A_{n1} en A_1 is in de span van die matrices. Met andere woorden S en T zijn equivalent, en dus zijn hun flattening ranks hetzelfde (want flattening ranks zijn monotoon onder restrictie). Met deze procedure kunnen we een tensor U vinden die een subtensor is van T, equivalent is aan T, en concise is.
\end{proof}

Not every space $\FF^{n_1} \otimes \FF^{n_2} \otimes \FF^{n_3}$ contains concise tensors, as can be seen from the following implication (which is in fact known to be an equivalence \cite[Theorem~2]{MR2776439}):

\begin{lemma}\label{cl:concise-formats}
For any $n_1,n_2,n_3\in\NN$, if there is a concise tensor in $\FF^{n_1} \otimes \FF^{n_2} \otimes \FF^{n_3}$, then $n_1 \leq n_2\cdot n_3$, $n_2 \leq n_1\cdot n_3$, and $n_3 \leq n_1\cdot n_2$.
\end{lemma}
\begin{proof}
    Suppose $T \in \FF^{n_1} \otimes \FF^{n_2} \otimes \FF^{n_3}$ is concise. Then $n_{\ell_1} = \rank(T_{\ell_1})$ where $T_{\ell_1}$ is the previously defined flattening of $T$. Since $T_{\ell_1}$ is an $n_{\ell_1} \times (n_{\ell_2}\cdot n_{\ell_3})$ matrix, we have $\tensorrank(T_{\ell_1}) \leq n_{\ell_2} \cdot n_{\ell_3}$. Thus $n_{\ell_1} \leq n_{\ell_2} \cdot n_{\ell_3}$ for every choice of distinct ${\ell_1},{\ell_2},{\ell_3} \in [3]$.
\end{proof}
%
%
%
%

%
%
%
%
%
%
%
%

\subsubsection*{Unit tensors, tensor rank, subrank, and slice rank}
For every $r \in \NN$ we define the tensor $\langle r \rangle = \sum_{i=1}^r e_i \otimes e_i \otimes e_i \in \FF^r \otimes \FF^r \otimes \FF^r$. We call~$\langle r \rangle$ the \emph{unit tensor} of rank $r$.
We say a tensor has \emph{tensor rank one} if it is of the form $u \otimes v \otimes w$ for nonzero vectors $u,v,w$.
The tensor rank of a tensor $T$ is the smallest number $r$ such that $T$ can be written as a sum of $r$ tensors with tensor rank one. 
%
We denote tensor rank by $\tensorrank(T)$. Equivalently, the tensor rank $\tensorrank(T)$ is the smallest number $r$ such that $T \leq \langle r\rangle$. The \emph{subrank} of $T$ is the largest number $r$ such that $\langle r\rangle \leq T$. We denote subrank by $\subrank(T)$. 
%\ILnote{maybe we should add: "For $\ell\in[3]$, we denote by $\subrank_\ell(T)$ the maximal rank of any linear combination of the $\ell$-slices of $T$ (see \autoref{subsec:maxrank}". without it -- we haven't defined it yet but it appears in the lemma below.} 
We say a tensor has \emph{slice rank one} if any of its flattening ranks equals~one. The slice rank of a tensor $T$ is the smallest number $r$ such that $T$ can be written as a sum of $r$ tensors with slice rank one.
%
Subrank, slice rank, the flattening ranks and tensor rank satisfy $\subrank(T) \leq \slicerank(T) \leq \tensorrank_\ell(T) \leq \tensorrank(T)$ for every $\ell\in [3]$.

It is also easy to see that these measures are monotone under restrictions:

\begin{lemma}\label{lem:monotone-restrict}
    For any two tensors $T$ and $S$, if $T\geq S$, then we have $\tensorrank(T)\geq \tensorrank(S)$, $\slicerank(T)\geq \slicerank(S)$, $\subrank(T)\geq \subrank(S)$, and for every $\ell\in[3]$, %$\subrank_\ell(T)\geq \subrank_\ell(S)$ and 
    $\tensorrank_\ell(T)\geq \tensorrank_\ell(S)$.
\end{lemma}


\subsubsection*{Kronecker product and asymptotic tensor parameters}
%
Asymptotic tensor parameters are defined by ``regularizing'' a tensor parameter over large powers under the Kronecker product. For any two tensors 
\begin{align*}
T &= \sum_{i,j,k} T_{i,j,k}\, e_i \otimes e_j \otimes e_k \in \FF^{n_1} \otimes \FF^{n_2} \otimes \FF^{n_3}\\
S &= \sum_{a,b,c} S_{a,b,c} \, e_a\otimes e_b \otimes e_c \in \FF^{m_1} \otimes \FF^{m_2} \otimes \FF^{m_3}
\end{align*}
the Kronecker product $T \boxtimes S \in \FF^{n_1 m_1} \otimes \FF^{n_2 m_2} \otimes \FF^{n_3 m_3}$ is defined as
\[
T \boxtimes S = \sum_{i,j,k} \sum_{a,b,c} T_{i,j,k} S_{a,b,c}\, (e_i \otimes e_{a}) \otimes (e_j \otimes e_{b}) \otimes (e_k \otimes e_{c}).
\]
We define $T^{\boxtimes n}$ as the product of $n$ copies of $T$. The asymptotic tensor rank is defined as $\asymprank(T) = \lim_{n \to \infty} \tensorrank(T^{\boxtimes n})^{1/n}$, which, because of sub-multiplicativity of $\tensorrank$, equals the infimum $\inf_{n} \tensorrank(T^{\boxtimes n})^{1/n}$ (by Fekete's lemma). The asymptotic subrank is defined as $\asympsubrank(T) = \lim_{n \to \infty} \subrank(T^{\boxtimes n})^{1/n}$, which, because of super-multiplicativity of the subrank $\subrank$, equals the supremum $\sup_{n} \subrank(T^{\boxtimes n})^{1/n}$. Over the complex numbers, the asymptotic slice rank is defined as $\asympslicerank(T) = \lim_{n\to\infty} \slicerank(T^{\boxtimes n})^{1/n}$. By the characterization in \cite{MR4495838} this limit exists. Over other fields, we generally do not know whether the limit $\lim_{n\to\infty} \slicerank(T^{\boxtimes n})^{1/n}$ exists, except for oblique tensors. When we do not know whether the limit exists we will instead consider the liminf and refer to this as the ``asymptotic slice rank''.\footnote{The results we have over finite fields for asymptotic slice rank hold equally well using liminf or limsup.}

%

%
%
%
%
%
%
%
%
%
%
%
%
%
%
%
%
%
%
%

\section{Max-rank of slice spans and asymptotic subrank lower bound}\label{sec:max-rank}

In this section we will prove a lower bound on the asymptotic subrank.
The main goal is to obtain a lower bound on the asymptotic subrank of concise tensors that only depends on the dimensions of the tensor and that moreover grows as the dimensions grow. Unlike previous lower bound methods (like the laser method mentioned in the introduction) our lower bound does not rely on any special structure of the tensor. Indeed the property of being concise is very mild as any tensor can be made concise by embedding it in the smallest tensor space that it fits in.  %

Concretely we will prove that the asymptotic subrank of a concise tensor is at least the cube root of the smallest dimension of the tensor. We will give two proofs of this in this paper, one in this section and another one in \autoref{subsec:cube-root-via-pivots} (which uses different ideas which we then use to get stronger results for symmetric tensors).

%\refnote{9}{The following paragraph sounds like a repetitions of earlier things. also \phantomsection\label{ref:9-repetition-21}\ref{ref:9-repetition-22}}
The proof of the asymptotic subrank lower bound in this section relies on the notion of max-rank of a matrix subspace (\autoref{subsec:maxrank}). %Any tensor naturally defines three matrix subspaces by slicing the tensor in the three different directions and taking spans. 
We prove that there is a strong relation between the max-ranks of slice spans of any tensor. Namely, the product of every pair of them is large (\autoref{subsec:QiQj}). This result, combined with known facts about matrix multiplication tensors, leads to the asymptotic subrank bound (\autoref{subsec:lb-mamu}).

%
%

%

%

\subsection{Max-rank of matrix subspaces and slice spans of a tensor}\label{subsec:maxrank}

%
    %
    We denote by $\FF^{n_1} \otimes \FF^{n_2}$ the space of $n_1 \times n_2$ matrices over $\FF$.
    We define the \emph{max-rank} of any matrix subspace $\mathcal{A} \subseteq \FF^{n_1} \otimes \FF^{n_2}$ as
    %
    \[
\maxrank(\mathcal{A}):=\max\{\rank(A):A\in\mathcal{A}\}.
    \]
    %
    %
    %
    %
%
%
For any tensor $T = \sum_{i,j,k} T_{i,j,k}\, e_i \otimes e_j \otimes e_k \in \FF^{n_1} \otimes \FF^{n_2} \otimes \FF^{n_3}$ we define the matrices %
%
%
%
%
%
%
\begin{align*}
T^{(1)}_i &= \sum_{j,k} T_{i,j,k}\, e_j \otimes e_k \quad (i \in [n_1])\\
T^{(2)}_j &= \sum_{i,k} T_{i,j,k}\, e_i \otimes e_k \quad (j \in [n_2])\\
T^{(3)}_k &= \sum_{i,j} T_{i,j,k}\, e_i \otimes e_j \quad (k \in [n_3]).
\end{align*}
For every $\ell \in [3]$, we call $\{T^{(\ell)}_i : i \in [n_\ell] \}$ the \emph{$\ell$-slices} of $T$. %
%
%
%
%
%
%
%
%
%
%
%
%
%
%
%
%
%
%
%
%
%
%
%
%
%
%
%
%
%
%
%
%
%
%
%
%
%
%
%
%
%
%
We define the $\ell$-slice span of $T$ as the matrix subspace
$\mathcal{A}_\ell = \linspan \{T^{(\ell)}_i : i \in [n_\ell] \}$
and denote its max-rank by
\[
\subrank_\ell(T) = \maxrank(\mathcal{A}_\ell).
\]
%
%
We note that the flattening rank $\tensorrank_\ell(T)$ defined earlier, is equal to the dimension of the subspace $\mathcal{A}_\ell$. 

%



%
%



%
%

%
%
%
%
%
%
%
%

%
%
%

%
%
%
%
%
%
%
%
%
%
%
%
%
%
%
%

%

%
%
%
%
%
%
%
%
%
%
%
%
%




\subsection{Lower bound on products of max-ranks of slice spans}\label{subsec:QiQj}

%

%

%

%

%\refnote{9}{The following paragraph sounds like a repetitions of earlier things. also \phantomsection\label{ref:9-repetition-22}\ref{ref:9-repetition-21}}
In this section we prove the aforementioned fundamental property of the max-ranks of the slice spans of a tensor, showing that they cannot all be small. This will be a crucial ingredient for proving the asymptotic subrank lower bound afterwards.

\begin{theorem}\label{th:uncertainty-principle-n}
    Let $T \in \FF^{n_1} \otimes \FF^{n_2} \otimes \FF^{n_3}$ be concise. Let ${\ell_1}, {\ell_2}, {\ell_3} \in [3]$ be distinct. Suppose that $|\FF|>n_{\ell_3}$. Then 
    \[
    \subrank_{\ell_1}(T)\subrank_{\ell_2}(T) \geq n_{\ell_3}.
    \]
\end{theorem}

%

%

%

A corollary of this is that there are two distinct directions with large max-rank:

\begin{corollary} %
\label{cor:twoLargeRankDirs}
    Let $T\in\FF^{n_1}\otimes\FF^{n_2}\otimes\FF^{n_3}$ be concise.
    Suppose $n_1 \leq n_2 \leq n_3$. Suppose that $|\FF| > n_3$. 
    %
    Then there are distinct ${\ell_1},{\ell_2}\in[3]$ such that
    \begin{align*}\subrank_{\ell_1}(T) &\geq \sqrt{n_3},\\
    \subrank_{\ell_2}(T) &\geq \sqrt{n_1}.
    \end{align*}
\end{corollary}
%
\begin{proof}%
    %
    We have by \autoref{th:uncertainty-principle-n} that
    \begin{align*}
    \subrank_1(T) \subrank_2(T) &\geq n_3\\
    \subrank_2(T) \subrank_3(T) &\geq n_1\\
    \subrank_1(T) \subrank_3(T) &\geq n_2 \geq n_1.
    \end{align*}
    Thus one of $\subrank_1(T)$ and $\subrank_2(T)$ is at least $\sqrt{n_3}$, and either $\subrank_3(T)$ is at least~$\sqrt{n_1}$ or both $\subrank_1(T)$ and $\subrank_2(T)$ are at least $\sqrt{n_1}$.
\end{proof}

%

%
%
%
%
%
%
%
%
%
%
%
%
%
%

%
%
%
%
%
%
%
%
%
%
%
%
%
%
%
%
%
%
%
%
%
%
%
%
%
     
%

%
%

%
%
%
%
%
%
%

%

%

%
%
%

%

%
%

%
%
%
%

%
%
%
%
%


%
%
%
%
%
%
%
%
%

%
%

%

%
We will now work towards the proof of \autoref{th:uncertainty-principle-n}.

For matrices $A \in \FF^{n\times m}$ and $B \in \FF^{n \times k}$ we denote by $[A; B] \in \FF^{n \times (m+k)}$ the concatenation of $A$ and $B$. For any matrix $A$ we denote by $A|_{[d]}$ the matrix consisting of the first $d$ columns of $A$. We denote by $\col(A)$ the span of the columns of $A$.

%\ILnote{Comment 9. Fixed below the instances of $[A;(BU)|_{[s]}$ which all missed the closing square bracket (missed it, and not placed the restriction inside, as the comment says).}
\begin{lemma}\label{lem:colspan}
    Let $A \in \FF^{n\times m}$ and $B \in \FF^{n \times k}$. There is a nonzero %$k^2$-variable 
    polynomial $p\in \FF[x_{1,1},x_{1,2},,\ldots,x_{i,j},\dots,x_{k,k}]$ of degree $s = \rank([A;B]) - \rank(A)$ such that for any matrix $U \in \FF^{k\times k}$ satisfying $p(U)\neq 0$, we have that
    \[
    \col([A;B]) = \col([A; (BU)|_{[s]}]).
    \]
\end{lemma}
\begin{proof}
    Let $r = \rank([A;B])$. Then, there is an $r\times r$ submatrix of $[A;B]$ with rank~$r$ induced by $r-s$ columns of~$A$ and~$s$ columns of~$B$.
    Let~$R\subseteq [n]$ be the set of row indices of this submatrix, and let $C_A\subseteq[m], C_B\subseteq \{m+1,\dots,m+k\}$ be the respective column indices of~$A$ and~$B$ in this submatrix.
    Let $p(X)$ be the determinant of the submatrix of $[A;BX]$ induced by the rows in~$R$ and columns in $C_A\cup \{m+1,\dots,m+s\}$, where $X = (x_{i,j})_{i,j=1}^k$ is a matrix of variables.
    
    We claim that~$p$ is nonzero. There is a permutation matrix~$V\in \FF^{k \times k}$ that swaps the columns in~$C_B$ with the columns in $\{m+1,\dots, m+s\}$. Since the submatrix of $[A;BV]$ induced by the rows in~$R$ and columns in $C_A\cup \{m+1,\dots,m+s\}$ is then of full rank, it follows that~$p(V)\ne 0$. 
        
    Let $U \in \FF^{k \times k}$ be a matrix such that $p(U) \neq 0$. 
    Then we have $\rank([A; (BU)|_{[s]}]) \geq r$. From $\col([A; (BU)|_{[s]}])\subseteq \col([A;B])$ and $\rank([A;B]) = r$, we find the equality $\col([A;B]) = \col([A; (BU)|_{[s]}])$.
\end{proof}

% \refnote{10}{Wouldn't it be more intuitive to merge 2.4-2.5 into one lemma where you directly add $s_i$ columns at a time?}

% \JBnote{The proofs were merged and only the second lemma remains, as suggested.}

% \begin{lemma}\label{lem:schwartz-zippel}
%     Let $A_1, \ldots, A_\ell \in \FF^{n \times k}$ be matrices. Suppose $|\FF| > \rank([A_1; \cdots; A_\ell])$. Let 
%     \[
%     s_i = \rank([A_1; \cdots ; A_{i}]) - \rank([A_1; \cdots ; A_{i-1}]).
%     \]
%     There is a matrix $U \in \FF^{k \times k}$ such that for all $i \in [\ell]$ we have 
%     \[
%     \col([A_1 ; A_2 ; \cdots ; A_{i-1} ; (A_i U)|_{[s_i]}]) = \col([A_1; \cdots ; A_i]).
%     \]
% \end{lemma}
% \begin{proof}
% From Lemma~\ref{lem:colspan}, we get that for each $i\in [\ell]$, there exists a nonzero~$k^2$-variable polynomial $p_i(X)\in \FF[x_{1,1},x_{1,2},\dots,x_{k,k}]$ of degree~$s_i$ such that for any $U\in \FF^{k\times k}$ satisfying~$p_i(U)\ne 0$,
% \[
%     \col([A_1 ; A_2 ; \cdots ; A_{i-1} ; (A_i U)|_{[s_i]}]) = \col([A_1; \cdots ; A_i]).
%     \]
% Let $q(X)=  p_1(X)\cdots p_\ell(X)$.
% This is a nonzero polynomial of degree $s_1 + \cdots +s_\ell = \rank([A_1; \cdots; A_\ell])<|\FF|$.
% It follows from the Schwartz--Zippel lemma \cite{DEMILLO1978193, schwartz1980fast, zippel1979probabilistic} that there is a~$U\in \FF^{k\times k}$ such that~$q(U) \ne 0$.
% This is to say that for all~$i\in[\ell]$ simultaneously,~$p_i(U)\ne 0$, which implies the result.
% \end{proof}

\begin{lemma}\label{lem:max-rank-induction}
    Let $A_1, \ldots, A_m \in \FF^{n \times k}$ be matrices. Suppose $|\FF| > \rank([A_1; \cdots; A_m])$. Let 
    \[
    s_i = \rank([A_1; \cdots ; A_{i}]) - \rank([A_1; \cdots ; A_{i-1}]).
    \]
    There is a matrix $U \in \FF^{k \times k}$ such that
\begin{equation}\label{eq:cspan2}
    \col([(A_1 U)|_{[s_1]} ; (A_2 U)|_{[s_2]} ; \cdots ; (A_m U)|_{[s_m]}]) = \col([A_1; \cdots ; A_m]).
\end{equation}
\end{lemma}
% \begin{proof}
% %
% Let~$U\in \FF^{k\times k}$ be a matrix as in Lemma~\ref{lem:schwartz-zippel}. We give a proof by induction.
% The base case is $\col((A_1 U)|_{[s_1]}) = \col(A_1)$,
% which follows directly from \autoref{lem:schwartz-zippel}.
% For the induction step, we assume
% \begin{equation}\label{eq:ih}
%     \col([A_1 ; A_2 ; \cdots ; A_i]) = \col([(A_1U)|_{[s_1]} ; \cdots ; (A_iU)|_{[s_i]}]).
% \end{equation}
% We now use that if $\col(A) = \col(B)$, then $\col([A ; C]) = \col([B;C])$ holds for any matrices $A,B,C$ with the same number of rows. Applying this to \eqref{eq:ih} gives
% \begin{multline}\label{eq:x}
%     \col([A_1 ; A_2 ; \cdots ; A_i; (A_{i+1} U)|_{[s_{i+1}]}])\\ = \col([(A_1U)|_{[s_1]} ; \cdots ; (A_iU)|_{[s_i]}; (A_{i+1} U)|_{[s_{i+1}]}]).
% \end{multline}
% From \autoref{lem:schwartz-zippel} we know that 
% \begin{equation}\label{eq:sz}
%     \col([A_1; \cdots ; A_{i+1}]) = \col([A_1 ; A_2 ; \cdots ; A_i ; (A_{i+1} U)|_{[s_{i+1}]}]).
% \end{equation}
% %
% %
% Then \eqref{eq:x} and \eqref{eq:sz} together give
%     \[
%     \col([A_1 ; A_2 ; \cdots ; A_i; A_{i+1}]) = \col([(A_1U)|_{[s_1]} ; \cdots ; (A_iU)|_{[s_i]}; (A_{i+1}U)|_{[s_{i+1}]}]).
%     \]
% This proves the claim by induction.
% \end{proof}


\begin{proof}
We begin by showing that there exists a matrix $U \in \FF^{k \times k}$ such that for all $i \in [m]$ we have 
\begin{equation}\label{eq:cspan1}
    \col([A_1 ; A_2 ; \cdots ; A_{i-1} ; (A_i U)|_{[s_i]}]) = \col([A_1; \cdots ; A_i]).
\end{equation}
From Lemma~\ref{lem:colspan}, we get that for each $i\in [m]$, there exists a nonzero~$k^2$-variable polynomial $p_i(X)\in \FF[x_{1,1},x_{1,2},\dots,x_{k,k}]$ of degree~$s_i$ such that for any $U\in \FF^{k\times k}$ satisfying~$p_i(U)\ne 0$,
\[
    \col([A_1 ; A_2 ; \cdots ; A_{i-1} ; (A_i U)|_{[s_i]}]) = \col([A_1; \cdots ; A_i]).
    \]
Let $q(X)=  p_1(X)\cdots p_m(X)$.
This is a nonzero polynomial of degree $s_1 + \cdots +s_m = \rank([A_1; \cdots; A_m])<|\FF|$.
It follows from the Schwartz--Zippel lemma \cite{DEMILLO1978193, schwartz1980fast, zippel1979probabilistic} that there is a~$U\in \FF^{k\times k}$ such that~$q(U) \ne 0$.
This is to say that for all~$i\in[m]$ simultaneously,~$p_i(U)\ne 0$, which implies the existence of the desired matrix~$U$.

Let~$U\in \FF^{k\times k}$ be a matrix so that~\eqref{eq:cspan1} holds for each $i\in [m]$.
Next, we show by induction on~$i\in[m]$ that~$U$ also satisfies~\eqref{eq:cspan2}.
%Let~$U\in \FF^{k\times k}$ be a matrix as in Lemma~\ref{lem:schwartz-zippel}. We give a proof by induction.
The base case is $\col((A_1 U)|_{[s_1]}) = \col(A_1)$,
which follows directly.
For the induction step, we assume
\begin{equation}\label{eq:ih}
    \col([A_1 ; A_2 ; \cdots ; A_i]) = \col([(A_1U)|_{[s_1]} ; \cdots ; (A_iU)|_{[s_i]}]).
\end{equation}
We now use that if $\col(A) = \col(B)$, then $\col([A ; C]) = \col([B;C])$ holds for any matrices $A,B,C$ with the same number of rows. Applying this to \eqref{eq:ih} gives
\begin{equation}\label{eq:x}
    \col([A_1 ; A_2 ; \cdots ; A_i; (A_{i+1} U)|_{[s_{i+1}]}]) = \col([(A_1U)|_{[s_1]} ; \cdots ; (A_iU)|_{[s_i]}; (A_{i+1} U)|_{[s_{i+1}]}]).
\end{equation}
From \eqref{eq:cspan1} we know that 
\begin{equation}\label{eq:sz}
    \col([A_1; \cdots ; A_{i+1}]) = \col([A_1 ; A_2 ; \cdots ; A_i ; (A_{i+1} U)|_{[s_{i+1}]}]).
\end{equation}
%
%
Then \eqref{eq:x} and \eqref{eq:sz} together give
    \[
    \col([A_1 ; A_2 ; \cdots ; A_i; A_{i+1}]) = \col([(A_1U)|_{[s_1]} ; \cdots ; (A_iU)|_{[s_i]}; (A_{i+1}U)|_{[s_{i+1}]}]).
    \]
This proves the claim by induction.
\end{proof}



\begin{proof}[Proof of \autoref{th:uncertainty-principle-n}]
We will prove $n_1 \leq \subrank_2(T) \subrank_3(T)$.
Let $A_1, \ldots, A_{n_3} \in \FF^{n_1 \times n_2}$ be the $3$-slices of $T$.
Let~$U$ be as in \autoref{lem:max-rank-induction}.
Let 
\[
s_i = \rank([A_1; \cdots ; A_{i}]) - \rank([A_1; \cdots ; A_{i-1}]).
\]
Conciseness of~$T$ gives that $s_1 + \cdots + s_{n_3} = n_1$.
It follows from this and \autoref{lem:max-rank-induction} that the columns of the submatrices $(A_iU)|_{[s_i]}$ are linearly independent.
This implies that for each $i\in[n_3]$, 
there is a 3-slice of rank at least $s_i$,
%
and in particular $\subrank_3(T)\geq\max_i s_i$. %
%
Let $S$ be the tensor with 3-slices $A_1U, A_2U, \ldots, A_{n_3}U$. Then $\subrank_2(S) \leq \subrank_2(T)$ since the 2-slices of $S$ are linear combinations of the 2-slices of $T$. 
The matrix $M = [(A_1U)|_{[1]};\cdots ; (A_{n_3}U)|_{[1]}]$ consisting of every first column of the 3-slices of $S$ equals the first 2-slice of~$S$. Since the linearly independent columns in the 3-slices of~$S$ are flushed to the left, we find that $M$ has at least $|\{i\in [n_3]: s_i\ne 0\}|$ linearly independent columns, and thus $\rank(M) \geq |\{i\in [n_3]: s_i\ne 0\}|$,
%
%
which gives a lower bound on~$\subrank_2(S)$ and thus on $\subrank_2(T)$.
Since
\[
n_1  = s_1 + \cdots + s_{n_3}\leq |\{i\in [n_3]: s_i\ne 0\}|\, \max_i s_i \leq \subrank_2(T) \subrank_3(T),
\]
the result follows. %\ILnote{I think we should add: "for $\ell_1=2,\ell_2=3,\ell_3=1$ and similarly for other choices." and add to the last equation above the inequality $\leq\subrank_2(T)\subrank_3(T)$}.
%
%
%
%
%
%
%
%
%
%
%
\end{proof}

%
%
%

%
%

\newcommand{\yesTextColour}{green!60!black}
\newcommand{\noTextColour}{red!60!black}
%
%
%
%
%
%

%

%
%
%

%
%
%

%
%
%
%
%
%
%
%

\newcommand{\yesCellColour}{green!20}
\newcommand{\noCellColour}{red!20}
%
%
%
%
%
%
%
%
%
%
%
%
%
%
%
%

%
%
%

%
%

%
%

%

%
%

%

%
%
%

%
%
%
%
%
%
%

%
%
%
%

%
%
%
%
%
%
%
%
%

%
%
%
%
%
%

%
%
%
%
%
%
%
%

%
%
%
%
%
%
%



%
%
%
%
%
%
%
    
%

%
%
%
%
%
%
%
%
%
%
    
%
%
    
%
%
%
%
%
%
%
%
%
%

%
%

%

%
%

%
    %
    %
    %
    %
%

\subsection{Examples of optimality of max-rank bounds}
The following examples show that the bound in \autoref{th:uncertainty-principle-n} is tight. In particular, the right hand side cannot be super-linear in $n_k$, and in fact the leading coefficient $1$ is optimal. %

\begin{example}[{Null-algebra \cite[page~168]{strassen1991degeneration}}] \label{ex:strassenNullAlgebra}
%
%
%
%
For every $n \in \NN$,
there is a concise tensor $T \in \FF^n \otimes \FF^n \otimes \FF^n$ such that $\subrank_1(T) = \subrank_3(T) = n$ and $\subrank_2(T) = 2$. In particular, $\subrank_2(T) \subrank_3(T) = 2n$.
%
%
    %
    To define this tensor, let
    \[T :=e_1\otimes e_1\otimes e_1 + \sum_{i=2}^n e_1\otimes e_i\otimes e_i + \sum_{i=2}^ne_i\otimes e_i\otimes e_1.\]
    We see directly that $\subrank_1(T) = \subrank_3(T) = n$. 
    Because the non-zero elements of the $2$-slices are confined to the first column and the first row, we have $\subrank_2(T)=2$. We also mention that for all $n\geq5$ the tensor $T$ has asymptotic subrank $\asympsubrank(T)=2\sqrt{n-1}$ \cite[page~169]{strassen1991degeneration}. %
\end{example}

\begin{example}%
    \label{ex:generalisedNullAlg}
    This example extends \autoref{ex:strassenNullAlgebra} (by taking $c=1$).
    Let $c$ be a constant and let $n$  be any multiple of $c$.
    %
    There is a concise tensor $T \in \FF^{n} \otimes \FF^n \otimes \FF^n$ such that
    %
    $\subrank_1(T)=n$, $\subrank_2(T)\leq c+1$ and $\subrank_3(T)\leq n/c+1$.
    In particular,
    \[
    \subrank_2(T)\subrank_3(T)\leq  \tfrac{(c+1)}{c} n+c+1.
    \]
    To define this tensor, let
    %
    \begin{align*}
    T:=\sum_{i=1}^n e_1\otimes e_i\otimes e_i+&\sum_{k=1}^{c}\sum_{j=1}^{n/c}e_{(k-1)\frac{n}{c}+j}\otimes e_j\otimes e_k\\
    =\sum_{i=1}^n e_1\otimes e_i\otimes e_i+&\sum_{i=0}^{n-1} e_{i+1}\otimes e_{(i \bmod \frac{n}{c}) + 1}\otimes e_{\lfloor i / (\frac{n}{c})\rfloor+1}.
    \end{align*}
    %
    We will prove that $\subrank_2(T)\leq c+1$ and $\subrank_3(T)\leq n/c+1$.
    %
    Looking at direction $2$ (interpreting direction $1$ as enumerating the rows and $3$ as enumerating the columns), we can see all slices are restricted to the first row (first sum) or the first $c$ columns (the other sum) 
    for non-zero entries, so each matrix in their span, is of rank $\subrank_2(T)\leq c+1$. 
    %
    Similarly, all slices in direction $3$ (interpreting direction $2$ as enumerating the rows) are restricted to the first $n/c$ rows or the first column, meaning $\subrank_3(T)\leq n/c+1$.
     
    %

    %
\end{example}

%
%
%
%
%
%
%

%

%

%
%
%

%

%
%

\begin{example}\label{ex:JopsEx}
    Let $n$ be a square.
    There is a concise tensor $T \in \FF^{n} \otimes \FF^n \otimes \FF^n$ such that for all~$\ell\in[3]$ we have
    $\sqrt{n} \leq \subrank_\ell(T)\leq 2\sqrt{n}$. In particular, for all ${\ell_1} \neq {\ell_2}$ we have $n\leq\subrank_{\ell_1}(T)\subrank_{\ell_2}(T)\leq4n$. 

    To define this tensor, let $f:[n]\to[\sqrt{n}]$ and $g:[n]\to[\sqrt{n}]$ be any two functions such that %\refnote{11}{missing "that"}\ILnote{added it} 
    the function $(f,g):[n]\to[\sqrt{n}]\times[\sqrt{n}] : i \mapsto (f(i), g(i))$ is injective and such that for every $i\in[\sqrt{n}]$ we have $f(i)=g(i)=i$. Let
    \[T:=\sum_{i=1}^{\sqrt{n}}e_i\otimes e_i\otimes e_i+\sum_{i=\sqrt{n}+1}^n e_i\otimes e_{f(i)}\otimes e_{g(i)} + e_{f(i)}\otimes e_i\otimes e_{g(i)} + e_{f(i)}\otimes e_{g(i)}\otimes e_i\]
    Then $T$ has a diagonal subtensor of size $\sqrt{n}$, so $\sqrt{n} \leq \subrank(T) \leq \subrank_\ell(T)$.
    On the other hand, $\subrank_\ell(T) \leq 2\sqrt{n}$, which can be seen using \autoref{th:Qi-SRi} %\refnote{12}{$SR_i$ is only defined below, on page 41} 
    as the support of all its slices is confined to the first $\sqrt{n}$ rows and the first $\sqrt{n}$ columns, so $\slicerank_\ell(T)\leq 2\sqrt{n}$ (see \autoref{subsec:mincov} for the definition of $\slicerank_\ell$). From injectivity of $(f,g)$ it follows that $T$ is concise. %
\end{example}

\subsection{Lower bound on the asymptotic subrank via max-ranks and matrix multiplication tensors}\label{subsec:lb-mamu}

%

We will now prove a lower bound on the asymptotic subrank of concise tensors in terms of the dimensions of the tensor. We will use the max-rank inequality (\autoref{th:uncertainty-principle-n}) proven before and basic properties of the matrix multiplication tensors that we will discuss now. The main result here is the cube-root lower bound $\asympsubrank(T) \geq (\min_\ell n_\ell)^{1/3}$ for concise tensors $T \in \FF^{n_1} \otimes \FF^{n_2} \otimes \FF^{n_3}$ (\autoref{th:cube-root-lb}). With a simpler argument we can already obtain the bound $\asympsubrank(T) \geq (\min_\ell n_\ell)^{1/4}$ (\autoref{th:fourth-root}), which we will prove first. This fourth-root lower bound is already sufficient for our use in the proof of the discreteness theorem. In \autoref{sec:lower-bound-asympsubrank} we will prove an even better square-root lower bound under a symmetry assumption.

For $a,b,c \in \NN$ we denote by $\langle a,b,c\rangle \in \FF^{ab} \otimes \FF^{bc} \otimes \FF^{ca}$ the %
matrix multiplication tensor $\langle a,b,c \rangle = \sum_{i=1}^a \sum_{j=1}^b \sum_{k=1}^c e_{i,j} \otimes e_{j,k} \otimes e_{k,i}$. 
These tensors have the multiplicative property $\langle a,b,c\rangle \otimes \langle x,y,z\rangle = \langle ax, by, cz\rangle$. They are naturally related to the parameters~$\subrank_i$ as follows:

\begin{lemma}\label{lem:Qi-mamu}
Let $T$ be a tensor and $r \in \NN$. The following holds:
\begin{itemize}
\item $\subrank_1(T) \geq r$ if and only if $T \geq \langle 1, 1,r \rangle$.
\item $\subrank_2(T) \geq r$ if and only if $T \geq \langle r, 1,1 \rangle$.
\item $\subrank_3(T) \geq r$ if and only if $T \geq \langle 1, r, 1 \rangle$.  
\end{itemize}
\end{lemma}

\begin{lemma}\label{lem:Qi-square}
    If there are distinct ${\ell_1},{\ell_2}$ such that $\subrank_{\ell_1}(T), \subrank_{\ell_2}(T) \geq r$, then $\subrank(T^{\boxtimes 2}) \geq r$.
\end{lemma}

\begin{proof}
Suppose $\subrank_1(T), \subrank_2(T) \geq r$ (the proof for the other cases goes the same). Then we have that
$
T \geq \langle1, 1, r\rangle =  \sum_{i=1}^r e_1 \otimes e_i \otimes e_i$ and 
$T \geq \langle r, 1, 1\rangle = \sum_{i=1}^r e_i \otimes e_1 \otimes e_i$.
%
We multiply these inequalities to get
$T^{\boxtimes 2} \geq \langle r,1,r\rangle= \sum_{i,j=1}^r e_i \otimes e_j \otimes (e_i \otimes e_j).
$
We apply the projection $(e_i \otimes e_j) \mapsto \delta_{i=j} e_i$ to the third tensor leg to get
$
T^{\boxtimes 2} \geq \sum_{i=1}^r e_i \otimes e_i \otimes e_i
$
which proves the claim.
\end{proof}

We are now ready to prove the fourth-root lower bound on the asymptotic subrank. After that we will prove the stronger cube-root lower bound (\autoref{th:cube-root-lb}).

\begin{theorem}\label{th:fourth-root}
    Let $T \in \FF^{n_1} \otimes \FF^{n_2} \otimes \FF^{n_3}$ be concise. Then $\asympsubrank(T) \geq (\min_\ell n_\ell)^{1/4}$.
\end{theorem}
\begin{proof}
The proof is an application of \autoref{th:uncertainty-principle-n}.
That theorem only holds over large enough fields however. 
Since the asymptotic subrank does not change under field extension \cite[Theorem~3.10]{strassen1988asymptotic}, we may assume that our base field $\FF$ is infinite (by working over the algebraic closure of the original field) so that the field size condition of \autoref{th:uncertainty-principle-n} is satisfied. %
We use that there are distinct $\ell_1, \ell_2 \in [3]$ such that 
%\refnote{13}{missing square root on $n_i$}\ILnote{maybe it was missing in an earlier version?}\Jshort{yes we fixed this earlier} 
$\subrank_{\ell_1}(T), \subrank_{\ell_2}(T) \geq \sqrt{\min_\ell n_\ell}$ (\autoref{cor:twoLargeRankDirs}).
Then $\asympsubrank(T^{\boxtimes 2}) \geq \sqrt{\min_\ell n_\ell}$ (\autoref{lem:Qi-square}) which gives the claim. %
\end{proof}




%

It turns out that the above construction is not optimal. We will now give a slightly more elaborate construction that leads to the better cube-root bound:

\begin{theorem}\label{th:cube-root-lb}
    Let $T \in \FF^{n_1} \otimes \FF^{n_2} \otimes \FF^{n_3}$ be concise. Then $\asympsubrank(T) \geq (\min_{\ell} n_\ell)^{1/3}$.
\end{theorem}

The proof of \autoref{th:cube-root-lb} makes use of basic properties of matrix multiplication tensors which we discuss in the following two lemmas. 
%\refnote{14}{The following is a repetition of matrix multiplication notation description -- it is used above in this page already} %Recall that we use the notation~$\langle a,b,c\rangle$ for~$a,b,c \in \NN$ to denote the matrix multiplication tensors.
%\Jshort{Removed it.}

\begin{lemma}\label{lem:cube-mamu}
    $T^{\boxtimes 3} \geq \langle \subrank_2(T),\subrank_3(T),\subrank_1(T) \rangle.$
\end{lemma}

\begin{proof}
    From \autoref{lem:Qi-mamu} we have %three inequalities
    %
    %
    %
    %
    %
    %
    $T \geq \langle \subrank_2(T),1,1\rangle$, 
    $T \geq \langle 1,\subrank_3(T),1\rangle$, and  
    $T \geq \langle 1,1,\subrank_1(T)\rangle$.
    %
We multiply these inequalities using the Kronecker product to get the claim.
\end{proof}


\begin{lemma}\label{lem:asympsub-min}
    $\asympsubrank(T^{\boxtimes 3}) \geq \min_{{\ell_1},{\ell_2}} \subrank_{\ell_1}(T) \subrank_{\ell_2}(T).$
\end{lemma}
\begin{proof}
    It is known that for every $a,b,c \in \NN$ the asymptotic subrank of the matrix multiplication tensor $\langle a,b,c\rangle$ is given by $\asympsubrank(\langle a,b,c\rangle) = \min \{ab, bc, ac\}$ \cite[Equation~4.6]{strassen1988asymptotic}. We combine this with \autoref{lem:cube-mamu} to get the claim.
\end{proof}

\begin{proof}[Proof of \autoref{th:cube-root-lb}]
As in the proof of \autoref{th:fourth-root} we may assume that the base field $\FF$ is infinite.
%
%
%
We use $\asympsubrank(T^{\boxtimes 3}) \geq \min_{{\ell_1},{\ell_2}} \subrank_{\ell_1}(T) \subrank_{\ell_2}(T)$ (\autoref{lem:asympsub-min}) and we use that for every distinct ${\ell_1}, {\ell_2}, {\ell_3} \in [3]$ we have $\subrank_{\ell_1}(T)\subrank_{\ell_2}(T) \geq n_{\ell_3}$ (\autoref{th:uncertainty-principle-n}) to obtain
$\asympsubrank(T^{\boxtimes 3}) \geq \min_{{\ell_1},{\ell_2}} \subrank_{\ell_1}(T) \subrank_{\ell_2}(T) \geq \min_{\ell_3} n_{\ell_3}$. %
\end{proof}






%
%
%
%
%
%
%

%
%
%

%


%

%
%
%
%
%
%
%
%
%
%
    
%
    
%

%
%
%
%
%
%
%
%
%


%
%
%

%
%
%
%
%
%
%
%

%
%
%
%
%
%

%
%
%
%
%
%
%
%
%

%
    
%
%

%
%
%
%
%
%

%
%
%
%
%
%
%
%
%



\section{Min-rank of slice spans and asymptotic subrank of narrow tensors} %
%
\label{sec:lower-bound-imbalanced}


%
%

%

In \autoref{sec:max-rank} %
we lower bounded the asymptotic subrank of a concise tensor by a growing function in the smallest local dimension (\autoref{th:cube-root-lb}). Namely, we showed that for every concise $n_1 \times n_2 \times n_3$ tensor the asymptotic subrank is at least $(\min_\ell n_\ell)^{1/3}$. 
We will use that result as an ingredient to prove the discreteness theorem in \autoref{sec:discreteness}. However, that bound will not be enough as we must also understand how asymptotic subrank behaves when the smallest dimension $\min_\ell n_\ell$ is fixed (we elaborate on this point in \autoref{sec:discreteness}). This is the topic of this section.

To briefly explain what is missing, we note that it is a priori possible that there is a sequence of concise $c \times n_1 \times n_2$ tensors with $c$ constant and at least one of $n_1, n_2$ growing\footnote{We note that in this context, since we consider only concise tensors, if $c$ is constant and at least one of $n_1, n_2$ is growing, then both $n_1$ and $n_2$ must be growing by \autoref{cl:concise-formats}.}
such that the asymptotic subrank (or asymptotic slice rank) has an accumulation point. Our result will rule this out in the appropriate sense. In particular, we prove (\autoref{theorem:concise-cn1n2-assubrank-c}) that for every $c \geq 2$ there is a number $N(c)$ such that for every $n_1, n_2$ with $\max \{n_1, n_2\} > N(c)$,
for every concise $c \times n_1 \times n_2$ tensor, the asymptotic subrank equals $c$. This will be a crucial ingredient in our proof of the discreteness theorem in \autoref{sec:discreteness}.



%

%
Whereas the proof of the asymptotic subrank lower bound in \autoref{sec:max-rank} relied for a large part on the max-rank of matrix subspaces, here we will need the related notion of min-rank of a matrix subspace.
In \autoref{subsec:minrank} we introduce the notion of min-rank, which we will use throughout. In \autoref{subsec:minrank-subrank} we prove that if the slices of a concise tensor have large min-rank, then the tensor has large subrank. In \autoref{subsec:maxrank-minrank} we see how from matrices with large max-rank we can construct matrices with large min-rank that are moreover diagonal. In \autoref{subsec:supermult} we then prove that min-rank is super-multiplicative when at least one of the matrix subspaces is diagonal, and finally we generalize the super-multiplicativity property to make it work with the output of \autoref{subsec:maxrank-minrank}.
%\ILnote{Comment 11. Removed the parentheses from the assumption of a subsoace being diagonal. I didn't find more instances of these parentheses, although the comment claims they exist...}

In \autoref{subsec:main-unb} we combine the lemmas on min-rank to prove the main result (\autoref{theorem:concise-cn1n2-assubrank-c}). Finally, as an extra, in \autoref{subsec:dimtwo}, we determine precisely what happens in the smallest case $c = 2$, and in particular determine that $N(2) = 2$.

%

To avoid potential confusion we remark that throughout we will call any matrix $A\in \FF^{n_1} \otimes \FF^{n_2}$ (that is not necessarily square) \emph{diagonal} if its support, which is defined as $\supp(A) := \{(i,j) : A_{i,j} \neq 0\}$, is a subset of $\{(i,i): 1\leq i\leq \min\{n_1, n_2\}\}$.
%

%
%


\subsection{Min-rank of matrix subspaces and its properties}\label{subsec:minrank} %
%
%

%
\newcommand{\StrExC}{T_{\textnormal{NA},c,n}}
%
%
%
%
%
%
%

%
%
%

%

%

%

%
The min-rank of a matrix subspace is simply the smallest rank of any non-zero matrix in the space. 


\begin{definition}
    %
    We define the \emph{min-rank} of a matrix subspace $\mathcal{A} \subseteq \FF^{n_1} \otimes \FF^{n_2}$ as
    %
    \[
    \minrnk(\mathcal{A}):=\min\{\rank(A):A\in\mathcal{A}\setminus\{0\}\}.
    \]
    For any collection of matrices $A_1,
    \ldots,A_c \in \FF^{n_1} \otimes \FF^{n_2}$ we define \[
    \minrnk(A_1,\dots,A_c):=\minrnk(\linspan\{A_1,\dots,A_c\}).
    \]
\end{definition}



%
%
%

%
%
%
%

We will use the following straightforward equivalence. % throughout this section.
%Given matrices $A_1, \ldots, A_c \in \FF^{n_1} \otimes \FF^{n_2}$ we call a linear combination $\alpha_1 A_1 + \cdots + \alpha_c A_c$ with coefficients $\alpha_i \in \FF$ \emph{nontrivial} if at least one of the coefficients $\alpha_i$ is nonzero.

\begin{lemma}\label{lem:minrank-nontrivial-lin-comb}
    Let $A_1,\dots,A_c \in \FF^{n_1} \otimes \FF^{n_2}$ be matrices. %
    For every $r\geq1$, the following two statements are equivalent:
    \begin{enumerate}[\upshape(a)]
        \item The matrices $A_1,\dots,A_c$ are linearly independent and $\minrnk(A_1,\dots,A_c)\geq r$.
        \item For every $u \in \FF^c \setminus \{0\}$ we have $\rank(u_1 A_1 + \cdots + u_c A_c) \geq r$.
    \end{enumerate} 
\end{lemma}
\subsubsection{Large min-rank implies large subrank}\label{subsec:minrank-subrank}

In this section we prove that if the slices of a tensor have large min-rank, then the tensor has large subrank, in the following sense:

%\Jnote{Need to look at the following}

\begin{restatable}{lemma}{lemD}
%
\label{cla-high-rank-high-subrank}
Let $T\in\FF^{n_1}\otimes\FF^{n_2}\otimes\FF^{n_3}$ be a tensor and $A_1,\ldots,A_c\in\FF^{n_1}\otimes\FF^{n_2}$ be $c$ 3-slices of it. %
If $A_1,\ldots,A_c$ are linearly independent and  $\minrnk(A_1,\dots,A_c)\geq 2c(c-1)$,
%\ILnote{the 10 should be $2c(c-1)$} 
then $\subrank(T)\geq c$.    
%
\end{restatable}
%\ILnote{Comment 5. Review A missed the difference between the lemma and the remark. I omitted the remark and changed the lemma to include it already. (the remark is a comment currently)}
%
%
%
%

%
%
%

%
%
%
%

% \begin{remark}
%     It is immediate from \autoref{cla-high-rank-high-subrank} that for any tensor $T\in\FF^{n_1}\otimes\FF^{n_2}\otimes\FF^{n_3}$, if any $c$ linearly independent 3-slices $A_1, \ldots, A_c$ of $T$ satisfy  $\minrnk(A_1, \ldots, A_c)\geq 10c^2$, then $\subrank(T)\geq c$.
% \end{remark}

%

%

%

%
%

%
%
%

\begin{proof}
%We can assume $A_1,\dots,A_c$ are the only $3$-slices (i.e. $n_3=c$) by deleting the other $3$-slices by a restriction.
%
%; some of them will be deleting a row or a column and some of them will be a row or a column operation, i.e. adding a scaled row or column to another in all $3$-slices at once (note that a row operation, for example, is actually adding a scaled $1$-slice to another, meaning performing the same row-operation in all the $3$-slices). 
%We will do this in $c$ steps. 
For $i \in [c]$ and $\alpha \in \{0,1\}$, we say that an $m_1\times m_2$ matrix $A$ has property $P(i,\alpha)$ if $A_{i,i} = \alpha$ and $A_{j, i} = 0$ for all $j \in [m_1]\setminus \{i\}$ and $A_{i,j} = 0$ for all $j \in [m_2] \setminus \{i\}$.

We may assume $n_3 = c$.
We will be changing $T$ with a series of restrictions after which the slices $A_1, \ldots, A_c$ will satisfy: for every $j \in [c]$ the matrix $A_j$ has property $P(j, 1)$ and for every $k \in [c] \setminus \{j\}$ the matrix $A_k$ has property $P(j, 0)$. Then $A_1, \ldots, A_c$ restricted to the first $c$ rows and columns form the unit tensor $\langle c\rangle$.

We will carry out the above inductively, obtaining by induction on $0 \leq i \leq c$, that 
\begin{itemize}
\item $A_1, \ldots, A_c$ are linearly independent and $\minrank(A_1, \ldots, A_c) \geq (c-i)\cdot 2(c-1)$
\item for all $1 \leq j \leq i$, 
\begin{itemize}
    \item $A_j$ has property $P(j, 1)$
    \item for all $k \in [c] \setminus \{j\}$, $A_k$ has property $P(j, 0)$.
\end{itemize}  
\end{itemize}
Throughout the induction, the restrictions we use may change the matrices $A_1, \ldots, A_c$, but we will keep denoting them by these names. In particular, the matrices $A_1, \ldots, A_c$ may become smaller. Namely, in every step of the induction we will remove at most $c-1$ rows and at most $c-1$ columns from all matrices (simultaneously). Note that the case $i=0$ is the assumption of the lemma.

Let $1 \leq i\leq c$ and assume the induction claim holds for $i-1$. 
In the following, all operations on rows and columns we apply to all matrices at the same time.
By linear independence, $A_i$ is nonzero. For $j < i$, row $j$ and column $j$ of $A_i$ are zero (by induction hypothesis), so there are $j,k \geq i$ with $(A_i)_{j,k}$ nonzero.
Swap rows $i$ and $j$ and swap columns~$i$ and $k$ and scale so that $(A_i)_{i,i} = 1$. By adding scalar multiples of row $i$ to rows $j > i$ and adding scalar multiples of column $i$ to columns $k > i$ we make row $i$ and column~$i$ in $A_i$ zero. Then $A_i$ has property $P(i,1)$. 

Now we go over $j \in [c] \setminus \{i\}$ one by one. By adding a scalar multiple of $A_i$ to $A_j$ we get $(A_j)_{i,i} = 0$. Suppose the $i$th row of $A_j$ is nonzero. Then there is a $k>i$ such that $(A_j)_{i,k}$ is nonzero (by induction hypothesis). Add scalar multiples of the $k$th column to the other columns $k'>i$ to make $(A_j)_{i,k'}$ zero. Then remove the $k$th column. Then the $i$th row of $A_j$ is zero. Repeat the same procedure for the $i$th column of $A_j$ to make it zero, possibly removing one row. As a result, $A_j$ has property $P(i,0)$.

For every $j \in [c] \setminus \{i\}$, we removed at most one row and one column, so in total we removed at most $2(c-1)$ rows or columns. Removing a row or column decreases matrix rank by at most one, so (\autoref{lem:minrank-nontrivial-lin-comb})
for any $u\in\FF^{c}\setminus\{0\}$ we have $\rank(u_1A_1+\cdots+u_cA_c)\geq
%$\minrank(A_1, \ldots, A_c) \geq 
(c-(i-1))\cdot 2(c-1) - 2(c-1) = (c-i)\cdot 2(c-1)$. 
%\ILnote{Suggested edits from the email: for the previous sentence: "...by at most one, so for any $u\in\FF^{c}\setminus\{0\}$ we have $\rank(u_1A_1+\cdots+u_cA_c)\geq...$" and for the next sentence: "...are linearly independent and their min-rank is lower bounded by it (\autoref{lem:minrank-nontrivial-lin-comb})."}.
If $i < c$, this is nonzero so $A_1, \ldots, A_c$ are linearly independent and $\minrank(A_1,\ldots,A_c) \geq (c-i)\cdot 2(c-1)$ (\autoref{lem:minrank-nontrivial-lin-comb}). (This is clear for $i=c$ since then the matrices form the tensor $\langle c\rangle$.) %(When $i = c$ linear independence follows since the matrices together from $\langle c\rangle$.) 
This proves the induction claim for $i$.
\end{proof}

% The changes will be split into $c$ steps, each step $i$ making sure the $i$'th $3$-slice $A_i$ has $(A_i)_{i,i}=1$, $0$ in the rest of its $i$'th row and column, and that any other $3$-slice has only $0$ in these row and column. During the $i$'th step we will not be ruining the achievements of previous steps as we will be operating with row and column operation and row and column deletions only for rows and columns $j\geq i$. The previous achievements include having all the at-least-$i$ rows and columns $0$ in their coordinates strictly smaller than $i$ in all $3$-slices (the achievements regarding the rows "protect" the beginning of these columns, for example), so these operations change nothing regarding the rows and columns strictly less than $i$.

% Note that row and column operations, performed for all $3$-slices together, do not change the rank of any matrix in the linear span of them, and thus do not change their min-rank. In each step we will delete at most $c-1$ rows and and at most $c-1$ columns, decreasing the min-rank \ILnote{instead: the rank of any non-trivial linear combination} by at most $2(c-1)$ (as the rank of a matrix is the dimension of its row-span and of its column-span). Hence, at the beginning of each step the min-rank of the $3$-slices is at least $10c^2-2(c-1)c>10c^2-2c^2=8c^2>0$ \ILnote{If I haven't made a mistake this time -- it looks like we can change the 10 into a 2, right?}, meaning they are still linearly independent so no $3$-slice is $0$.

% For the $i$'th step, consider the $i$'th $3$-slice $A_i$, which is nonzero by last observation. By the guarantee from the previous steps we know all the rows and columns $j<i$ are zero, so there is an $j,k\geq i$ with $(A_i)_{j,k}\neq0$. swap between rows $i$ and $j$ and between columns $i$ and $k$ so that $(A_i)_{i,i}=1$ \ILnote{scale}.
% Using row and column operations zero-out the rest of its $i$'th row and column (note again that the ones tha will need to be changed are only with index $j>i$). Use this $3$-slice to zero-out the $(i,i)$'th elements in the other $3$-slices: $(A_j)_{i,i}=0$ for any $j\neq i$.

% Next, for every $j\neq i$, we would like to make the $i$'th column of $A_j$ zero. 
% For this we might need to delete a row: we made sure $(A_j)_{i,i}=0$. If the rest of the column is zero we don't do anything. Otherwise, there is some $k>i$ with $(A_j)_{k,i}\neq 0$. We subtract a scaled version of this row from any other row $k'$ that has $(A_j)_{k',i}\neq0$. Note that these can only be $k'>i$ by the guarantee from previous steps. Finally, we delete row $k$ from our tensor. 

% We follow a similar procedure to zero-out the $i$'th row of $A_j$, with possibly deleting a column from our tensor at the end.

% Note that after these we indeed achieve the required guarantee: we haven't ruined the guarantee from previous steps and we zeroed-out the $i$'th row and column from all $3$-slices apart from $(A_i)_{i,i}=1$. For each $3$-slice except $i$ we deleted at most one row and one column to achieve this, so during the step we indeed deleted at most $c-1$ rows and at most $c-1$ columns.

% When finishing the $c$'th step, we have the tensor $\langle c\rangle$ as the corner subtensor.
% \end{proof}

%
%

%
%
%

%
%

%
%

%
%
%

%
%
%

%
%
%


%
%
%

%
%

%


%

%

%

%

%
%
%
%
%

%
%

%
%



%


%
%
%

%
%
%


%
%
%
%

%
%
%
%

%
%
%

%
%
%

%

%

%

%
%

%
%
%

%
%
%
%

%
%
%
%
%

%

%
%
%

%
%
%
%
%

%
%
%
%
%
%
%

%

\subsubsection{From large max-rank to large min-rank and diagonal}\label{subsec:maxrank-minrank}

In this section we discuss how from any collection of matrices that contains \emph{at least one} matrix of large rank (large max-rank) we can produce a collection of slightly smaller matrices such that \emph{all} nonzero matrices in their span have large rank (large min-rank). Moreover, the resulting matrices are all diagonal, which will be crucial in our application.

For any matrix $A\in \FF^{n_1} \otimes \FF^{n_2}$ %
and subsets $I\subseteq [n_1]$, $J\subseteq [n_2]$ we denote by $A_{I\times J}$ the submatrix of $A$ with rows indexed by $I$ and columns indexed by~$J$.  For any subset $X \subseteq \FF^{n_1} \otimes \FF^{n_2}$ we denote by $X_{I \times J}$ the set $\{A_{I\times J} : A \in X\}$.


\begin{restatable}{lemma}{lemBx}
%\begin{lemma}
\label{lem:minrk_diag}
\label{claim:matrices-into-diagonal-with-large-rank}
For every $c\in\NN$, there is an $\eps(c)>0$ such that the following holds.
%Let $n_1\leq n_2$ be positive integers and 
Let $\mathcal A\subseteq \FF^{n_1\times n_2}$ be a matrix subspace of dimension~$c$.
Then, there exist $U\in \GL(\FF^{n_1})$, $V\in \GL(\FF^{n_2})$, a subset $J\subseteq [\min_\ell n_\ell]$, a basis $B_1,\dots,B_c$ for the space $U\mathcal AV$ and $b\in[c]$ such that:
\begin{itemize}
    \item $(B_1)_{J\times J},\dots,(B_b)_{J\times J}$  are diagonal and linearly independent;
    \item $\minrank\big((B_1)_{J\times J}, \ldots, (B_b)_{J\times J}\big) \geq \eps(c) \maxrank(\mathcal{A})$; % n_2$;
    \item $(B_{b+1})_{J\times J} = \cdots = (B_c)_{J\times J} = 0$.
\end{itemize}
%\end{lemma}
\end{restatable}


Before discussing the proof in detail, as a motivation and ingredient for later, we observe using the pigeonhole principle that any concise tensor has at least one slice of large rank. 

\begin{restatable}{lemma}{lemA}
%
\label{cla:high-rk}
Let $T\in\FF^{n_1}\otimes\FF^{n_2}\otimes\FF^{c}$ be a concise tensor. Then $T$ has at least one 3-slice 
of rank at least $\max\{n_1,n_2\}/c$.    
%
\end{restatable}

\begin{proof}
Suppose $n_1\geq n_2$. %
Concatenate the 3-slices $A_1, \ldots, A_c$ of $T$ to an $n_1\times cn_2$ matrix. Conciseness of $T$ implies that the rank of this matrix is $n_1$ so there are $n_1$ linearly independent columns. By the pigeonhole principle, there is an $i \in [c]$ such that $A_i$ contains at least $n_1/c$ linearly independent columns.   
\end{proof}


% %
% %

% \begin{restatable}{lemma}{lemB}
% %
% \label{claim:matrices-into-diagonal-with-large-rank}
%     Let $0<\eps<1$ and $0<\alpha\leq 1$. For every $c\geq 2$ there is an $N_0(\eps,\alpha,c)\in\NN$ such that for any $n_1,n_2>N_{0}(\eps,\alpha,c)$, and any matrices $A_1,\dots,A_c\in \FF^{n_1}\otimes \FF^{n_2}$ that contain a matrix of rank at least $\alpha \min\{n_1,n_2\}$ in their span, there are invertible matrices 
%     \[
%     R\in \GL_{n_1}(\mathbb{F}),\, C\in \GL_{n_2}(\mathbb{F}),\, S\in \GL_{c}(\mathbb{F})
%     \]
%     and subsets $I\subseteq[n_1],J\subseteq[n_2]$ such that  the matrices in
%     \[
%     X = \biggl\{\Bigl(R\Bigl(\sum_{i=1}^cS_{ij}A_i\Bigr)C\Bigr)\Big|_{I\times J}\biggr\}_{j\in [c]}
%     \]
%     are diagonal, not all zero, have min-rank $\minrnk(X) \geq\min\{n_1,n_2\}^\eps$ and the non-zero elements of $X$ are linearly independent.
% %
% \end{restatable}

% \begin{remark}
%     We stress that in \autoref{claim:matrices-into-diagonal-with-large-rank} it is crucial that we are taking $S$ to be an invertible operation, because we want to maintain the property that the matrices $A_1, \ldots, A_c$ are linearly independent. Even though some matrices will become zero when restricted to $I\times J$ we will use the linear independence of those matrices at a later stage.
% \end{remark}

% %
% %
% The proof of \autoref{claim:matrices-into-diagonal-with-large-rank} has two main parts, which we present separately as the following two lemmas, and which imply it immediately.
% \begin{lemma}\label{lem:matrices-into-diagonal-with-one-large-rank-mat}
%     Let $0<\alpha\leq 1$. For every $c\geq 2$, $n_1,n_2\in \NN$ and any matrices $A_1,\dots,A_c\in \FF^{n_1}\otimes \FF^{n_2}$ that contain a matrix of rank at least $\alpha \min\{n_1,n_2\}$ in their span, there are invertible matrices 
%     \[
%     R\in \GL_{n_1}(\mathbb{F}),\, C\in \GL_{n_2}(\mathbb{F}),\, S\in \GL_{c}(\mathbb{F})
%     \]
%     and subsets $I\subseteq[n_1],J\subseteq[n_2]$ such that  the matrices in
%     \[
%     X = \biggl\{\Bigl(R\Bigl(\sum_{i=1}^cS_{ij}A_i\Bigr)C\Bigr)\Big|_{I\times J}\biggr\}_{j\in [c]}
%     \]
%     are diagonal and one of them is of rank at least $\tfrac{1}{3^{c-1}}\alpha\min\{n_1,n_2\}$.
% \end{lemma}

% \JBnote{I wonder if Lemma~\ref{lem:matrices-into-diagonal-with-one-large-rank-mat} can be proved by induction on~$c$ from the following statement, which I believe is what underlies Algorithm~1: For any $A\in \FF^{n\times n}$ there are invertible matrices $U,V\in \FF^{n\times n}$ and a set $I\subseteq [n]$ of size at least $n/3$ such that the following holds. The principal submatrix $(UAV)_{I\times I}$ is diagonal and for every $x\in \FF^n$, we have that $(U\Diag(x)V)_{I\times I} = \Diag(x)_{I\times I}$.}

% \ILnote{Yes, that's exactly what's going on there. It's a different language for the same thing -- induction instead of iterations (with iteration invariant). If it's more accessible as an induction I have no objection for changing it.}

% %
% %
% %
% \begin{lemma}\label{lem:diag-matrices-with-one-large-rank-mat-into-large-minrank}
%     Let $0<\eps<1$ and $0<\alpha\leq 1$. For any $c\geq2$ there is an $N_0(\eps,\alpha,c)\in\NN$ such that for any $n_1,n_2>N_0(\eps,\alpha,c)$, and any matrices $A_1,\dots,A_c\in\FF^{n_1}\otimes\FF^{n_2}$ and sets of indices $I\subseteq[n_1],J\subseteq[n_2]$ such that $A_1|_{I\times J},\dots,A_c|_{I\times J}$ are diagonal and one of them is of rank at least $\tfrac{1}{3^{c-1}}\alpha\min\{n_1,n_2\}$, there is an invertible matrix $S\in \GL_c(\FF)$ and index sets $\tilde I\subseteq[n_1],\tilde J\subseteq[n_2]$ such that 
%     \[X=\left\{\left(\sum_{i=1}^cS_{ij}A_i\right)\Big|_{\tilde I\times \tilde J}\right\}_{j\in[c]}\]
%     are diagonal, not all zero, have min-rank $\minrnk(X)\geq \min\{n_1,n_2\}^\eps$ and the non-zero elements of $X$ are linearly independent.
% \end{lemma}
% %
% %
% %
% \begin{proof}[Proof of \autoref{lem:matrices-into-diagonal-with-one-large-rank-mat}]
%     Assume $n_1\leq n_2$ (otherwise swap their names).
    
%     Note that the invertible matrix $S$ in the statement of the lemma
%     is just replacing $A_i$'s with linear combinations of them, in a way that keeps their span unchanged, and this is how we will treat it. Similarly, $R$ is invertible row operations done for all the matrices together and $C$ is invertible column operations done for all the matrices together. 
%  %

%  %
%  We start by describing the way $R,C,S$ and $I,J$ are constructed and showing the diagonal property this construction guarantees, together with having a high-rank matrix. 
 
%  %
%  %

%  %
 
%  Without loss of generality (using the $S$'s) we may assume that $\rank(A_1) \geq \alpha n_1$ and that $(A_1)|_{[\alpha n_1]\times [\alpha n_1]} = \id_{\alpha n_1}$ (assume $\alpha n_1\in\NN$. We avoid writing $\lceil\alpha n_1\rceil$ to simplify notation).

% %
 
%  At the beginning, set $I=[\alpha n_1]$. We will iteratively remove indices from it as we go along. $J$ will just be the same set $J=I$. We work in two nested loops: the outer loop (the for-loop in \autoref{alg:diagonalise-large-minrank} below)  goes through the matrices $A_k$ for $k>1$ and makes each diagonal in $I\times I$. It does so by the inner loop (the While loop in \autoref{alg:diagonalise-large-minrank})  which goes over $I$ and changes the matrices by row and column operations and removes some of the indices from $I$. For example, to zero-out a row (except its diagonal element) we use a column which is non-zero in the desired row and subtract scaled versions of that column from all other columns. After that we take the subtracted column's index out of $I$. This process is detailed in the form of an algorithm in \autoref{alg:diagonalise-large-minrank} below.

% \begin{algorithm}[H]
% \caption{Diagonalise with a high-rank element: Construction of $R,C$ and the beginning construction of $I,J$}\label{alg:diagonalise-large-minrank}
%     \begin{algorithmic}
%     \Require $A_1|_{[\alpha n_1]\times[\alpha n_1]}=\id_{\alpha n_1}$
%     \State $R,C\gets\emptyset$ \Comment{$R,C$ will be sets of row and column operations, resp.}
%     \State $I_1\gets[\alpha n_1]$ 
%     \For{$k=2,\dots,c$}
%         \State $I_k\gets I_{k-1}$ \Comment{$I_{k}$ is the current set of indices, some of which we will remove}
%         \State $X_k\gets \emptyset$ \Comment{$X_k$ is the set of indices that we have dealt with}
%         \While{$I_k\setminus X_k\neq\emptyset$}
%             \State $i_k\gets \min I_k\setminus X_k$
%             \If{$\exists j_1\in I_k\setminus\{i_k\}:(A_k)_{i_k,j_1}\neq0$}
%                 \State Subtract a multiple of column $j_1$ from every column $j\in I_k\setminus\{i_k,j_1\}$,
%                 \State making $(A_k)_{i_k,j}$ zero
%                 \State Add these column operations to $C$
%                 \State Remove $j_1$ from $I_k$
%             \EndIf \Comment{Now row $i_k$ in ${A_k}|_{I_k\times I_k}$ is zero outside the diagonal}
%             \If{$\exists j_2\in I_k\setminus\{i_k\}:(A_k)_{j_2,i_k}\neq0$}
%                 \State Subtract a multiple of row $j_2$ from every row $j\in I_k\setminus\{i_k,j_2\}$,
%                 \State making $(A_k)_{j,i_k}$ zero
%                 \State Add these row operations to $R$
%                 \State Remove $j_2$ from $I_k$
%             \EndIf \Comment{Now column $i_k$ in ${A_k}|_{I_k\times I_k}$ is zero outside the diagonal}
%             \State Add $i_k$ to $X_k$
%         \EndWhile
%     \EndFor
%     \State $I\gets I_c$
%     \State $J\gets I$
%     \end{algorithmic}
% \end{algorithm}

% %
 
%  %
% We claim the following iteration invariant: At the end of each iteration of the inner loop (the While loop)
% %
% \begin{enumerate}[(1)]
%     \item  $A_k|_{I_k\times I_k}$'s rows and columns with indices at most $i_k$ are all zero except in their diagonal elements,
%     \item $A_{k'}|_{I_k\times I_k}$ for $k'<k$ are diagonal, 
%     \item and $A_1|_{I_k\times I_k}=\id_{|I_k|}$.
% \end{enumerate}
% %
% %
% We now go over the items of the invariant carefully to make sure they hold.
% The first item ($1$) holds by the conditionals and the operations inside of them -- which explicitly guarantee it, first for the row and then for the column. If we go into the first conditional, say, and subtract a scaled version of column $j_1$ from the other columns, then $(A_k)_{i_k,j_1}\neq0$ but by ($1$) for the previous iteration we know that $j_1>i_k$ (and similarly $j_2>i_k$ in the second conditional), and after this operation the $i_k$'th row of $A_k|_{I_k\times I_k}$ is zero except in its $i_k,j_1$ coordinates. Additionally, by item ($2$) of the invariant in $A_{k'}$ for $k'<k$ the $j_1$ column is nonzero only on the diagonal, so this operation can blemish only the $j_1$'th row of $A_{k'}$ and as we remove $j_1$ from $I_k$ we keep this item as well. Lastly, note that every row or column that may change in $A_1$ have their indices removed from $I_k$ when that happens so indeed we keep $A_1|_{I_k\times I_k}=\id_{|I_k|}$.

% By looking at the end of the final iteration, this invariant shows that when the algorithm terminates we have $A_1|_{I\times J}=\id_{|I|}$ and all of $A_k|_{I\times J}$ are diagonal. 
% We showed above that whenever we go into one of the conditionals the picked $j_1$ or $j_2$ must be $j_1,j_2>i_k$. Thus, every index put into $X_k$ cannot be removed from $I_k$ in a later iteration as a $j_1$ or $j_2$ (as long as we're in the same iteration $k$ of the outer loop). At the end of the iteration for $k$ of the outer loop we have $X_k=I_k$ and each index in that set was responsible to at most $2$ indices removed from $I_k$, implying $|I_k|\geq \frac{1}{3}|I_{k-1}|$. So $|I|=|I_c|\geq\frac{1}{3^{c-1}}|I_1|=\frac{1}{3^{c-1}}\alpha n_1$. Specifically, $\rank(A_1|_{I\times J})\geq\frac{1}{3^{c-1}}\alpha n_1$ by item ($3$) of the invariant.
% \end{proof}
 
%  %
%  %
 
%  %

%  %
%  %
% %

%  \begin{proof}[Proof of \autoref{lem:diag-matrices-with-one-large-rank-mat-into-large-minrank}]
%  We need to guarantee that every linear combination of the nonzero diagonal matrices will be of large rank (which is equivalent to being linearly independent with large min-rank by \autoref{lem:minrank-nontrivial-lin-comb}). 
%  %
%  To achieve this we will replace matrices that appear in a low rank combination with that low rank matrix. We shall then restrict our attention to the complement of the nonzero entries of each such matrix. 

%  Assume w.l.o.g. $n_1\leq n_2$.
%  We know one of the diagonal matrices $A_i$ is of rank at least $\tfrac{1}{3^{c-1}}\alpha n_1$, so we may remove from $I$ and $J$ any index $j$ for which $(A_i)_{j,j}=0$ and we still have all the hypothesis of the lemma met. We may also change their order so that $A_i$ is the first of them (by recording the change of order into the matrix $S$, in which case at the beginning of the algorithm below $S$ is not empty).

%  The process described above is detailed precisely in \autoref{alg:zero-or-highrank} below.

% %
% %

% %
% %
 
% \begin{algorithm}[H]
%     \caption{Make the elements zero or high rank: Construction of $S$ and finishing the construction of $I,J$}\label{alg:zero-or-highrank}
%     \begin{algorithmic}
%         \Require $(A_1)|_{I\times J}=\id$
%         \State $S\gets\emptyset$ \Comment{$S$ will hold instructions to replace a matrix by a} 
%         \State {\hspace{21em}linear combination that includes it}
%         \State $Z\gets\emptyset$
%         \For{$k=1,\dots,c$}
%             \If{$(A_k)|_{I\times J}\neq0$}
%                 \State Add $k$ to $Z$
%             \EndIf
%         \EndFor
%         \While{$\exists A\in\linspan\{A_k\}_{k\in Z}:\rank(A|_{I\times J})<n_1^\eps$}
%             \State Let $A=\sum_{k\in Z}a_k A_k$
%             \State $k_0\gets \max\{k\in Z:a_k\neq0\}$
%             \State $A_{k_0}\gets A$
%             \State Add this replacement to $S$
%             \State Remove the coordinates $i\in I$ for which $A_{i,i}\neq0$ from $I$ and $J$
%             \State Remove $k_0$ from $Z$
%         \EndWhile
%     \end{algorithmic}
% \end{algorithm}

% Note that $Z$ becomes smaller by one in each iteration of the While loop, so the loop has to end in at most $c$ iterations. Additionally, we will show that under some choice for $N_0(\eps,\alpha,c)$, $k_0>1$ in all the iterations, meaning $1\in Z$ at the end of the While loop and thus $Z\neq\emptyset$ at the end of the algorithm, and all matrices in the span of $\{A_k\}_{k\in Z}$ have rank at least $n_1^\eps$ when restricted to rows in $I$ and columns in $J$. The $A_k$'s are still all diagonal in $I\times J$ as taking linear combinations of diagonal matrices keep them diagonal, so all we have left is to explain the choice of $N_0(\eps,\alpha,c)$.

% To guarantee the index $k_0$ in the While loop is always $k_0>1$, we would like to have $|I|\geq n_1^\eps$ at all iterations. That way, $\rank((A_1)|_{I\times J})=\rank(\id_{|I|})\geq n_1^\eps$ and so $A$ with a smaller rank must have at least one matrix $A_k$ with $k>1$ and $a_k\neq0$.

% Recall $|I|\geq\frac{1}{3^{c-1}}\alpha n_1$ at the beginning of this algorithm, and note that in each iteration we remove at most $n_1^\eps$ coordinates from it (the rank of $A$). There are at most $c-1$ iterations before considering $A_1$, so $I$ is of size $|I|\geq\alpha n_1/3^{c-1}-(c-1)n_1^\eps$. As $\alpha n_1/3^{c-1}=\omega(cn_1^{\eps})$ (by the assumptions $\eps<1$ and $\alpha>0$), there is $N_0(\eps,\alpha,c)\in\NN$ such that for any $n_1\geq N_0(\eps,\alpha,c)$ we have $\frac{\alpha}{3^{c-1}}>c n_1^\eps$ and thus $|I|>n_1^\eps$ as required.
%  %
%  %
%  %
% \end{proof}

% %
% %

% \subsubsection{Induction proofs of the above}

\paragraph{Diagonalization.}
The proof of \autoref{lem:minrk_diag} has three parts.
In this first part, we show in \autoref{lem:diagonalize} that any collection of matrices can be made simultaneously diagonal on a large principal submatrix in such a way that the property of having large max-rank is preserved. 
%\Jnote{Added:}
For $a,b \in [n]$ let $E_{a,b} \in \FF^n \otimes \FF^n$ be the matrix with coefficient one at index $(a,b)$ and coefficient zero elsewhere.

\begin{lemma}\label{lem:diagonalize-ind2}
    Let $A \in \FF^{n \times n}$. Then there are $U,V \in \GL(\FF^n)$ and $i,j \in \{2,\ldots, n\}$ such that for $I = [n] \setminus\{i,j\}$,
    %\refnote{16}{the notation for $E_{a,b}$ is only defined much later, on page 39}\Jshort{Fixed above}
    \begin{enumerate}[\upshape (i)]
    \item $(UAV)_{I\times I}  = \alpha_{1,1} E_{1,1} + \sum_{a,b \geq 2} \alpha_{a,b} E_{a,b}$ for some $\alpha_{a,b} \in \FF$;
    \item for every diagonal matrix $B\in \FF^{n\times n}$, we have that $(U B V)_{I\times I} = B_{I\times I}$.
    \end{enumerate}
\end{lemma}
\begin{proof}
    Let $X = \{\ell \in \{2, \ldots, n\} : A_{1, \ell} \neq 0\}$. If $X = \emptyset$, then we let $V = \id$ and let~$j$ be any element of $[n]$. Otherwise, we let $j$ be any element of $X$, and we let $V$ be the matrix that, when multiplied with $A$ from the right, adds multiples of the $j$th column of $A$ to the $\ell$th column of $A$ such that the $(1, \ell)$ entry becomes zero, for every $\ell \in X \setminus \{j\}$%
    %\refnote{17}{$i$ should be $j$}\ILnote{fixed it}
    . In other words, we let
    \[
    V = \id - \sum_{\mathclap{\ell \in X \setminus \{j\}}} A_{1,j}^{-1} A_{1,\ell} E_{j, \ell}.
    \]
    Let $Y = \{\ell \in \{2, \ldots, n\} : A_{\ell, 1} \neq 0\}$. If $Y= \emptyset$ or $Y = \{j\}$, then we let $U = \id$ and we let $i$ be any element of $[n] \setminus \{j\}$. Otherwise, let $i$ be any element of $Y \setminus \{j\}$, and we let $U$ be the matrix that, when multiplied with $A$ from the left, add multiples of the $i$th row of $A$ to the $\ell$th row of $A$ such that the $(\ell, 1)$ entry becomes zero, for every $\ell \in Y \setminus \{i,j\}$.
    In other words,
    \[
    U = \id - \sum_{\mathclap{\ell \in Y \setminus \{i,j\}}} A_{i,1}^{-1} A_{\ell,1} E_{\ell, i}.
    \]
    Then claim (i) is satisfied. 
    To see that (ii) is satisfied we note that $U = \id + U_0$ and $V = \id + V_0$, where~$U_0$ supported by the $i$th column and of~$V_0$ is supported by the $j$th row.
    %Suppose~$B\in \FF^{n\times n}$ is a diagonal matrix.
    Then, for any diagonal matrix $B \in \FF^{n\times n}$,
    \begin{align*}
        UBV &= B + U_0B + BV_0 + U_0BV_0.
    \end{align*}
    Since $i\ne j$, it follows that $U_0BV_0 = 0$.
    Moreover, $U_0B$ is supported by the $i$th column and $BV_0$ is supported by the $j$th row, so $(U_0B)_{I \times I}=0$ and $(BV_0)_{I \times I}=0$, and so $(UBV)_{I \times I} = B_{I\times I}$.
    %the support of $U \Diag(x) V$ is contained in the support of $UV$, which is contained in $\{(i,i) : i \in [n]\} \cup \{(i,\ell): \ell \in [n]\} \cup \{(\ell,j): \ell \in [n]\}$. \Jshort{This is not enough to get (ii) because we wanted to really keep $\Diag(x)$.}
\end{proof}

\begin{lemma}\label{lem:diagonalize-ind}
    Let $A \in \FF^{n \times n}$. Then there are $U,V \in \GL(\FF^n)$ and $I \subseteq [n]$ such that $|I| \geq n/3$ and 
    \begin{enumerate}[\upshape (i)]
    \item $(UA V)_{I \times I}$ is diagonal;
    \item for every diagonal matrix $B \in \FF^{n \times n}$,  we have that $(U B V)_{I \times I} = B_{I \times I}$.
\end{enumerate}
\end{lemma}
\begin{proof}
    The proof is by repeatedly applying \autoref{lem:diagonalize-ind2}.
    Applying \autoref{lem:diagonalize-ind2} to $A$ gives $U,V,I$ such that $(UAV)_{I \times I}$ is of the form 
    \[
    \alpha_{1,1} E_{1,1} + \sum_{a,b\geq 2} \alpha_{a,b} E_{a,b}
    \]
    and such that for any diagonal $B$ we have $(UBV)_{I \times I} = B_{I \times I}$.
    Suppose we are given $A' \in \FF^{m \times m}$ of the form 
    \[
    \sum_{i=1}^k \alpha_{i,i} E_{i,i} + \sum_{a,b\geq k+1} \alpha_{a,b} E_{a,b}.
    \]
    Applying \autoref{lem:diagonalize-ind2} to $A'' = (A')_{\{k+1, \ldots m\} \times \{k+1, \ldots m\}}$ gives $U', V', I'$ such that $(U' A'' V')_{I' \times I'}$ is of the form $\alpha'_{1,1} E_{1,1} + \sum_{a,b\geq 2} \alpha'_{a,b} E_{a,b}$. Then 
    %\refnote{18}{Expecting $\id\oplus U'$ rather than $U'\oplus\id$, based on $A'$'s structure} 
    %\Jshort{I think he is right.}
    \[
    ((\id_k\oplus U') A' (\id_k\oplus V'))_{I' \times I'}
    \]
    is of the form 
    \[
    \sum_{i=1}^{k+1} \alpha_{i,i} E_{i,i} + \sum_{a,b\geq k+2} \alpha_{a,b} E_{a,b}
    \]
    and for any diagonal $B$ we have $((\id_k\oplus U')UBV(\id_k\oplus V'))_{I' \times I'} = B_{I' \times I'}$. In every step of the process $k$ goes up by one, and $|I|$ goes down by two, so that in the end we obtain the claim with $I \geq n/3$.
\end{proof}

\begin{lemma}\label{lem:diagonalize}
Let $A_1 = \id_n, A_2, \ldots, A_c \in \FF^{n \times n}$. Then there are $U,V \in \GL(\FF^n)$ and $I \subseteq [n]$ such that $|I| \geq 3^{-(c-1)} n$ and
\begin{enumerate}[\upshape (i)]
    \item $(U \id_n V)_{I\times I} = (\id_n)_{I\times I}$
    \item $(UA_i V)_{I \times I}$ is diagonal for all $i \in [c]$.
\end{enumerate}
\end{lemma}
\begin{proof}
    The proof is by induction over $c$ using \autoref{lem:diagonalize-ind}. 
    The base case $c = 1$ is clear.
    For the induction step, we are given $A_1, \ldots, A_m$, $U,V$ and $I \subseteq[n]$ with $|I| \geq 3^{-(m-1)}n$ such that for $B_i = (UA_i V)_{I \times I}$ we have $B_1 = \id_I$ and every $B_i$ is diagonal.
    Let $B_{m+1} = (U A_{m+1} V)_{I \times I}$. We apply \autoref{lem:diagonalize-ind} to $A = B_{m+1}$. This gives $U', V' \in \GL(\FF^I)$ and $J \subseteq I$ such that $|J| \geq \tfrac13 |I|$, $(U' B_{m+1} V')_{J \times J}$ is diagonal, and $(U' B_i V')_{J \times J} = (B_i)_{J \times J}$ for every $i \in [m]$. This proves the claim.
\end{proof}

\paragraph{From large max-rank to large min-rank.}
In this second part of the proof of \autoref{lem:minrk_diag}, we show in \autoref{lem:large-minrank} that any low-dimensional matrix subspace~$\mathcal A$ of diagonal matrices with large max-rank contains a principal-matrix subspace with large min-rank.
Moreover, in \autoref{lem:basis-extension-matrices}, we show that any matrix subspace has a basis that, when restricted to a particular principal submatrix, consists of a set of linearly independent matrices and all-zero matrices.

For any vector $w \in \FF^n$ let $\supp(w) = \{i \in [n]: w_i \neq 0\}$ and let $|w| = |\supp(w)|$. For any subspace $W \subseteq \FF^n$ let
\begin{align*}
\maxsupp(W) &= \max \{ |w| : w \in W\},\\
\minsupp(W) &= \min \{ |w| : w \in W \setminus \{0\}\}.
\end{align*}

\begin{lemma}\label{lem:minsupp}
Let $V \subseteq \FF^n$ be a subspace such that $\dim(V) \leq c$.  
%and suppose there is an element in $V$ with support size $s$. 
There is an $I \subseteq [n]$ such that $V_I$ is not zero and
\[
\minsupp(V_I) \geq \tfrac1c \maxsupp(V).
\]
\end{lemma}
\begin{proof}
    We will use the general fact that for arbitrary vectors $w_1, \ldots, w_\ell \in W$, if for every $j\in \{0, \ldots, \ell-1\}$ we have $\cup_{i \leq j} \supp(w_i) \neq \cup_{i \leq j + 1} \supp(w_i)$, then $w_1, \ldots, w_\ell$ are linearly independent.

    Let $k = \maxsupp(V)$ and $d = \dim(V)$.
    Before constructing $I$, we first construct a set $X \subseteq V$ by the following process, starting with $X$ empty.
    If there is a nonzero element $v_1 \in V$ such that $|v_1| < \tfrac1c k$, then add $v_1$ to $X$.
    If $X = \{v_1, \ldots, v_\ell\}$ and there is an element $v_{\ell + 1} \in V$ such that 
    $0 < |\cup_{i \leq \ell+1} \supp(v_i) \setminus \cup_{i \leq \ell} \supp(v_i)| < \tfrac1c k$, then add $v_{\ell+1}$ to $X$.
    By the aforementioned general fact, $X$ is linearly independent and thus $X$ is finite, namely $|X| \leq d$.

    % We claim that in fact $|X| \leq d-1$. Suppose $|X| = d$, then $V$ has a basis $v_1,\ldots, v_d$ satisfying
    % $|\cup_{i \leq \ell+1} \supp(v_i) \setminus \cup_{i \leq \ell} \supp(v_i)| < \tfrac1c k$ for every $\ell$
    % and so we have the inequality $| \cup_{i \in [d]} \supp(v_i)| < d \tfrac1c k \leq k$. However, there is an element $v \in V$ with $|v|  = k$ which cannot be in the span of these $v_1, \ldots, v_d$, which is a contradiction.
    
    We let $J = \cup_i \supp(v_i)$ and $I = [n] \setminus J$.
    
    First we claim that $V_I$ is not zero. We have $|J| < d\tfrac1c k \leq k$.
    Then for any element $v \in V$ with $|v| = k$ we have $|v_I| \geq k - |J| > 1$ and thus $v_I \neq 0$.
    
    Second, we claim $\minsupp(V_I) \geq \tfrac1c k$.  Suppose there is a nonzero element $w \in V_I$ such that $|w| < \tfrac1c k$, then this contradicts the way we constructed $X$.
\end{proof}


\begin{lemma}\label{lem:basis-extension}
    Let $V \subseteq \FF^n$ be a subspace of dimension $c$, and let $I \subseteq [n]$. Then there is a basis $v_1, \ldots, v_c$ for $V$ such that 
    \begin{itemize}
        \item $(v_1)_I, \ldots, (v_b)_I$ are linearly independent
        \item $(v_{b+1})_I = \cdots = (v_c)_I = 0$
    \end{itemize}
    for $b = \dim(V_I)$.
\end{lemma}
\begin{proof}
    Choose any basis $w_1, \ldots, w_c$ for $V$. Let $M$ be the matrix with these basis vectors as columns. The submatrix $M_{I \times[n]}$ has rank $b$ so there is a subset $J \subseteq [n]$ of size $b$ such that $M_{I \times J}$ has rank $b$.
    Let $\{v_1, \ldots, v_b\} = \{w_j : j \in J\}$ and $\{v_{b+1}, \ldots, v_c\} = \{w_j : j \not\in J\}$.
    Then $(v_1)_I, \ldots, (v_b)_I$ form a basis of $V_I$. By subtracting multiples of $v_1, \ldots, v_b$ from the vectors $v_{b+1}, \ldots, v_c$ we can ensure that $(v_{b+1})_I = \cdots = (v_c)_I = 0$.
\end{proof}

\begin{lemma}\label{lem:large-minrank}
    Let $\mathcal{A} \subseteq \FF^{n \times n}$ be a subspace of diagonal matrices such that $\dim(\mathcal{A}) \leq c$. There is an $I \subseteq [n]$ such that 
    \[
    \minrank(\mathcal{A}_{I\times I}) \geq \tfrac1c \maxrank(\mathcal{A}).
    \]
\end{lemma}
\begin{proof}
    This follows from \autoref{lem:minsupp} by identifying every diagonal matrix $A \in \mathcal{A}$ with the vector of its diagonal coefficients $v = (A_{1,1}, \ldots, A_{n,n})$ and $\rank(A) = |\supp(v)|$.
\end{proof}

\begin{lemma}\label{lem:basis-extension-matrices}
    Let $\mathcal{A} \subseteq \FF^{n_1 \times n_2}$ be matrix subspace of dimension $c$, and $J \subseteq [\min_\ell n_\ell]$. Then there is a basis $B_1, \ldots, B_c$ for $\mathcal{A}$ such that
    \begin{itemize}
        \item $(B_1)_{J \times J}, \ldots, (B_b)_{J \times J}$ are linearly independent
        \item $(B_{b+1})_{J \times J}= \cdots = (B_{c})_{J \times J} = 0$
    \end{itemize}
    for $b = \dim(\mathcal{A}_{J\times J})$.
\end{lemma}
\begin{proof}
This follows from applying \autoref{lem:basis-extension} to $V = \mathcal{A}$ and $I = J \times J$.
\end{proof}


\paragraph{Combining the above.} \!We now combine \autoref{lem:diagonalize}, \autoref{lem:large-minrank} and \autoref{lem:basis-extension-matrices} to prove \autoref{lem:minrk_diag}.

\begin{proof}[Proof of \autoref{lem:minrk_diag}]
Let $k = \maxrank(\mathcal A)$.
    By Lemma~\ref{cla:high-rk}, we have that $k \geq n_2/c$.
    After a basis transformation, we may assume that $\mathcal A$ has a basis $A_1 = \id_k,A_2,\dots,A_c$.
    Applying \autoref{lem:diagonalize} to the submatrices $(A_i)_{[k]\times [k]}$, we obtain matrices $U,V\in \GL(\FF^k)$ and a set $I\subseteq[k]$ of size at least $3^{-(c-1)}k$ such that
    \begin{enumerate}[(i)]
        \item $A_1' = \big((U\oplus \id_{n_1 - k})A_1(V\oplus \id_{n_2 - k})\big)_{I\times I} = (\id_k)_{I\times I}$;
    \item $A_i' = \big((U\oplus \id_{n_1 - k})A_i(V\oplus \id_{n_2 - k})\big)_{I\times I}$ is diagonal for all $i\in [c]$.
    \end{enumerate}
    %Let~$V\subseteq \FF^I$ be the vector space spanned by the diagonals of the submatrices~$A_i'$.
    %It follows from~(i) that~$\maxsupp(V) = |I|$.
    %Hence, by Lemma~\ref{lem:minsupp} and the fact that the matrices~$A_i'$ are diagonal that, 
    By \autoref{lem:large-minrank} there is a set~$J\subseteq I$ such that
    \beqn
    \minrnk\big((A_1')_{J\times J},\dots,(A_c')_{J\times J}\big) \geq \tfrac{1}{c} 3^{-(c-1)}k. % \geq \frac{1}{c^2} 3^{-(c-1)} n^2.
    \eeqn
    The claim follows from applying \autoref{lem:basis-extension-matrices} with $B_i = A'_i$.
\end{proof}

\subsubsection{Super-multiplicativity of min-rank}\label{subsec:supermult}

For subspaces of matrices $\mathcal{A}$ and $\mathcal{B}$ their tensor product $\mathcal{A} \otimes \mathcal{B}$ is the subspace of matrices spanned by the products $A \otimes B$ for all $A \in \mathcal{A}$ and $B \in \mathcal{B}$.
In this section we will prove that the min-rank is super-multiplicative under the tensor product, under the assumption that at least one of the matrix subspaces is \emph{diagonal}. (Without this assumption the min-rank can be strictly submultiplicative, see \autoref{ex:minrank-not-supermultiplicative}.)
We then apply this to the min-rank of powers of any diagonal subspace of matrices.
%\ILnote{Comment 12. Changed the parentheses above.}

%
%

\begin{lemma}\label{lem:minrank-diag-supermultiplicative}
    Let $\mathcal{A}$ and $\mathcal{B}$ be subspaces of matrices. Suppose that all matrices in $\mathcal{A}$ are diagonal.
    %
    Then %
$\minrnk(\mathcal{A}\otimes\mathcal{B})\geq\minrnk(\mathcal{A})\cdot\minrnk(\mathcal{B})$.
\end{lemma}
%
\begin{proof}%
    Let $C\in\mathcal{A}\otimes\mathcal{B}$ be a nonzero element. Then $C$ is of the form $C =\sum_{i=1}^\ell A_i\otimes B_i$ for some $A_i \in \mathcal{A}$, $B_i \in \mathcal{B}$ and $\ell \geq 1$. We may assume $\{A_i\}_{i=1}^\ell$ and $\{B_i\}_{i=1}^\ell$ are both linearly independent sets of matrices.  %

    Since every $A_i$ is diagonal, the matrix $C$ is block-diagonal with blocks $D_j$ given by $D_j = \sum_{i=1}^\ell (A_i)_{j,j} B_i$. Thus $\rank(C) = \sum_j \rank(D_j)$. For any $j$ such that $(A_1)_{j,j} \neq 0$, the matrix $D_j$ is a nontrivial linear combination of the $B_i$ and thus $\rank(D_j) \geq \minrnk(\mathcal{B})$. The number of $j$ such that $(A_1)_{j,j} \neq0$ is equal to $\rank(A_1)$ since $A_1$ is diagonal. %
    Thus
    \[
    \rank(C) = \sum_j \rank(D_j) \geq \sum_{\mathclap{j: (A_1)_{j,j} \neq 0}} \rank(D_j) %
    \geq \rank(A_1)\minrnk(\mathcal{B}) \geq \minrnk(\mathcal{A})\minrnk(\mathcal{B}),
    \]
    which proves the claim. %
    %
%
    %
%
    %
\end{proof}
%
%\ILnote{Comment 12. Changed in the example below the "is necessary" to "cannot be removed". We can also phrase it as: "The lemma cannot be made without any condition" if we prefer.}
\begin{example}\label{ex:minrank-not-supermultiplicative}
    The condition in \autoref{lem:minrank-diag-supermultiplicative} that at least one of the matrix subspaces consists of diagonal matrices cannot be removed. To see this we may take the subspace of matrices over the reals
    \[
\mathcal{A} = \Bigl\{\begin{pmatrix}
    a & b\\-b & a
\end{pmatrix} : a,b \in \RR\Bigr\}.
    \]
    Then $\minrnk(\mathcal{A}) = 2$. However, $\mathcal{A}^{\otimes 2}$ contains the element
    \[
    \begin{pmatrix}
        1 & 0\\0 & 1
    \end{pmatrix}^{\otimes 2}
    +
    \begin{pmatrix}
        0 & 1\\-1 & 0
    \end{pmatrix}^{\otimes 2}
     = \begin{pmatrix}
         1 & 0 & 0 & 1\\
         0 & 1 & -1 & 0\\
         0 & -1 & 1 & 0\\
         1 & 0 & 0 & 1
     \end{pmatrix}
    \]
    which has rank $2 < 2^2$, so $\minrnk(\mathcal{A}^{\otimes 2}) < \minrnk(\mathcal{A})^2$.
\end{example}


\begin{lemma}\label{claim:large-rank-lin-comb-tensorprods}
    Let $\mathcal{A}$ %
    be a subspace of matrices. Suppose that all matrices in $\mathcal{A}$ are diagonal. %
    Then for all $k\in\NN$ we have $\minrnk(\mathcal{A}^{\otimes k})\geq \minrnk(\mathcal{A})^k$.
\end{lemma}
\begin{proof}
    Repeatedly apply \autoref{lem:minrank-diag-supermultiplicative}.
\end{proof}
%
%

%

%

%
In the rest of this section we prove a slightly more general version of \autoref{claim:large-rank-lin-comb-tensorprods}.
The reason that we need to do this is that we want to take the output of \autoref{claim:matrices-into-diagonal-with-large-rank} (a collection of linearly independent matrices, such that restricted to specified rows and columns the matrices are diagonal and have large min-rank, or are zero) and use the super-multiplicativity property of \autoref{claim:large-rank-lin-comb-tensorprods} to construct a large collection of matrices (that are tensor products of the original matrices) with large min-rank. However, we cannot do this directly as some of the matrices after restriction are zero. The lemma we prove here deals with that.

%\refnote{19}{The notation $A_{I\times J}$ has been used many times up to this point} 
%\Jshort{Removed it}
%Recall that for any matrix $A \in \FF^{n_1} \otimes \FF^{n_2}$ and subsets $I\subseteq [n_1]$, $J\subseteq [n_2]$ we denote by $A_{I\times J}$ the submatrix of $A$ indexed by rows $I$ and columns $J$. For any subset $X \subseteq \FF^{n_1} \otimes \FF^{n_2}$ we denote by $X_{I \times J}$ the set $\{A_{I\times J} : A \in X\}$.


%
%

%
\begin{restatable}{lemma}{lemC}
%
\label{claim:linear-combinations-of-mixed-tensorprods-high-rank}
    Let $X = \{B_1, \ldots, B_b, C_{b+1} \ldots, C_c\} \subseteq \FF^{n_1} \otimes \FF^{n_2}$ be a set of linearly independent matrices. Let $I\subseteq [n_1]$, $J \subseteq [n_2]$, $|I|=|J|$, such that 
    \begin{itemize}
        \item the matrices $(B_i)_{I\times J}$, for $i = 1, \ldots, b$, are diagonal and linearly independent,
        \item the matrices $(C_i)_{I\times J}$, for $i = b+1, \ldots, c$, are zero.
    \end{itemize}
    For any $m \geq \ell\geq 1$, 
    let $Y$ %
    be the set of all order-$m$ Kronecker products of elements of~$X$ such that at least $\ell$ of the factors are from $B_1, \ldots, B_b$. Then %
    \[
    \minrnk(Y) \geq \minrnk(X_{I\times J})^{\ell}.
    \]
%
\end{restatable}
 
%
%
%
%
%
%

%
\begin{proof}
    %
    %
    Let $A$ be any nontrivial linear combination of elements of $Y$.
    Then
    \[
    A=\sum_{t=1}^r a_t\, D_{i_{t,1}}\otimes\cdots\otimes D_{i_{t,m}}
    \]
    for some nonzero $a_t \in \FF$ and $D_{i_{t, j}} \in X$. Every $D_{i_{t, j}}$ is either in $\{B_1, \ldots, B_b\}$, in which case we will call it a $B$-matrix, or in $\{C_{b+1}, \ldots, C_c\}$, in which case we will call it a $C$-matrix.
    %
    %
    %
    %
    Let $k$ be the maximal number of $B$-matrices that appear as factors in any of the summands. Then $k\geq \ell$. 
    Consider any summand with the maximal number of $B$-matrices in its factors. Suppose it is the first summand. %
    %
    For simplicity of notation, we assume the first $k$ factors of this summand are $B$-matrices and the last $m - k$ factors are $C$-matrices: %
    %
    \[
    D_{i_{1,1}} \otimes \cdots \otimes D_{i_{1,m}}  = B_{i_{1,1}}\otimes\cdots\otimes B_{i_{1,k}}\otimes C_{i_{1,k+1}}\otimes\cdots\otimes C_{i_{1,m}}.
    \]
    We restrict the first $k$ factors of each summand in the definition of $A$ to $I\times J$ to get %\refnote{20}{add "of each summand in the definition"}\ILnote{added it}%
    \[
    A' = \sum_{t=1}^r a_t \, (D_{i_{t, 1}})_{I\times J} \otimes \cdots \otimes (D_{i_{t, k}})_{I\times J} \otimes D_{i_{t, k+1}} \otimes \cdots \otimes D_{i_{t, m}}.
    \]
    If there is a $C$-matrix among $D_{i_{t, 1}}, \ldots, D_{i_{t, k}}$, then $(D_{i_{t, 1}})_{I\times J} \otimes \cdots \otimes (D_{i_{t, k}})_{I\times J}$ is zero. 
    %
    Moreover, nonzero summands cannot have any $B$-matrix among $D_{i_{t, k+1}}, \ldots, D_{i_{t, m}}$, as this would contradict the maximality of $k$. So in any non-zero summand the $D_{i_{t, 1}}, \ldots, D_{i_{t, k}}$ are $B$-matrices and the $D_{i_{t, k+1}}, \ldots, D_{i_{t, m}}$ are $C$-matrices. %So in any non-zero summand has $k$ $B$-matrices and then $m-k$ $C$-matrices.
    %
    %
    %
    We know that the first $k$ factors of the first summand of $A$ are $B$-matrices, so $A'$ is nonzero by the linear-independence of $\{(B_i)_{I\times J}\}^{\otimes k}\otimes\{C_i\}^{\otimes m-k}$.
    %
    %
%
    %
%
    %
%
    Using the notation $B'_i = (B_i)_{I\times J}$, we may rewrite $A'$ as
    %
    % \[
    % A' = \sum_{{j}\in[b]^k}B'_{j_1}\otimes\cdots\otimes B'_{j_k}\otimes\Bigl(\sum_{t=1}^r\delta_{i_{t,1}=j_1}\cdots\delta_{i_{t,k}=j_k}\cdot a_t\, C_{i_{t,k+1}}\otimes\cdots\otimes C_{i_{t,m}}\Bigr).
    % \]
    \[
    A' = \sum_{{j}\in[b]^k}B'_{j_1}\otimes\cdots\otimes B'_{j_k}\otimes M_j,
    \]
    where
    \[
    M_j = \sum_{t=1}^r\delta_{i_{t,1}=j_1}\cdots\delta_{i_{t,k}=j_k}\cdot a_t\, C_{i_{t,k+1}}\otimes\cdots\otimes C_{i_{t,m}}.
    \]
    %We have $a_1\neq0$ so for the multi-index that accompanies it, $(i_{1,1},\dots,i_{1,k})\in[b]^k$, the following is a non-trivial linear combination, meaning a non-zero matrix:
    Note that $S = M_{(i_{1,1},\dots,i_{1,k})}$ is a nonzero matrix, because 
    %Let $\tilde j=(i_{1,1},\dots,i_{1,k})\in[b]^k$ be the first $k$ indices of the first summand in $A$. 
    $a_1\neq0$ and the set $\{C_i\}_{i\in\{b+1,\dots,c\}}^{\otimes(m-k)}$ is linearly independent. %
    % %
    % %
    % \begin{equation}\label{eq:block-lin-comb}
    % M = \sum_{t=1}^r\delta_{i_{t,1}=i_{1,1}}\cdots\delta_{i_{t,k}=i_{1,k}}\cdot a_t\, C_{i_{t,k+1}}\otimes\cdots\otimes C_{i_{t,m}}
    % \end{equation} 
    % is nonzero. 
    %
    %
    %
    %
%
    %
    % We may think of $A'$ as a block matrix composed of $n_1^{m-k}\times n_2^{m-k}$ blocks of size $|I|^{k}\times |J|^{k}$ \refnote{21}{$n_1^{m-k}\times n_2^{m-k}$ block matrix was confusing. Suggested alternative phraseing}\ILnote{changed to the suggestion}. That is, the ``outer'' blocks correspond to linear combinations of products of $C_i$'s, and the ``inner'' blocks correspond to linear combinations of tensor products of $B'_i$'s. 
    %Consider the nonzero linear combination \autoref{eq:block-lin-comb} corresponding to ${\tilde j}$ and denote it by $S$. 
%
    Suppose $S_{x,y}\neq 0$.
    Then
    \[
    \sum_{{j}\in[b]^k}B'_{j_1}\otimes\cdots\otimes B'_{j_k} \cdot (M_j)_{x,y}
    \]
    is a nontrivial linear combination, so it has rank at least 
    \[
    \minrnk(X_{I\times J})^k\geq \minrnk(X_{I\times J})^\ell.
    \]
    This matrix is a submatrix of $A'$, and $A'$ is a submatrix of $A$. This proves the claim.
    %
    % Assume without loss of generality that the $(1,1)$ entry of $S$ is not zero ($S_{11}\neq0$). Block $(1,1)$ in $A'$ is a linear combination of the order-$k$ tensor products of the $B'_i$'s with a nonzero coefficient ($S_{11}$) for $B'_{i_{1,1}}\otimes\cdots\otimes B'_{i_{1,k}}$, meaning this linear combination is nontrivial. \autoref{claim:large-rank-lin-comb-tensorprods} then implies that this block has rank at least $\minrnk(X_{I\times J})^k\geq \minrnk(X_{I\times J})^\ell$ (by $k\geq\ell$) which in turn lower bounds the rank of $A'$ and thus of $A$. %, which proves the claim.
    %\refnote{22}{Page 27, last few lines. These were not very clear to me and I had to read a few times. I hope I now understand, but this step could be expanded slightly.} \Jshort{We rephrased the last part of the proof to make it more concrete.}
\end{proof}

%
%
%
%

%

%
    
%




%
%
    
%
%
    
%
    
%
%
%
%
%
%
%
%
%
%
%
%

%
%
%
%

%
%
%

%
    
%
%
%
%
%
%
%
%
%
%
%
%
%
%

%
%
%
%


\subsection{Narrow tensors have maximal asymptotic subrank}
\label{sec:narrow}\label{subsec:main-unb}



%

%


%
We will now %
%
prove the main theorem of this section, which says that for every constant~$c$, and large enough $n_1, n_2$, all tensors in $\FF^{n_1} \otimes \FF^{n_2} \otimes \FF^{c}$ have maximal asymptotic subrank.

\begin{theorem}%
\label{theorem:concise-cn1n2-assubrank-c}
    For every $c\geq2$ there is $N(c)\in\mathbb{N}$ such that for any $n_1,n_2\in\NN$ such that $\max\{n_1,n_2\}>N(c)$, every concise tensor $T\in\mathbb{F}^{n_1}\otimes\mathbb{F}^{n_2}\otimes\mathbb{F}^c$ satisfies $\asympsubrank(T)=c$.
\end{theorem}

%The proof of \autoref{theorem:concise-cn1n2-assubrank-c} relies on min-rank lemmas that we proved in
%\autoref{subsec:minrank}.
%For convenience of the reader we briefly restate here the ones we need. 
% \refnote{23}{Having the lemmas repeated is indeed convenient -- but it seems like the repetition is excessive}
% \Jnote{I removed them.}

% \lemA*

% \lemBx*

% \lemC*

% \lemD*

%

% \Jnote{Did we fix the Chernoff bound proof below?}
% \ILnote{I didn't try to see how you got to the bound that appears below ($m\geq8c^3$ and $\ell=\tfrac1{2c}m$). We can either do that or revert back to the more loose bound of having $\ell=\tfrac1{100c}m$ and $m>5\ln(4)c$. For the proof of this looser bound see the screenshot from our backup, I upleoded it here as the file "Chernoff\_old.png".}

\begin{proof}%[Proof of \autoref{theorem:concise-cn1n2-assubrank-c}]
Suppose $n_1 \geq n_2$. Suppose $n_1 > N(c)$ for an $N(c)$ to be determined in the proof.
Let $T \in \FF^{n_1} \otimes \FF^{n_2} \otimes \FF^{c}$ be concise.
Let $\mathcal{A} \subseteq \FF^{n_1 \times n_2}$ be the 3-slice span of $T$.
By \autoref{cla:high-rk} we have %there is a slice of $T$ of rank at least $\frac{1}{c} n_1$. In other words, 
$\maxrank(\mathcal{A}) \geq \frac{1}{c} n_1$.
We apply \autoref{lem:minrk_diag} to $\mathcal{A}$, giving matrices $B_1, \ldots, B_b$, $B_{b+1}, \ldots, B_c$ and a subset $J \subseteq [\min_\ell n_\ell]$.
To these, we apply \autoref{claim:linear-combinations-of-mixed-tensorprods-high-rank} for $m \geq \ell \geq 1$ (to be specified later) 
%$m = 8c^3$ and $\ell = \tfrac{1}{2c} m$ 
to obtain a set~$Y$ of slices from~$T^{\boxtimes m}$
%tensor $T'$ with 3-slice span $\mathcal{A}'$ 
satisfying
\beq\label{eq:mrkY}
\minrank(Y) \geq \big(\eps(c)\maxrank(\mathcal{A})\big)^\ell \geq (\eps(c)\tfrac1c n_1)^{\ell}. %\tfrac{1}{100c} m}.
\eeq
The size of~$Y$ equals the number of $m$-tuples in $[c]^m$ containing at least~$\ell$ elements from~$[b]$.
Taking $m \geq 8c$ and $\ell = \tfrac{1}{2c} m$, it then follows from the Chernoff bound (e.g., \cite[Equation~(7)]{MR1045520}) that~$|Y|\geq \tfrac12c^m$.

Since~$T$ is concise, so is~$T^{\boxtimes m}$ and therefore, the set~$Y$ is linearly independent.
Let~$T'$ be the tensor with 3-slices given by the elements in~$Y$.
Then $T^{\boxtimes m} \geq T'$.
Provided
\begin{equation}\label{eq:minrank}
\minrnk(Y) \geq 2(\tfrac12c^{m})^2,
\end{equation}
it follows from \autoref{cla-high-rank-high-subrank} that
\beq\label{eq:QTmbound}
\subrank(T^{\boxtimes m}) \geq \subrank(T') \geq \tfrac12c^m.
\eeq
By our choice of~$\ell$ and~\eqref{eq:mrkY}, we have that~\eqref{eq:minrank} holds if
\beqn
\big(\eps(c)\tfrac1c n_1\big)^{\tfrac{1}{2c}m}
\geq
2(\tfrac12c^{m})^2.
\eeqn
It is clear that there is an~$N(c)\in\NN$ such that this is satisfied for any~$n_1>N(c)$ and every~$m$.
Letting~$m$ go to infinity we then get from~\eqref{eq:QTmbound} that 
\beqn
\asympsubrank(T) = \lim_{m\to \infty}\subrank(T^{\boxtimes m})^{1/m} \geq c.\qedhere
\eeqn
\end{proof}

\begin{remark}
    To conclude this section, let us note that a similar approach fails for $k$-tensors, for $k>3$. The motivating question mentioned at the beginning of \autoref{sec:lower-bound-imbalanced}, whether there exist accumulation points in the image of asymptotic slicerank (over~$\CC$ and finite fields), makes sense for $k$-tensors for any $k$. In order to show there cannot exist such points in the same way for $k>3$ we would need to show there are only finitely many formats that allow asymptotic sliceranks strictly smaller than the maximal constant $c$. For $k=3$ we showed it for asymptotic subrank, which lower bounds asymptotic slicerank, but this is false for asymptotic subrank and any~$k>3$. For $k=3$ conciseness implies the formats that allow for a constant asymptotic subrank are only $(n_1,n_2,c)$ for a constant $c$ and $n_1,n_2$ growing. But for $k=4$ and above there can be formats of the form $(n_1,n_2,\dots,n_{k-2},c_1,c_2)$ for $c_1,c_2$ constants and $n_1,n_2,\dots,n_{k-2}$ growing. For these formats the straightforward generalisation 
    %
    of \autoref{theorem:concise-cn1n2-assubrank-c}, asking to show that for every $c_1,c_2$ constants there would be $N$ for which any concise $(n_1,n_2,\dots,n_{k-2},c_1,c_2)$-tensor with $n_1,n_2,\dots,n_{k-2}>N$ has asymptotic subrank $\min\{c_1,c_2\}$ does not hold for $k$-tensors, for any $k>3$:
    
    The $4$-tensor $\id_n\otimes \id_2=\sum_{i=1}^n e_i\otimes e_i\otimes e_1\otimes e_1+\sum_{i=1}^n e_i\otimes e_i\otimes e_2\otimes e_2$ is a concise $4$-tensor of format $(n,n,2,2)$ (for each $i\in[4]$, the set of $i$-slices is linearly independent) but it is a tensor-product as a $2$-tensor (for the bipartition $\{\{1,2\},\{3,4\}\}$) so its subrank and its asymptotic subrank are $1$ as a $2$-tensor and thus as a $4$-tensor as well.
    
    This example can be generalised to any $k>3$ and any $c_1=c_2=:c$ as $\langle n\rangle_{k-2}\otimes \id_c$ where $\langle n\rangle_{k-2}=\sum_{i=1}^n e_i\otimes\cdots\otimes e_i$ is the diagonal (unit) $(k-2)$-tensor of size $n$.
\end{remark}


\subsection{Concrete case: smallest dimension equals two} %
\label{subsec:dimtwo}

In the previous section we proved for any constant $c$ that for large enough $n_1,n_2$ all concise tensors in $\FF^{n_1} \otimes \FF^{n_2} \otimes \FF^c$ have asymptotic subrank $c$. In the special case $c=2$ we can prove the following stronger statement in which we make the constant $N(2)$ explicit, namely $N(2) = 2$, and we are able to replace asymptotic subrank by subrank.
%

\begin{theorem}\label{local-dim-2}
    Let $n_1,n_2 > 2$ and let $T\in\mathbb{F}^{n_1}\otimes\mathbb{F}^{n_2}\otimes\mathbb{F}^2$ be a concise tensor. Then
    \[\subrank(T)=2.\]
\end{theorem}
%
%

\begin{example}
    In \autoref{local-dim-2} it is required that at least one of the dimensions~$n_1, n_2$ is strictly larger than 2.
    Namely, there are concise tensors $T \in \FF^2 \otimes \FF^2 \otimes \FF^2$ such that $\subrank(T) = 1$, for instance the tensor $T=e_1 \otimes e_2 \otimes e_2 + e_2 \otimes e_1 \otimes e_2 + e_2 \otimes e_2 \otimes e_1$ that is often called the ``W-tensor''. 
\end{example}

\begin{remark}
    \autoref{local-dim-2} leaves open what happens for $n_1=2$ and $n_2 > 2$. 
    Over the complex numbers the classification %
    of Miyake and Verstraete \cite{miake2004multipartite22n} shows that indeed for every $n>2$ and every concise tensor $T \in \FF^{n} \otimes \FF^2 \otimes \FF^2$, $\subrank(T)=2$.
    %
    Over any algebraically closed field, Tobler \cite[Section 4.2]{Tobler1997Spezi-6235} showed that for every $n_2>n_1\geq2$ and every concise tensor $T \in \FF^{n_1} \otimes \FF^{n_2} \otimes \FF^2$, $\bordersubrank(T)=2$. This in particular implies that the asymptotic subrank is $2$ (\autoref{lem:border-asymp-subrank}). %
    
    %
    %

    %
    %
\end{remark}

%
%
%
%
%

%

%

%
%
    
%
%

%
\begin{lemma}\label{cl:1000andtus1Sub2} 
For any $s,t,u \in \FF$, the tensor with slices 
\[
A=\begin{pmatrix}1 & 0\\ 0 & 0\end{pmatrix}, B=\begin{pmatrix}t & u\\s & 1\end{pmatrix}
\]
has subrank $2$.
\end{lemma}
\begin{proof} We first subtract $s$ times the second column from the first column and $u$ times the second row from the first row to get 
\[
A=\begin{pmatrix}1 & 0\\ 0 & 0\end{pmatrix}, B=\begin{pmatrix}t' & 0\\ 0 & 1\end{pmatrix}
\]
for some $t'$.
We then subtract $t'\cdot A$ from $B$ to get 
\[
A=\begin{pmatrix}1 & 0\\ 0 & 0\end{pmatrix}, B=\begin{pmatrix}0 & 0\\ 0 & 1\end{pmatrix}.
\]
%
This is $\langle2\rangle$, the diagonal tensor of rank $2$.
\end{proof}

%
\begin{proof}[Proof of \autoref{local-dim-2}]
    Let $A,B$ be the two $n_1\times n_2$ slices. We diagonalize the $n_1\times 2n_2$ matrix $(A|B)$ by diagonalizing $A$ and doing the same row and column operations on~$B$. Observe that $\rank(A|B)=n_1$ as $T$ is concise. 
    Assume w.l.o.g.\ that the top $r\times r$ minor of $A$ is the identity and the rest is~$0$. Specifically, the last $n_1-r$ rows in it are $0$. %
%
    %
    We can now continue and diagonalize the last $n_1-r$ rows of $B$, noting that this will not change the structure of~$A$ (we are only doing row operations now).
    We will get a subset~$I$ of $B$'s columns such that restricted to some of the last $n_1-r$ rows, $B|_{\{r+1,\dots,r+\ell\}\times I}$ is the identity. 
    By doing more row operations we can assume that the first $r$ rows in the columns~$I$ in $B$ are $0$ (the $n_1-r$ rows in $A$ are zeros so we are free to subtract scaled versions of these rows from above rows to zero out the desired cells in $B$). Here's an illustration of our two $n_1\times n_2$ matrices (the dashed lines separate the first~$r$ rows and columns from the rest and the $I$ indices of columns are indicated beneath matrix $B$. Asterisks signify arbitrary values that we cannot assume anything about). In particular, a $0$ above the dashed line in $B$ represents an $r\times 1$ columns of zeros:
    \[
   A=\left(
   \begin{array}{c;{2pt/2pt}c}
       \id_r & 0 \\ \hdashline[2pt/2pt]
       0 & 0
   \end{array}
   \right) \quad
   B=\genfrac{}{}{0pt}{}{\left(
   \begin{array}{cccccc;{2pt/2pt}c}
        *&0&*&0&\cdots & 0& \\ \hdashline[2pt/2pt]
        *&1&*&0&\cdots&0&*\\
        *&0&*&1&\cdots&0&*\\
        *&\vdots&*&\vdots&\ddots&\vdots&\vdots\\
        *&0&*&0&\cdots&1&*\\
        *&0&*&0&\cdots&0&*
   \end{array}
   \right)}{ \begin{array}{ccccccc}
        &I_1&&I_2&\cdots&I_\ell&\hspace{1pt}  
   \end{array}}%
    \]
%
    We now have three cases to consider:

%
    %
    Case 1: $1<r<\min\{n_1,n_2\}$.
    Let $i\in I$ be arbitrary but different than $r$.
    Note that $B_{r+i,i}=1$ (it is part of the identity matrix on columns $I$). 
    Consider the restriction of our tensor to rows $r,r+i$ and columns $r$ and $i$. 
    This gives the tensor with slices
    %
    \[
    A = \begin{pmatrix}1 & 0\\ 0 & 0\end{pmatrix},
    B=\begin{pmatrix}t & 0\\ s & 1\end{pmatrix}
    \]
    for some $s$ and $t$. From \autoref{cl:1000andtus1Sub2} we know that this tensor has subrank 2.

    %
    Case 2: $r=\min\{n_1,n_2\}$.
We consider two subcases:
    \begin{itemize}
    \item Case 2a: The matrix $B$ has a nonzero nondiagonal coefficient. Assume w.l.o.g.\ that $B_{1,2}=t \neq 0$.  
Restrict our tensor to rows 1,3 and columns 2,3 ($n_1,n_2\geq3$).
We get slices
\[
A=\begin{pmatrix}0 & 0\\ 0 & 1\end{pmatrix}, B=\begin{pmatrix}t & *\\ * & *\end{pmatrix}.
\]
Using \autoref{cl:1000andtus1Sub2}, as $t\neq 0$, this tensor has subrank $2$.

    \item Case 2b: All nondiagonal coefficients of $B$ are $0$ ($\forall i\neq j: B_{i,j}=0$). 
As the tensor is concise, the matrix~$B$ is not a multiple of the matrix $A$. The diagonal coefficients of $A$ are all~$1$. Hence, w.l.o.g., $B_{1,1}\neq B_{2,2}$.
Restricting to rows and columns $\{1,2\}$ we get 
\[
A=\begin{pmatrix}1 & 0\\ 0 & 1\end{pmatrix},
B=\begin{pmatrix}a & 0\\ 0 & b\end{pmatrix}.
\]
Taking appropriate linear combinations of $A$ and $B$ we get 
\[
A=\begin{pmatrix}0 & 0\\ 0 & 1\end{pmatrix}, B=\begin{pmatrix}1 & 0\\ 0 & 0\end{pmatrix}.
\]
\end{itemize}

    %
    Case 3: $r=1$.
    In this case we repeat the argument above with $B$ in place of $A$. As $\rank(A)=1$, $n_1,n_2>2$, and the tensor is concise we must have $\rank(B)>1$.
    %
\end{proof}

\begin{remark}
    The proof of \autoref{local-dim-2} only uses conciseness with respect to two directions, namely that the matrices $A$ and $B$ are linearly independent (Case 2b) and that the matrix $(A|B)$ has full rank (all cases). An example (that is not concise) is
    \[
    A=\begin{pmatrix}1 & 1\\ 1 & 1\end{pmatrix}, B=\begin{pmatrix}1 & 1\\ 2 & 2\end{pmatrix}.
    \]
\end{remark}

%

%
%

%

%

%
%


%
%

%

%
%
%
%

%
%
%
%

%
%
%
%
%
%
    
%
%

%
%
%

%
%
%
%
%

%

%

%
%
    
%
%

%
%
%
%
%
%
%
%
%
%
%
%
%
%
%
%
%
%
%
%

%
%
%
%
%
%
%
%
%
%
%
%
%
%
%
%
%
%
%
%
%
%
%
%
%
%
%
%
%
%
%

%
%
%
%
%
%
%
%
%
%
%
%
%

%
%
%
%
%
%
%
%
%
%
%

%
%
%
%
%
%
%
%
%
%
%
%

%
%
%
%
%

%
%
%
%
%
%

%


%
%
%
%
%
%
%


%
%
%

%
%
%
%
%
%
%
%
%

%
%
%

%

%

%

%
%
%
%
%

%

%
%
    
%
%
%

%
%

%
%

%
%
%

%
%
%
%
%

%

%

%

%

%
%
%
%

%

%
%
%

%
%

%
%

%
%
%

%
%
%

%
%
%


%
%
%

%
%

%


%

%

%

%

%
%
%
%
%

%
%

%
%



%


%
%
%

%
%
%


%
%
%
%

%
%
%
%

%
%
%

%
%
%

%

%

%

%
%

%
%
%

%
%
%
%

%
%
%
%
%

%

%
%
%

%
%
%
%
%

%
%
%
%
%
%
%

%

%


%
%
%
%
%
%
%
%
%
%
%

%
%
%

%
%

%


%
%
%
%
%
%
%
%
%
%
%
%
%
%
%
%
%
%
%
%
%
%
%
%
%
%
%
%

%
%
%
%


%
%
%
%
%

%

%
%
    
%
    
%
    
%
    
%

%
%

%

%
    

%
%
    
%
%
%

%
 
%
%

%
 
%
%
%
 
%

%
 

%
%
%
%

%
%
    
%

%

%
%

%
%
    
%
%
    
%
    
%
%
%
%
%
%
%
%
%
%
%
%

%
%
%
%

%
%
%

%
    
%
%
%
%
%
%
%
%
%
%
%
%
%
%

%
%
%
%

%
%
    
%
    
%
    
%

%
%

%
%
%
%
    
%
    
%
%




\section{Discreteness of asymptotic tensor ranks}
\label{sec:discreteness}
%
%

In this section we prove a general sufficient condition for a tensor parameter to be discrete (using the lower bounds on the asymptotic subrank of concise tensors that we proved in the previous sections, \autoref{sec:max-rank} %
and \autoref{sec:lower-bound-imbalanced}). Then we show that this condition is satisfied (and thus obtained discreteness) for several natural tensor parameters and regimes.

\subsection{General sufficient condition for discreteness}

We recall that two tensors $S$ and $T$ are called equivalent if $S \geq T$ and $T \geq S$. Every tensor is equivalent to a concise tensor (\autoref{lem:equiv-concise}). We will be interested in functions~$f$ on tensors that are invariant under equivalence, meaning that if two tensors~$S$ and $T$ are equivalent, then $f(S) = f(T)$. This is a very mild condition, in particular any function that is monotone under restriction is also invariant under equivalence.

\begin{theorem}
    \label{th:gen-no-acc}
    Let $\FF$ be any field. Let $\mathcal{C}$ be any subset of order-three tensors over~$\FF$, such that for any $T \in \mathcal{C}$ there is an $S \in \mathcal{C}$ that is concise and equivalent to $T$.
    Let~$f : \mathcal{C} \to \RR_{\geq0}$ be any function such that
    \begin{enumerate}[\upshape (i)]
    \item\label{item:equiv} $f$ does not change under equivalence of tensors
    \item\label{item:lb} $f(T) \geq \asympsubrank(T)$ for every $T \in \mathcal{C}$
    \item\label{item:fin} For every $n_1, n_2, n_3\in \NN$,  $f$ takes finitely many values on $\mathcal{C} \cap (\FF^{n_1} \otimes \FF^{n_2} \otimes \FF^{n_3})$ 
    \item\label{item:ub} For every $n_1, n_2, n_3\in \NN$, $f(T) \leq \min_\ell n_\ell$ for every $T \in \mathcal{C} \cap (\FF^{n_1} \otimes \FF^{n_2} \otimes \FF^{n_3})$.
    \end{enumerate}
    Then the set of values that $f$ takes on $\mathcal{C}$ has no accumulation points.
\end{theorem}

%

\begin{proof}
    Suppose we have, for every $i \in \NN$, three numbers $a_i, b_i, c_i \in \NN$ and a tensor
    %
    $T_i \in \mathcal{C} \cap (\FF^{a_i} \otimes \FF^{b_i} \otimes \FF^{c_i})$ in such a way that the set $\{f(T_i): i \in \NN\}$ is infinite. We will show that the $f(T_i)$ cannot converge.
    
    We may assume that every tensor $T_i$ is concise by condition~(\ref{item:equiv}) and \autoref{lem:equiv-concise}.
    %
    %
    %
    By condition~(\ref{item:fin}), $f$ takes only finitely many values on $\mathcal{C} \cap (\FF^{n_1} \otimes \FF^{n_2} \otimes \FF^{n_3})$ for any choice of $n_1, n_2, n_3$, so the set of triples $\{(a_i, b_i, c_i): i \in \NN\}$ must be infinite. %
    
    We consider two cases. First,  suppose $\left\{\min \{a_i, b_i, c_i\}\right\}_i$ is bounded. Suppose for simplicity of notation that $a_i = \min \{a_i, b_i, c_i\}$ for every $i \in \NN$. Then $a_i$ takes only finitely many values. Take any infinite subsequence of the $T_i$ such that  $a_i$ is constant. %
    Then both $b_i$ and $c_i$ are unbounded, because $b_i\leq a_i c_i$ and $c_i\leq a_i b_i$ by \autoref{cl:concise-formats}. When~$b_i$ and $c_i$ are large enough we have $\asympsubrank(T_i) = a_i$ by \autoref{theorem:concise-cn1n2-assubrank-c}, %
    and hence $f(T_i) = a_i$ by condition~(\ref{item:lb}) and~(\ref{item:ub}). We conclude that $f(T_i)$ takes only finitely many values, which contradicts our assumption that $\{f(T_i): i \in \NN\}$ is infinite.
    
    Second, suppose $\min \{a_i, b_i, c_i\}$ is unbounded. 
    We have $f(T_i) \geq \asympsubrank(T_i)$ by condition~(\ref{item:lb}). We have $\asympsubrank(T_i) \geq \min \{a_i, b_i, c_i\}^{1/3}$ by \autoref{th:cube-root-lb}. It then follows that the values $f(T_i)$ do not converge.
\end{proof}

%

\subsection{Asymptotic subrank and asymptotic slice rank}

We will now use \autoref{th:gen-no-acc} to prove a discreteness theorem for asymptotic slice rank and asymptotic subrank. 

Before stating our theorem, we recall the notions of tight and oblique tensors, which are used in one part of the theorem. 
Let $T \in \FF^{n_1} \otimes \FF^{n_2} \otimes \FF^{n_3}$.
We recall that the support of $T$ is the set $\supp(T) = \{(i,j,k) : T_{i,j,k} \neq 0\}$.
%
The tensor $T$ is called \emph{tight} if (after an appropriate basis transformation) there exist injective maps $f_\ell : [n_\ell] \to \ZZ$ such that for every $(i_1,i_2,i_3) \in \supp(T)$ we have $f_1(i_1) + f_2(i_2) + f_3(i_3) = 0$. 
%
%
The tensor $T$ is called \emph{oblique} if (after an appropriate basis transformation) $\supp(T)$ is an antichain in the product order (i.e.~pointwise order) on $[n_1] \times [n_2] \times [n_3]$. Tight tensors are oblique.

\begin{theorem}
\label{th:no-acc}
\hfill
\begin{enumerate}[\upshape(i)]
\item Let $\mathbb{F}$ be any field and $S \subseteq \FF$ any finite subset. 
%\begin{itemize}
%\item 
Then the asymptotic subrank and asymptotic slice rank on order-three tensors over $S$ have no accumulation points.
%\item The asymptotic slice rank on order-three tensors over $S$ has no accumulation points.
%\end{itemize}
\item %Let $\mathbb{F}$ be the complex numbers. 
%\begin{itemize}
%\item 
The asymptotic slice rank on order-three tensors over~$\CC$ has no accumulation points.
%\end{itemize}
\item Let $\mathbb{F}$ be any field.
\begin{itemize}
\item The asymptotic subrank and asymptotic slice rank on tight order-three tensors over $\mathbb{F}$ have no accumulation points.
\item The asymptotic slice rank on oblique order-three tensors over $\mathbb{F}$ has no accumulation points.
\end{itemize}
\end{enumerate}
\end{theorem}

\begin{proof}
    The proof of each case follows from \autoref{th:gen-no-acc}. As slice rank is lower bounded by subrank, $f(T) \geq \asympsubrank(T)$ for every $T \in \mathcal{C}$, over any field $\FF$. Thus, conditions (i), (ii) and (iv) of  \autoref{th:gen-no-acc} are fulfilled. We next show that condition (iii) of \autoref{th:gen-no-acc} holds for each of the cases in \autoref{th:no-acc}, i.e. that in each of these cases $f$ takes only finitely many values on $\mathcal{C} \cap (\FF^{n_1} \otimes \FF^{n_2} \otimes \FF^{n_3})$.

    (i) Let $f$ be the asymptotic subrank or the asymptotic slice rank and let $\mathcal{C}$ be the set of all order-three tensors over the finite set $S$. Indeed $\mathcal{C}$ is closed under taking the concise version of any element by \autoref{lem:equiv-concise}, since any tensor is equivalent to a concise subtensor (which in particular again has coefficients in $S$). Then $f$ takes only finitely many values on $\mathcal{C} \cap (\FF^{n_1} \otimes \FF^{n_2} \otimes \FF^{n_3})$ simply because the set of tensors in $\FF^{n_1} \otimes \FF^{n_2} \otimes \FF^{n_3}$ with coefficients in $S$ is finite. 
    
    (ii) Let $f$ be the asymptotic slice rank and let~$\mathcal{C}$ be the set of all order-three tensors over~$\CC$.  Fix dimensions $n_1,n_2$ and $n_3$. By the characterization of asymptotic slice rank in terms of entanglement (moment) polytopes (via the quantum functionals) \cite[Corollary 5.7]{MR4495838} and the fact that for fixed $n_1,n_2,n_3$ there are finitely many such polytopes \cite[Theorem 2]{walter2013entanglement}, it holds that $f$ takes only finitely many values on $\mathcal{C} \cap (\FF^{n_1} \otimes \FF^{n_2} \otimes \FF^{n_3})$. 
    
    (iii) Let $f$ be the asymptotic subrank or the asymptotic slice rank and let $\mathcal{C}$ be the set of all tight order-three tensors over arbitrary $\FF$. Fix dimensions $n_1,n_2$ and $n_3$. By the characterization of asymptotic subrank in terms of the tight support (via the support functionals) \cite[Proposition 5.4]{strassen1991degeneration} and \cite[Theorem 5.9]{MR4495838}, and the fact there are only a finite number of possible supports, it holds that $f$ takes only finitely many values on $\mathcal{C} \cap (\FF^{n_1} \otimes \FF^{n_2} \otimes \FF^{n_3})$.
    
%    $f$ takes only finitely many values on $\mathcal{C} \cap (\FF^{n_1} \otimes \FF^{n_2} \otimes \FF^{n_3})$ because of the characterization of asymptotic subrank and asymptotic slice rank in terms of the tight support via the support functionals proven in \cite[Proposition 5.4]{strassen1991degeneration}. 

    Let $f$ be the asymptotic slice rank and let $\mathcal{C}$ be the set of all oblique order-three tensors over arbitrary $\FF$. By the characterization of asymptotic slice rank in terms of the oblique support \cite[Proposition 4]{sawin} it follows that $f$ takes only finitely many values on $\mathcal{C} \cap (\FF^{n_1} \otimes \FF^{n_2} \otimes \FF^{n_3})$. 
\end{proof}

\begin{remark}
    \autoref{th:no-acc} gives several regimes in which the asymptotic subrank and asymptotic slice rank are discrete, that is, between every two consecutive values there is a ``gap''.
    Using \autoref{ex:strassenNullAlgebra} we can say something about the size of these gaps. Namely,
    \autoref{ex:strassenNullAlgebra} gives for every large enough $n$ (namely $n\geq5$) a concise tensor $T \in \FF^n \otimes \FF^n \otimes \FF^n$ such that the asymptotic subrank and asymptotic slice rank of $T$ equal $2\sqrt{n-1}$. This means that the gap between two consecutive values is at most $2\sqrt{n-1} - 2\sqrt{n-2}$, which is at most $O(1/\sqrt{n})$.
    %
    %
    This shows in particular that the consecutive gaps of the asymptotic subrank and asymptotic slice rank tend to zero as $n$ grows.
\end{remark}

As sketched in the introduction, with similar ideas as above (but much simpler) we can prove the analogous statement for asymptotic rank over finite fields.
Recall that the asymptotic rank of a tensor~$T$ is defined as $\asymprank(T) = \lim_{n\to\infty} \tensorrank(T^{\boxtimes n})^{1/n}$.

\begin{theorem}\label{th:asympran-disc}
    Let $\FF$ be any field, $S\subseteq\FF$ a finite set, and $k \in \mathbb{N}$. The asymptotic rank on $k$-tensors over~$S$ has no accumulation points. 
    %
    In particular, the asymptotic rank on $k$-tensors over any finite field has no accumulation points.
\end{theorem}
%\ILnote{Changed the theorem (and proof) slightly to talk about any finite subset of any field, as we changed \ref{th:no-acc} above.}
\begin{proof} %\refnote{24}{The proof was essentially already given on page 3}\Jshort{I propose to ignore this, see response.}
We may assume all our tensors to be concise (\autoref{lem:equiv-concise}).
For a fixed format $(n_1, \ldots, n_k)$, there are only finitely many possible values of the asymptotic rank of tensors in $\FF^{n_1} \otimes \cdots \otimes \FF^{n_k}$ with coefficients from $S$ (simply because there are only finitely many tensors of that format with coefficients from $S$). Thus, to make a sequence of concise tensors $T_i$ such that the asymptotic rank $\asymprank(T_i)$ takes infinitely many different values, we need infinitely many different formats $(n_1,\ldots, n_k)$. In particular, $\max_\ell n_\ell$ needs to be unbounded. Since $\asymprank(T_i) \geq \max_\ell n_\ell$, the sequence $\asymprank(T_i)$ cannot converge.
\end{proof}


\section{Pivots of tensors and better lower bounds on the asymptotic subrank}\label{sec:lower-bound-asympsubrank}

In \autoref{sec:max-rank} we proved a cube-root lower bound on the asymptotic subrank of concise tensors, using a result on max-rank and matrix multiplication tensors.

In this section we develop a different strategy, using pivots, to first reprove the cube-root lower bound on asymptotic subrank of \autoref{sec:max-rank} and then prove a better lower bound under a symmetry condition that we introduce (pivot-matched).

%
%
%

%

Namely we show that the asymptotic subrank of a concise pivot-matched tensor is at least the square-root of the smallest dimension. In particular, this holds for symmetric concise tensors. %

%

%
%

%

%
%
%

%

%
%
%

%
%
%
%
%
%
%
%
%


%
%
%
%
%
%

%
%
%
%
%
%
%
%
%
%


%

%

\subsection{Pivots of tensors}
We will first discuss what a pivot of a matrix is and what the pivots of a matrix subspace are, then introduce the pivots of a tensor, then relate this notion to the border subrank and the asymptotic subrank.

\subsubsection{Pivot of a matrix and pivots of matrix subspaces}

Following Meshulam \cite{meshulam1985maximal}, we recall the notion of pivots of a matrix subspace and associated covering and packing parameters $\rho$ and $\sigma$, which are known to be equal by an application of Kőnig's theorem. In the next section we will use these notions in the context of tensors.

\begin{definition}
    For any $n_1 \times n_2$ matrix $A$ we call the smallest element $(i,j) \in [n_1] \times [n_2]$ in the lexicographic ordering on $[n_1] \times [n_2]$ for which $A_{i,j} \neq 0$ the \emph{pivot} of $A$, and we denote this element by $p(A)=(f(A),g(A))$. 
\end{definition}

%
%

\begin{lemma}\label{lem:pivot-set-invariant}
    Let $\mathcal{A}$ be a subspace of $n_1\times n_2$ matrices. There is a basis for $\mathcal{A}$, $\{A_1,\dots,A_{n_3}\}$, with distinct pivots $\forall i\neq j: p(A_i)\neq p(A_j)$.
    Additionally, for any such basis, the set of its matrices' pivots $\{p(A_1),\dots,p(A_{n_3})\}$ is the same.
\end{lemma}
To see this, vectorise the matrices row-by-row and stack them as the rows of a matrix $M$. Changing the basis amounts to invertible row operations on $M$ and having distinct pivots amounts to getting $M$ into its row-echelon form (up to the order of the rows). Uniqueness of the row-echelon form gives the uniqueness of the set of pivots.
%
    %
    %

    %
%

By this lemma we may define the set of pivots as a parameter of the linear subspace of matrices itself:

\begin{definition}\label{def:pivots-space}
    For any subspace of $n_1 \times n_2$ matrices $\mathcal{A}$, let $A_1, \ldots, A_{n_3}$ be any basis of $\mathcal{A}$ such that the pivots of the $A_i$ are distinct. We call the set of all pivots $\{p(A_1), \ldots, p(A_{n_3})\}$ the \emph{pivots} of $\mathcal{A}$, and we denote this set by $p(\mathcal{A})$. 
    %
\end{definition}

\begin{remark}\label{rem:pivots-as-cell-locations}
    Let $T$ be a tensor with distinct pivots of its $1$-slices, then this set of pivots is a special case of a linearly independent set of cell-locations for $T$ in direction~$1$.
\end{remark}

%
%
%
%
%

%

%
%
%

\begin{definition}\label{def:rho}
    For any subspace of $n_1 \times n_2$ matrices $\mathcal{A}$, we let $\rho(\mathcal{A})$ be the smallest number $r$ such that there are subsets $S_1 \subseteq [n_1], S_2 \subseteq [n_2]$ of size $|S_1| + |S_2| = r$ for which $p(\mathcal{A}) \subseteq (S_1 \times [n_2]) \cup ([n_1] \times S_2)$.
\end{definition}

\begin{definition}
    For any subspace of $n_1 \times n_2$ matrices $\mathcal{A}$, we let $\sigma(\mathcal{A})$ be the largest number $s$ such that there are matrices $A_1, \ldots, A_s \in \mathcal{A}$ for which for every $i\neq j$ we have $f(A_i) \neq f(A_j)$ and $g(A_i) \neq g(A_j)$.%
\end{definition}
%

The following is an application of Kőnig's theorem on the equivalence of maximum matchings and minimal vertex covers in bipartite graphs 
%

\begin{lemma}[\cite{meshulam1985maximal}] \label{lem:pivots-max-ind-min-lines}
    For any subspace of $n_1 \times n_2$ matrices $\mathcal{A}$ we have $\rho(\mathcal{A}) = \sigma(\mathcal{A})$.
\end{lemma}
%
%
%

%
%
%

%

%
%
    
%

%
%
%
%

%


\subsubsection{Pivots of a tensor}

Having discussed the notion of pivots of a matrix subspace, we will now associate to any tensor matrix subspaces obtained by taking the span of the slices of the tensor. We have considered slice spans before when defining the parameters $\subrank_\ell$, but here we will need to be more precise and distinguish six matrix subspaces instead of three (because the property we consider is not invariant under taking transpose). We then consider the pivots of these subspaces. We will relate these to the parameters $\subrank_\ell$ and to the border subrank. Combining these ideas we then give another proof of the cube-root lower bound on the asymptotic subrank.

\begin{definition}\label{def:pivots-tensors}
    Let $T = \sum_{i,j,k} T_{i,j,k}\, e_i \otimes e_j \otimes e_k \in \FF^{n_1} \otimes \FF^{n_2} \otimes \FF^{n_3}$ be a tensor.
    Define the subspaces
    \begin{align*}
        &\mathcal{A}_{1,2} = \linspan \{A_1, \ldots, A_{n_3}\} \subseteq \FF^{n_1} \otimes \FF^{n_2},\quad A_k = \sum_{i,j} T_{i,j,k}\, e_i \otimes e_j\\
        &\mathcal{A}_{2,1} = \linspan \{A_1, \ldots, A_{n_3}\} \subseteq \FF^{n_2} \otimes \FF^{n_1},\quad A_k = \sum_{i,j} T_{i,j,k}\, e_j \otimes e_i\\
        &\mathcal{A}_{1,3} = \linspan \{A_1, \ldots, A_{n_2}\} \subseteq \FF^{n_1} \otimes \FF^{n_3},\quad A_j = \sum_{i,k} T_{i,j,k}\, e_i \otimes e_k\\
        &\mathcal{A}_{3,1} = \linspan \{A_1, \ldots, A_{n_2}\} \subseteq \FF^{n_3} \otimes \FF^{n_1},\quad A_j = \sum_{i,k} T_{i,j,k}\, e_k \otimes e_i\\
        &\mathcal{A}_{2,3} = \linspan \{A_1, \ldots, A_{n_1}\} \subseteq \FF^{n_2} \otimes \FF^{n_3},\quad A_i = \sum_{j,k} T_{i,j,k}\, e_j \otimes e_k\\
        &\mathcal{A}_{3,2} = \linspan \{A_1, \ldots, A_{n_1}\} \subseteq \FF^{n_3} \otimes \FF^{n_2},\quad A_i = \sum_{j,k} T_{i,j,k}\, e_k \otimes e_j.
    \end{align*}
%
 %
%
    Let $\rho_{{\ell_1},{\ell_2}}(T) = \rho(\mathcal{A}_{{\ell_1},{\ell_2}})$.
\end{definition}

\begin{remark}
    Note that $\rho_{{\ell_1},{\ell_2}}(T)$ can depend on the order of ${\ell_1},{\ell_2}$, which defines which direction enumerates the rows and which the columns. This is a direct consequence of using the lexicographic order in the definition of the pivots. For example, the tensor~$T$ with %
    \[
    \mathcal{A}_{1,2} = \Bigl\{\begin{pmatrix}
        0&0&1\\
        1&0&0
    \end{pmatrix},\begin{pmatrix}
        0&1&0\\
        0&0&0
    \end{pmatrix}\Bigr\},
    \]
    has 
    \[
    \mathcal{A}_{2,1} = \Bigl\{\begin{pmatrix}
        0&1\\
        0&0\\
        1&0
    \end{pmatrix},\begin{pmatrix}
        0&0\\
        1&0\\
        0&0
    \end{pmatrix}\Bigr\}.
    \]
    Then $\rho_{1,2}(T)=1\neq2=\rho_{2,1}(T)$. %
\end{remark}

%

%

\subsection{Alternative proof of cube-root lower bound via pivots} \label{subsec:cube-root-via-pivots}

In this section we will give another proof of the cube-root lower bound on asymptotic subrank (\autoref{th:cube-root-lb}) using pivots. %

\begin{theorem}\label{th:pivot-uncertainty}
    Let $T \in \FF^{n_1} \otimes \FF^{n_2} \otimes \FF^{n_3}$ be concise.
    For every distinct ${\ell_1},{\ell_2},{\ell_3}$ we have 
    \[
    \rho_{{\ell_1},{\ell_2}}(T) \max \{\subrank_{\ell_1}(T), \subrank_{\ell_2}(T)\} \geq n_{\ell_3}.
    \]
\end{theorem}
\begin{proof}
%
%
%
%
%
We give the proof for ${\ell_1}=2, {\ell_2}=3, {\ell_3}=1$ (the proof for the other cases is the same).
%
By an invertible linear transformation of direction $1$ (linear combinations of the $1$-slices) we may assume that for every $i\in[n_1]$, the $i$th $1$-slice is the only one with a non-zero entry in its pivot $p(i)=(f(i),g(i))$. Now we will prove that in every row there are at most $\subrank_2(T)$ pivots and in every column at most $\subrank_3(T)$. %
Let $x$ in the image of $f$, and project the second leg of $T$ to $e_x$. Restricted in the third leg to the columns with pivots in row $x$, $\{k\in[n_3]:\exists i\in[n_1], (f(i),g(i))=(x,k)\}$, this gives the following diagonal submatrix of the $x$th $2$-slice:
\[
\sum_{i: f(i) = x} e_i \otimes e_x \otimes e_{g(i)},
\]
which implies $\subrank_2(T) \geq | \{i : f(i) = x\}|$ for every $x$ in the image of $f$. Similarly, we have the inequality $\subrank_3(T) \geq | \{i : g(i) = y\}|$ for every $y$ in the image of $g$.
%
By \autoref{def:rho}, $p(\mathcal{A}_{2,3}) \subseteq (S_2 \times [n_3]) \cup ([n_2] \times S_3)$ for some subsets $S_\ell \subseteq [n_\ell]$ such that $|S_2| + |S_3| = \rho_{2,3}(T)$.
%
%
%
%
%
%
%
%
%
%
%
%
Both of these combine to:
\begin{align*}
n_1 &= |\{(f(i), g(i)) : i\in [n_1]\}|\\ &\leq |\{(f(i),g(i)): i \in [n_1], f(i) \in S_2\}| + |\{(f(i),g(i)) : i \in [n_1], g(i) \in S_3\}|\\
&\leq |S_2| \subrank_2(T) + |S_3| \subrank_3(T)\\
&\leq (|S_2| + |S_3|) \max \{\subrank_2(T), \subrank_3(T)\}\\
&= \rho_{2,3}(T) \max \{\subrank_2(T), \subrank_3(T)\}.
\end{align*}
This proves the claim.
\end{proof}

%

%

%
%
%
%
%
%
%
%
%
%
%
%
%
%
%
%
%
%
%
%
%
%
%
%
%
%
%

%

%

%

\begin{definition}
The \emph{border subrank} of a tensor $T \in \FF^{n_1} \otimes  \FF^{n_2} \otimes \FF^{n_3}$, denoted by~$\bordersubrank(T)$, is the largest number~$r$ such that there are matrices $A(\eps), B(\eps), C(\eps)$ whose coefficients are Laurent polynomials in the formal variable $\eps$, such that
\[
(A(\eps) \otimes B(\eps) \otimes C(\eps))  T = \langle r\rangle + \eps S_1 + \eps^2 S_2 + \cdots + \eps^t S_t
\]
where $\langle r\rangle$ is the $r \times r \times r$ diagonal tensor and the $S_i$ are arbitrary tensors (with coefficients in $\FF$). 
\end{definition}

The following known relation between border subrank and asymptotic subrank can be found in \cite[Proposition 5.10]{strassen1988asymptotic} and \cite[Proposition 15.26]{burgisser1997algebraic}.
\begin{lemma}\label{lem:border-asymp-subrank}
    For any tensor~$T$ we have $\bordersubrank(T) \leq \asympsubrank(T)$.
\end{lemma}

\begin{theorem}\label{th:rho-bordersubrank}
    For any tensor $T$ and any ${\ell_1}\neq {\ell_2}$ we have
    $\rho_{{\ell_1},{\ell_2}}(T) \leq \bordersubrank(T)$.
\end{theorem}
\begin{proof}
    Let $r = \rho_{{\ell_1},{\ell_2}}(T)$.
    Then by \autoref{lem:pivots-max-ind-min-lines}, after taking appropriate linear combinations of the slices of $T$, $T$ has an $r \times r \times r$ subtensor with slices $A_1, \ldots, A_r$ whose pivots $p(A_i)$ are pairwise disjoint in both coordinates. We may assume that $p(A_i) = (i,i)$ and that the $i$'th row of $A_i$ is zero except for the pivot coordinate $(i,i)$. Thus $A_i$ is of the form
    \[
    A_i = e_i \otimes e_i 
    %
    + \sum_{j > i} \sum_{k} (A_i)_{j,k}\, e_j \otimes e_k.
    \]
    So $T$ restricts to the tensor
    \[
    S = \sum_i e_i \otimes e_i \otimes e_i 
    %
    + \sum_{j > i} \sum_{k} (A_i)_{j,k}\, e_i \otimes e_j \otimes e_k.
    \]
    We define the matrices $A(\eps), B(\eps), C(\eps)$ as the ones that map
    \begin{align*}
    &A(\eps) : e_i \mapsto \eps^i e_i\\
    &B(\eps) : e_i \mapsto \eps^{-i} e_i\\
    &C(\eps) : e_i \mapsto e_i.
    \end{align*}
    Then 
    \[
(A(\eps) \otimes B(\eps) \otimes C(\eps))S = \sum_i e_i \otimes e_i \otimes e_i + \sum_{j > i} \eps^{j-i} \sum_{k} (A_i)_{j,k}\, e_i \otimes e_j \otimes e_k
    \]
    so $\bordersubrank(T) \geq r$.
\end{proof}



%

%
%
%
%
%
%
%
%
%
%
%
%
%
%
%
%
%
%
%
%

\begin{theorem} \label{thm:cube-root-via-pivots}
Let $T \in \FF^n \otimes \FF^n \otimes \FF^n$ be concise. Then
$\asympsubrank(T) \geq n^{1/3}$.
\end{theorem}

%
%

\begin{proof}
We consider two cases.
For the first case, suppose there are ${\ell_1},{\ell_2}$ such that $\rho_{{\ell_1},{\ell_2}}(T) \geq n^{1/3}$. Then $\asympsubrank(T) \geq n^{1/3}$, because generally we have $\asympsubrank(T) \geq \bordersubrank(T) \geq \rho_{{\ell_1},{\ell_2}}(T)$ (\autoref{th:rho-bordersubrank}). 
For the second case, suppose for all ${\ell_1} \neq {\ell_2}$ we have $\rho_{{\ell_1},{\ell_2}}(T) < n^{1/3}$. Then $\max \{ \subrank_{\ell_1}(T), \subrank_{\ell_2}(T) \} \geq n^{2/3}$ for all distinct ${\ell_1},{\ell_2}$ (\autoref{th:pivot-uncertainty}). From this follows that there are distinct ${\ell_1},{\ell_2}$ such that $\subrank_{\ell_1}(T) \geq n^{2/3}$ and  $\subrank_{\ell_2}(T) \geq n^{2/3}$. Then $\subrank(T^{\boxtimes 2}) \geq n^{2/3}$ (\autoref{lem:Qi-square}) so $\asympsubrank(T) \geq n^{1/3}$.
\end{proof}


We finish this subsection with a relation between the border subrank and the parameters~$\subrank_\ell$ that may be of independent interest.

%\Jnote{The way to fix the below would be to introduce $\bordersubrank_e(T)$ where $e$ is an approximation degree (as in Section 2 of https://arxiv.org/pdf/1705.09379). Then $\max_e \bordersubrank_e(T) = \bordersubrank(T)$. Then we get the theorem for $\bordersubrank_e(T)$  assuming a lower bound on $|\FF|$ in terms of $\bordersubrank_e(T)$ and $e$.}

\begin{theorem} \label{thm:commrnk_bounds_bordersubrnk} 
Let $T$ be any tensor.
%Suppose $|\FF|>\bordersubrank(T)+1$. 
Suppose $|\FF|=\infty$. 
For every $\ell \in [3]$ we have
\[\bordersubrank(T) \leq \subrank_\ell(T).\]
\end{theorem}

\begin{proof}
    %
    We will give the proof for $\ell=3$.
    Let $q=\bordersubrank(T)$. Then there are matrices $A(\eps)$, $B(\eps)$, $C(\eps)$ whose coefficients are Laurent polynomials over $\FF$ in the variable $\eps$ %
    such that
    \[
    (A(\eps)\otimes B(\eps)\otimes C(\eps))(T)=\langle q\rangle+\eps S(\eps) %
    \]
    for some tensor $S(\eps)$ whose coefficients are polynomials over $\FF$ in $\eps$.
    %
    Note that summing the 3-slices of $\langle q\rangle$ gives $\id_q$.
    Towards taking the sum of the 3-slices on both sides, define $C'(\eps) = (1,\ldots,1) C(\eps)$.
    %Taking the sum of its $3$-slices is acting by yet another local linear transformation on the direction $3$, so we can absorb it into the third transformation to get a new matrix~$C'(\eps)$ such that 
    Then
    $(A(\eps)\otimes B(\eps)\otimes C'(\eps))(T)=\id_q+\eps N(\eps)$ for some matrix $N(\eps)\in M_{q}(\FF[\eps])$. The determinant of the matrix $\id_q+\eps N(\eps)$ is a nonzero polynomial in $\eps$. %, say of degree $d$. 
    %at most $q$ 
    %\refnote{25}{Why is its degree bounded? It seemingly can be strictly larger than $q$.} \Jshort{We have changed the field assumption to infinite field and adjusted the proof.}
    %, as it is evaluated to a nonzero element for $\eps=0$. 
    Thus, it has finitely many roots and we have another evaluation point for which it is non-zero. %, as long as $|\FF|>q+1$. 
    Let $x\in\FF$ be such an evaluation point. Then $A(x),B(x),C'(x)$ are linear transformations over~$\FF$ for which we get a $3$-slice $(A(x)\otimes B(x)\otimes C'(x))(T)= \id_q +xN(x)\in M_q(\FF)$ of rank~$q$, showing $\subrank_3(T)\geq q=\bordersubrank(T)$.
\end{proof}

Combining \autoref{th:rho-bordersubrank} and \autoref{thm:commrnk_bounds_bordersubrnk}, assuming that $|\FF| = \infty$, 
%> \bordersubrank(T)+1$,
we find for any directions ${\ell_1},{\ell_2},{\ell_3} \in [3]$ such that ${\ell_1}\neq {\ell_2}$ the relation
%
$\rho_{{\ell_1},{\ell_2}}(T) \leq \bordersubrank(T) \leq \subrank_{\ell_3}(T).$
%

%
%
%

%
%
%
%

%
%

%

\subsection{Square-root lower bound for symmetric tensors} \label{sub:sqrt-assub-equal-pivots}

For symmetric tensors we can improve the cube-root lower bound of \autoref{thm:cube-root-via-pivots} to a square-root lower bound, using the notion of pivots.

%

%

%
\begin{theorem}\label{cor:symmetric}
    Let $T \in \FF^{n} \otimes \FF^n \otimes\FF^n$ be symmetric and concise. Then  $\asympsubrank(T) \geq \sqrt{n}$.
\end{theorem}

In fact, we will prove a stronger theorem which implies \autoref{cor:symmetric}. For this we will define the notion of being pivot-matched.
%

First, for any collection of matrices $A_1, \ldots, A_n \in \FF^{m \times k}$ we define the pivot map 
\[
p = (f,g) : [n] \to [m] \times [k]
\]
by mapping $i\in [n]$ to the pivot $p(A_i) = (f(A_i), g(A_i))$ of $A_i$.
%

%

%

%

\begin{definition}[Pivot-matched] \label{def:pivot-matched}
    We call $T\in\FF^{n}\otimes\FF^{n}\otimes\FF^{n}$ \emph{pivot-matched} if there are tensors $T_A$ and $T_B$ isomorphic 
    %\refnote{26}{What does "isomorphic" mean here?}\Jshort{Added it to the preliminaries} 
    to $T$ such that, 
    %
%
    letting $A_i = \sum_{j,k} (T_A)_{i,j,k}\, e_j \otimes e_k \in \FF^{n} \otimes \FF^{n}$ for $i \in [n]$ be the 1-slices of $T_A$, and $p_A=(f_A,g_A)$ the pivot map of these matrices, and
    letting $B_k = \sum_{i,j} (T_B)_{i,j,k}\, e_i \otimes e_j \in \FF^{n}\otimes \FF^{n}$ for $k \in [n]$ be the 3-slices of $T_B$, and $p_B=(f_B,g_B)$ the pivot map of these matrices, we have that $p_A$ and $p_B$ are injective and  for all $i \in [n]$ we have $f_A(i)=f_B(i)$.
%
%
%
%
%
\end{definition}

The condition that $p_A$ and $p_B$ are injective can always be satisfied by \autoref{lem:pivot-set-invariant}.

We note that in the above definition, since pivots are defined using the lexicographic ordering, to make our proofs work it is crucial to define the slices $A_i$ as stated (and not as their transpose) and the same for the $B_i$.

%



%
%
%
%

    %
    %
    
    %
%

\begin{remark}
    By reordering the $1$-slices of $T_A$ and the $3$-slices of $T_B$, the definition is equivalent to requiring a seemingly weaker requirement of having the same number of slices with pivots in each row for the two sets of slices, i.e.~$\forall i\in[n]:|f_A^{-1}(i)|=|f_B^{-1}(i)|$.
\end{remark}

%
%
%
%
%
%
%
%
%
%
%
%
%
%
%

\begin{theorem} \label{thm:samePivotsInRow_sqrtLowerBound_asymsub}
    Let $T\in\FF^{n}\otimes\FF^{n}\otimes\FF^{n}$ be a concise tensor that is pivot-matched. Then 
    \[\asympsubrank(T)\geq\sqrt{n}.\]
\end{theorem}

%
%
%
%

\begin{remark}
    We do not know whether \autoref{cor:symmetric} can be used to prove $\asympsubrank(T) \geq \sqrt{n}$ for general non-symmetric concise $T \in \FF^{n \times n \times n}$. 
    A natural approach is to symmetrize $T$ by taking $\mathrm{sym}(T) = \boxtimes_{\pi \in S_3} \pi T$ (where we let the symmetric group $S_3$ act by permuting the three tensor legs) and apply \autoref{cor:symmetric} to $\mathrm{sym}(T)$.
    We note, however, that the asymptotic subrank of a tensor generally does not behave nicely with respect to symmetrization (in contrast with matrix rank). 
    In particular, if we let $T = \langle n,1,1\rangle$, which has asymptotic subrank $1$, then the product of the cyclic permutations of $T$, $T \boxtimes (1,2,3) T \boxtimes (1,2,3)^2 T$, equals the matrix multiplication tensor $\langle n,n,n\rangle \in \FF^{n^2 \times n^2 \times n^2}$ which has asymptotic subrank~$n^2$. In other words, tensoring $T$ with permutations of $T$ can drastically increase the asymptotic subrank.
\end{remark}

%
\begin{lemma}\label{lem:upper-triangular}
    Suppose $T\in\FF^n\otimes\FF^n\otimes\FF^n$ is a concise tensor that is %
    pivot-matched. 
    Then
    $T\boxtimes T$ restricts to a tensor of the form
\begin{multline}\label{eq:req-form}
\sum_{i=1}^n \big[e_i\otimes e_{f(i)}\big]\otimes\big[e_{f(i)}\otimes e_{g(i)}\big] \otimes \big[e_{h(i)}\otimes e_i\big]
+ \sum_{i=1}^n\sum_{j<k}\sum_{l,u,v=1}^n c_{ijkluv}\big[e_i\otimes e_j\big]\otimes \big[e_k\otimes e_l\big]\otimes \big[e_u\otimes e_v].
\end{multline}
for $f,g,h:[n]\to[n]$ functions such that $i\mapsto(f(i),g(i))$ is injective.
\end{lemma}

\begin{proof}
Let $A_1,\dots,A_n$ be the $1$-slices of $T_A$ as in \autoref{def:pivot-matched}, and $B_1,\dots,B_n$ the $3$-slices of $T_B$. Without loss of generality, we may assume that
\beq\label{eq:pivot1}
(A_i)_{p_A(j)} = (B_i)_{p_B(j)} = \delta_{ij}
\eeq
where $p_A=(f_A,g_A)$ is the pivot map for the $1$-slices of $T_A$ (rows enumerated by direction~$2$) and $p_B=(f_B,g_B)$ is for the $3$-slices of $T_B$ (rows by direction $1$), otherwise we can take invertible linear combinations of the slices to guarantee it.

%
%
%
%
Using these two slicings, we have the decomposition
\beqn
T^{\boxtimes 2}\geq T_A\boxtimes T_B
=
\sum_{i,u=1}^n\sum_{j,v=1}^n\sum_{k,w=1}^n (A_i)_{jk}(B_w)_{uv} (e_i\otimes e_u)\otimes (e_j\otimes e_v)\otimes (e_k\otimes e_w).
\eeqn

We apply the following three linear maps to the respective legs of this tensor:
\begin{align*}
U &= \sum_{i=1}^n E_{i,i}\otimes E_{f_A(i),f_A(i)}\\
V &= \sum_{v=1}^n E_{f_B(v),f_B(v)}\otimes E_{g_B(v),g_B(v)}\\
W &= \sum_{w=1}^nE_{g_A(w),g_A(w)}\otimes E_{w,w}.
\end{align*}
where $E_{i,i}$ is the projector onto $e_i$.
A small calculation then gives
\begin{multline*}
T^{\boxtimes 2}\geq(U\otimes V\otimes W)(T_A\boxtimes T_B)
=\\
\sum_{i,v,w=1}^n
(A_i)_{f_B(v)\,g_A(w)}\, (B_w)_{f_A(i)\,g_B(v)}\, 
(e_i\otimes e_{f_A(i)})\otimes (e_{f_B(v)}\otimes e_{g_B(v)})\otimes (e_{g_A(w)}\otimes e_w).
\end{multline*}
This is the final tensor. We now proceed to show how it fits the desired form \eqref{eq:req-form}.

By definition of $f_A$ as giving the first non-zero rows of the $A$ slices, we get non-zero contributions only from the $(i,v)$-pairs satisfying $f_B(v) \geq f_A(i)$ from the first factor after the triple sum.
This already shows that there is an upper-triangular structure present, by looking at the second component of the first leg and the first component of the second leg.

To finish showing $(U\otimes V\otimes W)(T_A\boxtimes T_B)$ is in the desired form we focus on the pairs $(i,v)$ satisfying $f_A(i) = f_B(v)$, aiming to show that this fits the first sum in \eqref{eq:req-form}. 
%\refnote{27}{Suggesting to number the equation at the statement of the lemma and referring to it here}\Jshort{Done}. %
Let~$S$ denote the set of such pairs.
Restricted to $S$, the triple sum becomes
\beqn
\sum_{(i,v)\in S}\sum_{w=1}^n 
(A_i)_{f_A(i)\,g_A(w)}\, (B_w)_{f_B(v)\,g_B(v)}\, 
(e_i\otimes e_{f_A(i)})\otimes (e_{f_B(v)}\otimes e_{g_B(v)})\otimes (e_{g_A(w)}\otimes e_w).
\eeqn
By~\eqref{eq:pivot1}, we have that 
\beqn
(B_w)_{f_B(v)\,g_B(v)} = (B_w)_{p_B(v)} = \delta_{vw}.
\eeqn
Hence, the sum over~$w$ collapses to a single term and we are left with the tensor 
\beqn
\sum_{(i,v)\in S}
(A_i)_{f_A(i)\,g_A(v)}\, 
(e_i\otimes e_{f_A(i)})\otimes (e_{f_B(v)}\otimes e_{g_B(v)})\otimes (e_{g_A(v)}\otimes e_v).
\eeqn
Fix $i,j$ such that $f_A(i) = j$.
Then~$v$ ranges over the set $f_B^{-1}(j)$, so all the $B$-slices whose pivot lies on row~$j$.
But by assumption, these are precisely the $A$-slices whose pivot lies on row~$j$.
This implies that the coordinate $(j,g_A(v))$ is the pivot of some $A$-slice.
Again by~\eqref{eq:pivot1} we  get that
\beqn
(A_i)_{f_A(i)\,g_A(v)} = \delta_{iv}.
\eeqn
Thus our tensor looks like
\beqn
\sum_{i=1}^n 
(e_i\otimes e_{f_A(i)})\otimes (e_{f_B(i)}\otimes e_{g_B(i)})\otimes (e_{g_A(i)}\otimes e_i).
\eeqn
Note that the maps $i \mapsto (i, f_A(i))$, $i \mapsto (f_B(i), g_B(i))$ and $i \to (g_A(i), i)$ are injective, as required
(i.e.\ this first sum is a diagonal of length~$n$).
\end{proof}



%
%

%
%
%
%
%
%
%
%
%
%

%

%

%
%
%
%
%
%
%
%

%
%
%
%
%
%

%
%
%
%
%
%

%
%
%
%
%
%
%
%
%
%
%
%
%
%
%
%
%
%
%
%
%
%
%
%
%
%
%
%
%
%
%
%
%
%
%
%
%
%
%
%
%
%
%
%
%
%
%
%
%



%

\begin{lemma} \label{lem:border-subrank-for-upper-triag}
Suppose that $T\in\FF^{n^2}\otimes\FF^{n^2}\otimes\FF^{n^2}$ restricts to a tensor $S$ of the form
\begin{multline*}
\sum_{i=1}^n \big[e_i\otimes e_{f(i)}\big]\otimes\big[e_{f(i)}\otimes e_{g(i)}\big] \otimes \big[e_{h(i)}\otimes e_i\big]
+ \sum_{i=1}^n\sum_{j<k}\sum_{l,u,v=1}^n c_{ijkluv}\big[e_i\otimes e_j\big]\otimes \big[e_k\otimes e_l\big]\otimes \big[e_u\otimes e_v]
\end{multline*}
for $f,g,h:[n]\to[n]$ functions such that $i\mapsto(f(i),g(i))$ is injective.
Then,
    \[\bordersubrank(T) \geq n.\]
\end{lemma}

\begin{proof}
We define the maps
\begin{align*}
&A(\eps) : [e_i \otimes e_j] \mapsto \eps^{-j} [e_i \otimes e_j]\\
&B(\eps) : [e_i \otimes e_j] \mapsto \eps^{i} [e_i \otimes e_j]\\
&C(\eps) : [e_i \otimes e_j] \mapsto [e_i \otimes e_j].
\end{align*}
Then
\begin{multline*}
(A(\eps) \otimes B(\eps) \otimes C(\eps))S = 
\sum_{i=1}^n \big[e_i\otimes e_{f(i)}\big]\otimes\big[e_{f(i)}\otimes e_{g(i)}\big] \otimes \big[e_{h(i)}\otimes e_i\big]\\
+ \sum_{i=1}^n\sum_{j<k}\sum_{l,u,v=1}^n \eps^{k-j} c_{ijkluv}\big[e_i\otimes e_j\big]\otimes \big[e_k\otimes e_l\big]\otimes \big[e_u\otimes e_v],
\end{multline*}
where the first sum is a diagonal of length~$n$.
Therefore, $\bordersubrank(T) \geq n$.
\end{proof}

\begin{proof}[Proof of \autoref{thm:samePivotsInRow_sqrtLowerBound_asymsub}]
    By \autoref{lem:upper-triangular}, $T^{\boxtimes 2}$ restricts to the upper-triangular form that is the hypothesis for \autoref{lem:border-subrank-for-upper-triag}, which in turn states that it degenerates to a diagonal tensor of size $n$. By \autoref{lem:border-asymp-subrank} this lower bounds its asymptotic subrank $\asympsubrank(T^{\boxtimes 2})\geq\bordersubrank(T^{\boxtimes 2})\geq n$. So $\asympsubrank(T)\geq\sqrt{n}$.
\end{proof}





\section{Asymptotic slice rank versus asymptotic subrank}\label{slicerank-vs-subrank}


%

In this section we prove a lower bound on the asymptotic subrank in terms of the asymptotic slice rank, thus bounding how different the asymptotic subrank and asymptotic slice rank can be. This requires some of the methods developed so far, and some extra ingredients. These extra ingredients are the notion of the minimal covering number of a matrix subspace, which is also known as non-commutative rank in the literature, and a Flanders-type relation to the max-rank.

%
\subsection{Minimal covering number of matrix subspaces and slice spans}\label{subsec:mincov}

For any matrix subspace $\mathcal{A} \subseteq \FF^{n_1} \otimes \FF^{n_2}$ we define the \emph{minimal covering number} $\mincov(\mathcal{A})$ as 
%
%
    %
    the smallest number $m_1 + m_2$ such that there are subspaces $V_1 \subseteq \FF^{n_1}$ and $V_2 \subseteq \FF^{n_2}$ with $\dim V_1 = m_1$ and $\dim V_2 = m_2$ such that $\mathcal{A} \subseteq V_1 \otimes \FF^{n_2} + \FF^{n_1} \otimes V_2$.
%
In the literature, the parameter $\mincov$ is also known as the non-commutative rank~\cite{MR2123060}, and a matrix subspace $\mathcal{A}$ satisfying $\mincov(\mathcal{A}) \leq r$ is sometimes called $r$-decomposable~\cite{meshulam1985maximal}.

%


%
%
%
\begin{lemma}\label{th:Qi-SRi-matrix}
    For any matrix subspace $\mathcal{A}$, $\maxrank(\mathcal{A}) \leq \mincov(\mathcal{A})$.
\end{lemma}
\begin{proof}
%
    %
    %
    There are subspaces $V_1\subseteq\FF^{n_1}$ and $V_2\subseteq\FF^{n_2}$ with $\dim(V_1)+\dim(V_2)=\mincov(\mathcal{A})$ %
    %
    %
    and such that $\mathcal{A} \subseteq V_1\otimes\FF^{n_2}+\FF^{n_1}\otimes V_2$. Any matrix in $\mathcal{A}$ can then be written as $A+B$ for some $A\in V_1\otimes\FF^{n_2}$ and $B\in\FF^{n_1}\otimes V_2$. Then $\rank(A)\leq\dim{V_1}$ and $\rank(B)\leq\dim(V_2)$. This gives $\rank(A+B)\leq\rank(A)+\rank(B)\leq\mincov(\mathcal{A})$, as required.
\end{proof}

% \Jnote{DA Ref 1 comment 2: The sentence after Lemma 6.1 should refer to Lemma 6.1}
% \Jnote{Fixed}
In \autoref{subsec:flanders} we discuss reverse versions of 
\autoref{th:Qi-SRi-matrix}. %\autoref{th:Qi-SRi}.

For any tensor $T \in \FF^{n_1} \otimes \FF^{n_2} \otimes \FF^{n_3}$, let $\mathcal{A}_\ell$ denote the slice spans of $T$ for $\ell \in [3]$. We define $\slicerank_\ell(T) = \mincov(\mathcal{A}_\ell)$ for $\ell \in [3]$.
\autoref{th:Qi-SRi-matrix} directly gives:

\begin{lemma}\label{th:Qi-SRi}
    For any tensor $T$ and every $\ell \in [3]$, we have $\subrank_\ell(T) \leq \slicerank_\ell(T)$.
\end{lemma}

Finally, it is easily seen how $\subrank_\ell$ is a relaxation of $\subrank$, and similarly how $\slicerank$ is a relaxation of $\slicerank_\ell$, leading to the following inequalities: %

\begin{lemma}\label{lem:ranks-vs-ranksi}
    For any tensor $T$ and every $\ell\in[3]$, we have $\subrank(T) \leq \subrank_\ell(T)$.
\end{lemma}
\begin{proof}
    Let $r=\subrank(T)$. Then $T\geq\langle r\rangle$ and specifically there are matrices $A,B,C$ such that $(A\otimes B\otimes C)T$ is a matrix of rank $r$, and a linear combination of the $\ell$-slices acted on by a matrix from the left and a matrix on the right. So the linear combination itself is of rank at least $r$.
\end{proof}

\begin{lemma}
    For any tensor $T$ and every $\ell\in[3]$, we have $\slicerank(T) \leq \slicerank_\ell(T)$.
\end{lemma}    
%
%
%
%
%
\begin{proof}
    %
%
    %
    %This inequality is immediate by noting that $\slicerank_i(T)$ is the minimal size of a slicerank decomposition with the added restriction that the used slices are only $j$-slices with $j\neq i$. 
    %\refnote{28}{This was not very immediate to me -- had to think a bit about it}\Jshort{gave some more details}
%
    Let $T \in \FF^{n_1} \otimes \FF^{n_2} \otimes \FF^{n_3}$.
    The slice rank $\slicerank(T)$ can be characterized \cite{sawin} as the smallest number $r$ such that there are subspaces $V_\ell \subseteq \FF^{n_\ell}$ for which $\sum_\ell \dim V_\ell = r$ and $T \in V_1 \otimes \FF^{n_2} \otimes \FF^{n_3} + \FF^{n_1} \otimes V_2 \otimes \FF^{n_3} + \FF^{n_1} \otimes \FF^{n_2} \otimes V_3$.
    %
    Suppose $\ell = 3$. Let $s = \slicerank_3(T)$. Then $\mincov(\mathcal{A}) = s$ where $\mathcal{A}$ is the $3$-slice span of $T$. Then $\mathcal{A} \subseteq V_1 \otimes \FF^{n_2} + \FF^{n_1} \otimes V_2$ where $\dim V_1 = m_1$, $\dim V_2 = m_2$ and $m_1 + m_2 = s$. So $T \in V_1 \otimes \FF^{n_2} \otimes \FF^{n_3} + \FF^{n_1} \otimes V_2 \otimes \FF^{n_3}$, which gives $\slicerank(T) \leq s$.
\end{proof}
\subsection{Flanders-type bounds on the max-rank}\label{subsec:flanders}

%
%

We know for any matrix subspace $\mathcal{A}$ that $\maxrank(\mathcal{A}) \leq \mincov(\mathcal{A})$ %
(\autoref{th:Qi-SRi}). The goal of this section is to discuss inequalities in the reverse direction. While these are not new, we provide self-contained proofs for the convenience of the reader. Since the tensor parameters $\subrank_\ell$ and $\slicerank_\ell$ are defined in terms of $\maxrank$ and $\mincov$ we immediately get similar inequalities for those. We start with the most basic reverse inequality, which is %essentially 
due to Flanders~\cite{flanders1962spaces}:

\begin{lemma}\label{th:sr-q-2-matrix} %
    %
    Let $\mathcal{A}$ be a matrix subspace over the field $\FF$.
    Suppose $|\FF| > \maxrank(\mathcal{A})$. Then we have that $\mincov(\mathcal{A}) \leq 2 \maxrank(\mathcal{A})$.
\end{lemma}

%

% \begin{proof}
% %
%     Let $A$ be a maximal-rank matrix in $\mathcal{A}$, and $r:=\rank(A)=\maxrank(\mathcal{A})$. We may diagonalise $A$ with row and column operations, which correspond to multiplying $A$ from left and right by invertible matrices $R,C$. Look at $\mathcal{B}:=\{R X C:X\in\mathcal{A}\}$ which is a matrix space with $\maxrank(\mathcal{B})=\maxrank(\mathcal{A})$ (as multiplying a matrix by an invertible matrix does not change its rank) and $\mincov(\mathcal{B})=\mincov(\mathcal{A})$ (as if $\mathcal{A}\subseteq V_1\otimes\FF^{n_2}+\FF^{n_1}\otimes V_2$ then $\mathcal{B}\subseteq C(V_1)\otimes\FF^{n_2}+\FF^{n_1}\otimes R(V_2)$ where $C(V_1)=\{Cv:v\in V_1\}$ and similarly for $R(V_2)$, and as images of invertible matrices they have the same dimensions as the domains and the implication can go the other way, using their inverses). The diagonalised $A$, $\tilde{A}\in\mathcal{B}$, is now a block-diagonal matrix with $\id_r$ in its top-left corner and $0$ everywhere else. We claim that similar to $\tilde A$, any matrix in $\mathcal{B}$ must have a zero block as its  bottom-right $(n_1-r)\times(n_2-r)$ submatrix. 
%     Otherwise, if $B_{i,j}\neq0$ for $i\geq n_1-r$ and $j\geq n_2-r$, look at the determinant of $(x\tilde{A}+B)|_{([r]\cup\{i\})\times([r]\cup\{j\})}$ where the restriction means the submatrix of the noted coordinates and $x$ is an indeterminate. The determinant is a polynomial in $x$ of degree at most $r$ and it includes the monomial $B_{i,j}x^r$ which is nonzero and thus it is a nonzero polynomial. As a degree-$r$ polynomial it has at most $r$ different zeros, and as $|\FF|>\maxrank(\mathcal{A})=r$ there is a nonzero evaluation for it, meaning there is an $a\in\FF$ for which $a\tilde{A}+B$ is of rank at least $r+1$ in contradiction for the maximality of $\tilde{A}$'s rank. So all matrices in $\mathcal{B}$ are in $\linspan\{e_i\}_{i=1}^r\otimes\FF^{n_3}+\FF^{n_2}\otimes \linspan\{e_i\}_{i=1}^r$, showing $\mincov(\mathcal{A})=\mincov(\mathcal{B})\leq 2r=2\maxrank(\mathcal{A})$.
% \end{proof}

\autoref{th:sr-q-2-matrix} requires the base field $\FF$ to be large enough. For arbitrary fields, using a slightly larger constant on the right hand side, a similar inequality is still known to be true by a result of Haramaty and Shpilka \cite[Lemma 3.5]{haramaty2010structure} (Lemma~3.7 in the arXiv version).

\begin{lemma}\label{thm:sri_leq_4Qi-matrix} %
    %
    Let $\mathcal{A}$ be any matrix subspace. Then $\mincov(\mathcal{A}) \leq 4 \maxrank(\mathcal{A})$.
\end{lemma}
%
The result in \cite[Lemma 3.5]{haramaty2010structure} is stated in the more general setting of quadratic polynomials. We follow their proof closely, but present it for the special case of matrices (bilinear forms).
\begin{proof}
%
%
%
    Throughout the proof we will consider matrices as vectors in the tensor product $\FF^{n_1}\otimes\FF^{n_2}$.
    Denote $r=\maxrank(\mathcal{A})$. We will find subspaces $U_1\subseteq\FF^{n_1},V_1\subseteq\FF^{n_2}$, both of dimension $r$, such that restricting $\mathcal{A}$ to some complement subspaces for them $U_1^\perp,V_1^\perp$
    %
    we get maximum rank in $\mathcal{A}|_{U_1^\perp\times V_1^\perp}$ at most $\frac{r}{2}$ (as we are working 
    %\refnote{29}{typo}\ILnote{fixed} 
    over arbitrary fields, we are not assuming an inner product. $U_1^\perp,V_1^\perp$ are complements w.r.t.\ an arbitrary direct sum decomposition of $\FF^{n_1},\FF^{n_2}$ as explained in detail below). Repeating this process, we get by induction subspaces $U_1,U_2\dots,U_m$ and $V_1,V_2,\dots,V_m$ such that restricting $\mathcal{A}$ to complements of them all we have the trivial space of only the $0$ matrix. The dimension of both $U_i$ and $V_i$ is at most $\frac{r}{2^{i-1}}$ so the dimension of both $\sum_{i=1}^m U_i$ and $\sum_{i=1}^m V_i$ is at most $\sum_{i=1}^m\frac{r}{2^{i-1}}\leq 2r$. Then $\mathcal{A}\subseteq (\sum_{i=1}^m U_i)\otimes\FF^{n_2}+\FF^{n_1}\otimes(\sum_{i=1}^m V_i)$ shows the desired $\mincov(\mathcal{A})\leq 2r+2r=4r$.

    We have left to show that indeed we can find $U_1,V_1$ as described above. Let $A\in\mathcal{A}$ be a maximal-rank element, so $\rank(A)=r$. Let $A=\sum_{i=1}^r u_i\otimes v_i$ be a rank-decomposition of $A$. Define $U_1:=\linspan\{u_1,\dots,u_r\}$ and $V_1:=\linspan\{v_1,\dots,v_r\}$. By the fact that it is a rank decomposition we have that $\{u_1,\dots, u_r\}$ and $\{v_1,\dots,v_r\}$ are both linearly independent sets of vectors, showing $\dim(U_1)=\dim(V_1)=r$. Complete them to bases $u_1,\dots,u_{n_1}$ and $v_1,\dots,v_{n_2}$ of $\FF^{n_1}$ and $\FF^{n_2}$ respectively. Define $U_1^\perp:=\linspan\{u_{r+1},\dots,u_{n_1}\}$ and $V_1^\perp:=\linspan\{v_{r+1},\dots,v_{n_2}\}$ and denote the projectors onto them with respect to these bases by $\Pi_{U_1^\perp}$ and $\Pi_{V_1^\perp}$ respectively. Define the restricted subspace $\mathcal{A}|_{U_1^\perp\times V_1^\perp}:=\{(\Pi_{U_1^\perp}\otimes\Pi_{V_1^\perp})(B):B\in\mathcal{A}\}$. 
    Let $B\in\mathcal{A}$ be any element in $\mathcal{A}$. Denote its restriction by $B|_{U_1^\perp\times V_1^\perp}$ and denote the rank of this restriction by $s:=\rank(B|_{U_1^\perp\times V_1^\perp})$. We want to show $s\leq\frac{r}{2}$. Write the restriction as a rank-decomposition $B|_{U_1^\perp\times V_1^\perp}=\sum_{i=1}^s x_i\otimes y_i$ for elements $x_i\in U_1^\perp$ and $y_i\in V_1^\perp$. Then $\{u_1,\dots,u_r,x_1,\dots,x_s\}$ and $\{v_1,\dots,v_r,y_1,\dots,y_s\}$ are both linearly independent. 
    Thus, to write the unrestriced $B$ we only need to introduce tensor products of either elements from $U_1$ with $\ell^{(2)}_i\in\FF^{n_2}$ or tensor products of $\ell^{(1)}_i\in\FF^{n_1}$ with elements from~$V_1$, that is: 
    %
    $B=\sum_{i=1}^s x_i\otimes y_i+\sum_{i=1}^r( u_i\otimes \ell^{(2)}_i+\ell^{(1)}_i\otimes v_i)$. Decompose $\ell^{(2)}_i=\tilde{v}_i+\tilde{y}_i+\tilde{\ell}^{(2)}_i$ where $\tilde{v}_i\in V_1$, $\tilde{y}_i\in \linspan\{y_1,\dots,y_s\}$ and $\tilde{\ell}^{(2)}_i$ is independent of both subspaces (or $0$). Do the same for the $\ell^{(1)}_i$'s. Denote $\tilde{y}_i=\sum_{j=1}^s b_{i,j}y_j$ and $\tilde{x}_i=\sum_{j=1}^s a_{i,j}x_j$ for $a_{i,j},b_{i,j}\in\FF$, and define ${\ell'}^{(1)}_j:=\sum_{i=1}^r b_{i,j}u_i\in U_1$ and similarly ${\ell'}^{(2)}_j:=\sum_{i=1}^r a_{i,j}v_i\in V_1$. Then 
    %\begin{multline*}
    \[
    B=\underbrace{\sum_{i=1}^s (x_i+{\ell'}^{(1)}_i)\otimes(y_i+{\ell'}^{(2)}_i)}_{=:C}
    \,+\, \underbrace{\sum_{i=1}^r(u_i\otimes(\tilde{v}_i+\tilde{\ell}^{(2)}_i)+(\tilde{u}_i+\tilde{\ell}^{(1)}_i)\otimes v_i)-\sum_{i=1}^s({\ell'}^{(1)}_i\otimes{\ell'}^{(2)}_i)}_{=:D}.
    \]
    %\end{multline*}
    Note that $\{x_i+{\ell'}^{(1)}_i\}_{i=1}^s\cup\{u_i\}_{i=1}^r\cup\{\tilde{\ell}^{(1)}_i:i\in[r],\tilde{\ell}^{(1)}_i\neq0\}$ is linearly independent and so is $\{y_i+{\ell'}^{(2)}_i\}_{i=1}^s\cup\{v_i\}_{i=1}^r\cup\{\tilde{\ell}^{(2)}_i:i\in[r],\tilde{\ell}^{(2)}_i\neq0\}$. In particular we have $\rank(C)=s$ and $B=C\oplus D$. 
    Then $s+\rank(D)=\rank(C)+\rank(D)=\rank(B)\leq r$ so $\rank(D)\leq r-s$. Recall that $A\in U_1\otimes V_1$ so $C\in \linspan\{x_i+{\ell'}^{(1)}_i\}_{i=1}^s\otimes \linspan\{y_i+{\ell'}^{(2)}_i\}_{i=1}^s$ and $D+A$ are supported on linearly independent vectors (in both factors), implying the rank is additive on them:
    $r\geq\rank(B+A)=\rank(C+D+A)=\rank(C)+\rank(D+A)\geq s+(\rank(A)-\rank(D))\geq s+(r-(r-s))=2s$ where the first inequality is by $A+B\in\mathcal{A}$, the second is generally true for any two matrices and the third is the rank of $A$ and the inequality for the rank of $D$ we established earlier.
\end{proof}

From the above follows immediately the corresponding inequalities for the parameters $\subrank_\ell$ and $\slicerank_\ell$.  %

\begin{lemma}\label{th:sr-q-2}
   Let $T \in \FF^{n_1} \otimes \FF^{n_2} \otimes \FF^{n_3}$. Suppose $|\FF| > \subrank_\ell(T)$. Then $\slicerank_\ell(T) \leq 2\subrank_\ell(T)$ for every $\ell \in [3]$.
    %
    %
\end{lemma}
\begin{lemma}\label{thm:sri_leq_4Qi}
    Let $T \in \FF^{n_1} \otimes \FF^{n_2} \otimes \FF^{n_3}$. Then $\slicerank_\ell(T) \leq 4 \subrank_\ell(T)$ for every $\ell \in [3]$.
    %
\end{lemma}

\subsection{Lower bound on asymptotic subrank in terms of slice rank}\label{sec:gap}
%
%

In previous sections we have proved lower bounds on the asymptotic subrank of concise tensors that only depend on the dimensions of the tensor.
Here we shift gears, and use some of the methods that we have developed so far, including the Flanders-type bounds of \autoref{subsec:flanders}, to prove a lower bound on the asymptotic subrank in terms of the slice rank. In fact, we will prove a lower bound on the border subrank in terms of slice rank. As a consequence, we obtain a bound on the gap between asymptotic subrank and asymptotic slice rank.

%
%

%
%
%




%

%

%

%

\begin{lemma}\label{lem:b-subr-sr}
    Let $T \in \FF^{n_1} \otimes \FF^{n_2} \otimes \FF^{n_3}$. Then %
    $\bordersubrank(T^{\boxtimes 3})\geq
    \frac{3}{64}\slicerank(T)^2$.
\end{lemma}

\begin{proof}
    By \autoref{lem:cube-mamu} we have $T^{\boxtimes 3}\geq\langle\subrank_{1}(T),\subrank_2(T),\subrank_3(T) \rangle$.
    Strassen~\cite[Theorem~6.6]{MR882307} proved  for any $e\leq h\leq \ell$ that
    %
    %
    %
    \[
    \bordersubrank(\langle e,h,\ell\rangle)\geq \begin{cases}
    eh-\lfloor\frac{(e+h-\ell)^2}{4}\rfloor &e+h\geq\ell\\
    eh &\textnormal{otherwise}.
    \end{cases}
    \]
    In both cases we have $\bordersubrank(\langle e,h,\ell\rangle)\geq \frac{3}{4}eh$ (the worst case being $\ell=h$).
    Combining these we get 
        \[\bordersubrank(T^{\boxtimes 3})\geq \tfrac34\subrank_1(T)\cdot\subrank_2(T)
        \geq\tfrac34(\tfrac14\slicerank_1(T))\cdot(\tfrac14\slicerank_2(T))\geq \tfrac34(\tfrac14\slicerank(T))^2
        \]
    for which we used \autoref{thm:sri_leq_4Qi} and \autoref{lem:ranks-vs-ranksi}.
\end{proof}

%
%
%\ILnote{Comment 8. I added the following paragraph and changed the theorem below to use limsup rather than the limit. I don't think the proof should be changed at all.}
Unlike rank and subrank, slice rank is not sub- or super-multiplicative under Kronecker product, so Fekete's lemma cannot be used to guarantee its asymptotic version's existence. 
%In \cite{Christandl_2021} it was shown that over the complex numbers the asymptotic slice rank always exists, but 
%In order to keep the theorem as general as possible, we consider the limsup, rather than the limit, of the regularisation in the theorem below. 
To state the following theorem, let $\asympsupslicerank(T):=\limsup_{m\to\infty} \slicerank(T^{\boxtimes m})^{1/m}$. Of course, this is just $\asympslicerank(T)$ if the limit exists.

\begin{theorem}
\label{th:q-sr-bound}
Let $T \in \FF^{n_1} \otimes \FF^{n_2} \otimes \FF^{n_3}$. Then
% Suppose the limit $\asympslicerank(T) = \lim_{n\to\infty} \slicerank(T^{\boxtimes m})^{1/m}$ exists. Then 
\[
\asympsupslicerank(T) \geq \asympsubrank(T) \geq \asympsupslicerank(T)^{2/3}.
\]
\end{theorem}
\begin{proof}
The first inequality follows directly from the fact that slice rank upper bounds subrank. For the second inequality, we have from \autoref{lem:border-asymp-subrank} and \autoref{lem:b-subr-sr} that $\asympsubrank(T) \geq \bordersubrank(T^{\boxtimes 3})^{1/3} \geq (\tfrac{3}{64}\slicerank(T)^2)^{1/3}$. Now replace $T$ by $T^{\boxtimes m}$, take the $m$th root on both sides and let $m$ go to infinity to get the claim.
\end{proof}

\begin{remark}\label{th:a-subr-sr}
 The constant in \autoref{lem:b-subr-sr} can be slightly improved if we replace the border subrank by the asymptotic subrank.
 Namely, for any $3$-tensor $T$, we have $\asympsubrank(T) \geq (\frac{1}{4} \slicerank(T))^{2/3}$.
 The proof is as follows.
 We have
    $\asympsubrank(T^{\boxtimes 3}) \geq \min_{{\ell_1},{\ell_2}} \subrank_{\ell_1}(T) \subrank_{\ell_2}(T)$ (\autoref{lem:asympsub-min}). Combining with $\subrank_\ell(T) \geq \tfrac14 \slicerank_\ell(T)$ (\autoref{thm:sri_leq_4Qi}) and $\slicerank_\ell(T) \geq \slicerank(T)$ (\autoref{lem:ranks-vs-ranksi}) gives
    \[
    \asympsubrank(T)^3 = \asympsubrank(T^{\boxtimes 3}) \geq \min_{{\ell_1},{\ell_2}} (\tfrac14\slicerank_{\ell_1}(T)) (\tfrac14\slicerank_{\ell_2}(T)) \geq (\tfrac14\slicerank(T))^2.
    \]
    Taking the cube root on both sides gives the claim. Another situation in which we can improve the constant is when the base field $\FF$ is large enough, which allows using $\subrank_\ell(T) \geq \tfrac12 \slicerank_\ell(T)$ (\autoref{th:sr-q-2}) instead of $\subrank_\ell(T) \geq \tfrac14 \slicerank_\ell(T)$ (\autoref{thm:sri_leq_4Qi}) in the above proofs, leading to a stronger bound.
\end{remark}


%%% AUTHOR: optional acknowledgments here
\section*{Acknowledgments} %%  you may comment this out if no Ackno
We thank the Centre de Recherches Mathématiques Montréal, where this work was initiated. 
We also thank Jurij Volcic and Vladimir Lysikov for helpful and stimulating discussions. We thank Koen Hoeberechts for useful comments. MC thanks the National Center for Competence in Research SwissMAP of the Swiss National Science Foundation and the Section of Mathematics at the University of Geneva for their hospitality. IL thanks the QMATH Center at the University of Copenhagen for the hospitality during his stay.

%%% AUTHOR:
%%% Bibliography goes here. Note that the arXiv cannot process bibtex
%%% or biber bibliographies.  Example of acceptable bibliograpy format:
% \bibliographystyle{amsplain}
% \begin{thebibliography}{99}
% \bibitem{bergelson-johnson-moreira}
% Vitaly Bergelson, John H. Johnson Jr., and Joel Moreira.
% \newblock New polynomial and multidimensional extensions of classical partition
%   results.
% \newblock 2015, arXiv:1501.02408.

% \bibitem{cilleruelo}
% Javier Cilleruelo.
% \newblock Combinatorial problems in finite fields and {S}idon sets.
% \newblock {\em Combinatorica}, 32(5):497--511, 2012.

% \end{thebibliography}

\bibliographystyle{alphaurl}
\bibliography{refs}


%% AUTHOR: You can generate such a bibliography from a .bib file by 
%% running pdflatex/bibtex/pdflatex/pdflatex and then pasting the .bbl file
%% between \begin{thebibliography} and \end{bibliography}

%%% AUTHOR: Include a short description of each author following the
%%% structure below. Use the same short tags used previously.  
%%% Use \imageat{} and \imagedot{} instead of "@" and "." in
%%% email addresses-this replaces the symbols with graphics to avoid 
%%% e-mail address harvesting from the .pdf file
\begin{dajauthors}
\begin{authorinfo}[briet]
Jop Briët\\
Centrum Wiskunde \& Informatica\\
Amsterdam, The Netherlands\\
j\imagedot{}briet\imageat{}cwi\imagedot{}nl\\
%\url{https://homepages.cwi.nl/~jop/}
\end{authorinfo}
\begin{authorinfo}[christandl]
Matthias Christandl\\
University of Copenhagen\\
Copenhagen, Denmark\\
christandl\imageat{}math\imagedot{}ku\imagedot{}dk
\end{authorinfo}
\begin{authorinfo}[leigh]
  Itai Leigh\\
  Tel Aviv University\\
  Tel Aviv, Israel\\
  itai\imagedot{}leigh\imageat{}mail\imagedot{}huji\imagedot{}ac\imagedot{}il
\end{authorinfo}
\begin{authorinfo}[shpilka]
  Amir Shpilka\\
  Tel Aviv University\\
  Tel Aviv, Israel\\
  shpilka\imageat{}tauex\imagedot{}tau\imagedot{}ac\imagedot{}il
\end{authorinfo}
\begin{authorinfo}[zuiddam]
  Jeroen Zuiddam\\
  University of Amsterdam\\
  Amsterdam, The Netherlands\\
  j\imagedot{}zuiddam\imageat{}uva\imagedot{}nl
\end{authorinfo}

\end{dajauthors}

\end{document}

%%% Local Variables:
%%% mode: latex
%%% TeX-master: "daj-template"
%%% End:
